
\medterm{1p36 Deletion Syndrome} A genetic disorder caused by loss of a piece of chromosome 1. It may cause developmental and intellectual disability, small head size, low muscle tone, seizures, hearing loss, heart problems, and distinctive facial features; severity varies.

\synonyms
Monosomy 1p36; Chromosome 1p36 Deletion Syndrome

\medterm{22q11.2 Deletion Syndrome} A genetic condition caused by a missing piece of chromosome 22. It may cause heart defects, palate problems, reduced immune defenses, low calcium, distinctive facial features, learning or developmental difficulties, and increased risk of mental illness. The effects differ widely between people.

\synonyms
DiGeorge Syndrome Spectrum; Velocardiofacial Syndrome (VCFS); 22q11.2DS

\medterm{D And C} Dilation and curettage, a procedure that opens the cervix and removes tissue from the uterine lining for treatment or diagnosis.

\medterm{D And E} Dilation and evacuation, a procedure that gently opens the cervix and removes pregnancy tissue from the uterus using suction and instruments. It is performed by trained clinicians for selected pregnancy-related indications.

\medterm{D \& C} Dilation and curettage, a procedure in which the cervix is widened and tissue is removed from the uterine cavity for treatment or diagnosis.

\medterm{D2 Receptor} A protein on cells that responds to dopamine, a brain chemical involved in movement, motivation, and reward. Some medicines work partly by blocking or stimulating this receptor.

\medterm{Da Costa's Syndrome} An older term for recurring chest discomfort, awareness of heartbeat, tiredness, and breathlessness without a clear structural heart problem. Current evaluation should still exclude heart, lung, blood, and anxiety-related causes.


\medterm{Dacarbazine} An alkylating chemotherapy medicine used for Hodgkin lymphoma, melanoma, and selected other cancers; important toxicities include severe nausea, bone-marrow suppression, and liver injury.

\medterm{Daclizumab} A monoclonal antibody that blocks the interleukin-2 receptor on activated immune cells. It was used for selected immune conditions but was withdrawn or restricted in many settings because of serious liver and other risks.

\medterm{Dacomitinib} A targeted medicine that blocks several members of the epidermal-growth-factor receptor family and is used for selected EGFR-mutated lung cancers. Diarrhea, skin reactions, mouth sores, and lung or liver problems may occur.

\medterm{Dacrocytes (Teardrop Cells)} Teardrop-shaped red blood cells seen on a blood smear. They may occur when red blood cells pass through scarred or crowded bone marrow, and can be seen in myelofibrosis, bone-marrow cancers, and severe thalassemia.

\synonyms
Teardrop Cells

\medterm{Dacryoadenitis} Inflammation of a tear-producing gland, causing painful swelling near the outer upper eyelid. It may be caused by infection, inflammation, or a condition affecting the whole body.

\medterm{Dacryocystitis} Infection or inflammation of the tear sac near the inner corner of the eye, usually because the tear-drainage passage is blocked. It can cause pain, redness, swelling, and pus and may need prompt treatment.

\medterm{Dacryocystocele} A fluid-filled swelling near the inner corner of the eye caused by blockage and widening of the tear-drainage system. In a newborn it may interfere with breathing if it extends into the nose and needs prompt assessment.

\medterm{Dacryocystography} An imaging test in which contrast material is put into the tear-drainage passages to show narrowing, blockage, abnormal connections, or leakage. It may help plan treatment for persistent tearing or repeated tear-sac infection.

\medterm{Dacryocystorhinostomy} Surgery creating a new passage from the tear sac into the nose to bypass a blocked tear duct. It is used for persistent tearing or repeated tear-sac infections.

\medterm{Dacryocyte Analysis} Examination or counting of teardrop-shaped red blood cells on a blood smear; numerous dacryocytes may occur with marrow fibrosis, infiltration, severe anemia, or splenic disease.

\medterm{Dacryolithiasis} Formation of a stone or concretion within the lacrimal drainage system, which may cause excessive tearing, discharge, recurrent dacryocystitis, or obstruction.

\medterm{Dacryostenosis} Narrowing or blockage of the passages that drain tears from the eye into the nose. It can cause constant tearing, sticky discharge, or repeated infection of the tear sac.

\medterm{Dactinomycin} An anticancer medicine that binds DNA and interferes with gene activity and cell growth. It is used for selected cancers and can cause bone-marrow suppression, mouth sores, nausea, and tissue injury if it leaks outside a vein.

\medterm{Dactylitis} Swelling and inflammation of an entire finger or toe, producing a sausage-like appearance. It can occur with sickle-cell disease, psoriatic arthritis, infection, or other inflammatory conditions.

\medterm{Dactylomegaly} Abnormal enlargement of a finger or toe, which may occur as an isolated finding or as part of an endocrine, skeletal, vascular, or inflammatory disorder.

\medterm{Dactylosymphysis} A birth difference in which two neighboring fingers or toes are joined together by skin, bone, or both. The connection may limit movement or affect growth and hand or foot function.

\medterm{Daidzein} A naturally occurring isoflavone found in soy and other legumes that can be converted by intestinal bacteria into related compounds.

\medterm{Daily Bowel Care Program} A planned routine for preventing and treating constipation or fecal incontinence, often including timing, fluids, diet, activity, medication, and rectal measures when prescribed.

\medterm{Daily Value} A reference amount used on U.S. Nutrition Facts labels to show how much a nutrient in one serving contributes to a typical daily diet.

\medterm{Dakin Solution} A dilute sodium hypochlorite solution used as a topical antiseptic for selected contaminated or infected wounds. Concentration and exposure time must be controlled because stronger solutions can injure viable tissue.

\medterm{Dalteparin} An injectable blood thinner used to prevent or treat clots in veins and the lungs. It reduces clotting through the body’s natural anticoagulant system; bleeding is the main risk, and kidney disease may require dose adjustment.

\medterm{Dalteparin Sodium} The sodium salt of dalteparin, a low-molecular-weight heparin used to prevent or treat venous thromboembolism; bleeding is the principal serious adverse effect.

\medterm{Daltonism} Color-vision deficiency, especially difficulty distinguishing certain colors because of abnormal or absent retinal cone function.
\medterm{DALY} Disability-adjusted life year, a measure combining years of healthy life lost because of premature death and disability.

\medterm{Damage To The Amygdala} Damage to the amygdala can affect fear processing, emotional learning, memory, threat recognition, and social behavior. Causes include trauma, stroke, infection, tumor, epilepsy, or degenerative disease, and effects depend on the side and extent of injury.

\medterm{Damage-Control Resuscitation} Emergency treatment for severe injury that quickly controls bleeding, gives appropriate blood products, prevents dangerous cooling, and corrects clotting problems while avoiding harmful excess fluid.

\medterm{Damage-Control Surgery} A planned series of operations for a critically ill or severely injured person who cannot safely tolerate a long repair. The first operation controls bleeding or contamination; after stabilization, a later operation completes the repair.

\medterm{Damus-Kaye-Stansel Procedure} A heart operation for selected children with a single working ventricle and a blocked or narrowed route from the ventricle to the aorta. It connects the pulmonary artery to the aorta to create a wider route for blood leaving the heart.

\synonyms
DKS Procedure; Damus-Stansel-Kaye Operation

\medterm{Danazol} A synthetic androgen with antigonadotropic effects, used selectively for endometriosis, hereditary angioedema, or other conditions; androgenic, lipid, liver, and thrombotic adverse effects limit use.

\medterm{Dandruff} A common scalp condition causing small white or yellow flakes and sometimes itching. It is often related to oily skin, seborrheic dermatitis, or the scalp's reaction to the normal yeast \textit{Malassezia}; it is not usually contagious. Persistent severe flaking, redness, sores, or hair loss may indicate another scalp disorder.

\medterm{Dandruff And Scalp Conditions} Scalp flaking and itching may result from dandruff, seborrheic dermatitis, psoriasis, contact allergy, fungal infection, or head lice. Persistent redness, sores, pain, or hair loss needs medical assessment.

\medterm{Dandy Fever} An older name for dengue fever, a mosquito-borne viral infection causing sudden fever, headache, severe muscle or joint pain, nausea, or rash. Severe dengue can cause bleeding, shock, and organ failure.

\medterm{Dandy-Walker Malformation} A condition present at birth in which part of the cerebellum does not develop normally and one of the brain's fluid spaces becomes enlarged. It may cause fluid buildup in the brain, delayed development, balance problems, or seizures.

\medterm{Dandy-Walker Syndrome} A birth difference in the cerebellum and nearby fluid spaces of the brain. It may cause enlarged head size from fluid buildup, delayed development, balance problems, seizures, or no symptoms.

\medterm{Dane Particle} The complete infectious form of hepatitis B virus. It contains the virus's genetic material inside a protein shell and outer envelope, allowing it to enter liver cells and reproduce.
\medterm{Dangerous Abbreviations} Abbreviations that can be misread and cause medication or clinical errors; safer practice is to write the complete term when an abbreviation is ambiguous.

\medterm{Danon Disease} A rare inherited disorder, usually caused by an LAMP2 gene change, in which cells cannot properly recycle materials. It may cause heart-muscle disease, muscle weakness, intellectual disability, or eye problems.

\medterm{Dantrolene} A skeletal-muscle relaxant that reduces calcium release from the sarcoplasmic reticulum; it is the specific treatment for malignant hyperthermia and may cause weakness or liver injury.

\medterm{Dapsone} A sulfone antimicrobial used for leprosy, selected dermatitis, and other infections; haemolysis, especially with G6PD deficiency, methaemoglobinaemia, agranulocytosis, and severe skin reactions are important risks.

\medterm{Daptomycin} An antibiotic used for selected serious infections caused by gram-positive bacteria, including resistant strains. It can damage muscles and kidneys, so treatment requires appropriate testing and monitoring.

\medterm{Daptomycin Resistance} Reduced or absent response of bacteria to daptomycin, often because changes in the bacterial membrane prevent the antibiotic from working. It can develop during treatment and limits options for serious infections.

\medterm{Darbepoetin Alfa} A long-acting medicine that stimulates the bone marrow to make more red blood cells. It is used for selected anemia related to kidney disease or cancer treatment, but can raise blood pressure and the risk of blood clots.

\medterm{Darier's Disease} An inherited skin disorder causing firm, greasy, crusted bumps in areas such as the chest, back, scalp, or skin folds. Heat and sweating may worsen it, and nails or mouth lining can also be affected.

\medterm{Dark Adaptation} The gradual adjustment that allows the eyes to see better after moving from bright light into darkness. Poor dark adaptation may result from retinal disease, vitamin A deficiency, aging, or some medicines.

\medterm{Dartos} A thin muscle layer beneath the skin of the scrotum. It tightens or relaxes the scrotal skin in response to temperature and other signals, helping regulate the position of the testicles.

\medterm{Darunavir} An antiretroviral medicine that blocks HIV protease, an enzyme needed to make new viruses. It is used with other HIV medicines and has important interactions and possible liver or skin reactions.

\medterm{Darwin's Tubercle} A harmless small bump on the upper rim of the outer ear. It is a common normal variation in ear shape and usually needs no treatment.

\medterm{Dasatinib} A targeted cancer medicine that blocks abnormal tyrosine-kinase signals, especially those caused by BCR::ABL1. It treats selected blood cancers and may cause low blood counts, fluid around the lungs, bleeding, or abnormal heart rhythms.

\medterm{DASH Diet} An eating pattern designed to lower blood pressure by emphasizing vegetables, fruits, whole grains, beans, nuts, and low-fat dairy while limiting salt, saturated fat, and sugary foods. Individual needs may require adjustment.


\medterm{Dating Ultrasound} An ultrasound examination, usually performed in early pregnancy, used to estimate gestational age and calculate the expected date of delivery.
First-trimester crown--rump length is generally the most accurate ultrasound measurement for dating.

\synonyms
Pregnancy dating scan; gestational-age ultrasound.

\medterm{Daubenton's Plane} An anatomical reference plane extending from the lower margin of the orbit to the upper
margin of the external auditory opening, used in describing skull orientation.

\medterm{Daunorubicin} A chemotherapy medicine that damages cancer-cell DNA and is used mainly for selected blood cancers. It can lower blood counts and, especially after repeated doses, weaken or damage the heart.

\medterm{Daunorubicin Hydrochloride} The hydrochloride salt of daunorubicin, an anthracycline chemotherapy drug that interferes with DNA and is used mainly for acute leukemias. Important risks include bone-marrow suppression and heart damage.

\medterm{Dav Regimen} An abbreviation used for more than one cancer-treatment combination. In older leukemia reports, DAV meant daunorubicin, cytarabine, and etoposide; in newer reports, it may mean daunorubicin, cytarabine, and venetoclax. The source must identify the drugs and cancer before the regimen is used.

\medterm{David Procedure} An operation that replaces an enlarged or weakened aortic root while keeping the person's own aortic valve. The valve is supported inside a graft and the coronary arteries are reattached; bleeding, stroke, valve leakage, or recurrent aortic disease are possible.

\synonyms
David Operation; Valve-Sparing Root Replacement; Reimplantation of Aortic Valve

\medterm{Daw} An abbreviation meaning dispense as written, indicating that the pharmacist should dispense the prescribed product rather than automatically substitute a generic equivalent when permitted by law.

\medterm{Dawn Phenomenon} A rise in blood glucose during the early morning caused by normal hormone changes that help the body wake up. It may be more noticeable in people with diabetes whose insulin supply cannot match the rise.

\medterm{Day Hospital} A facility that provides multidisciplinary medical, psychiatric, rehabilitative, or diagnostic care during the day without overnight admission. It can support intensive treatment while allowing the patient to return home.

\medterm{Day-Blindness} Poor vision in bright light, also called hemeralopia. It can result from retinal disease, cone dysfunction, or other eye disorders.

\medterm{D\&C Instruments} Instruments used to dilate the cervix and remove tissue from the uterus during dilation and curettage.

\medterm{Dcc Gene} A gene that helps guide developing nerve fibers and may also help control cell growth. Harmful changes in this gene can cause some developmental nerve disorders and may contribute to certain cancers.

\medterm{DCIS} Alternate index label; see \textit{Ductal carcinoma in situ}.


\medterm{DDH} Alternate index label; see \textit{Hip dysplasia}.

\medterm{D-Dimer} A substance released into the blood when a blood clot breaks down. A high level can occur with a blood clot, but also with infection, pregnancy, recent surgery, cancer, or other conditions. A normal result can help rule out a clot in carefully selected patients whose initial risk is low or moderate.

\medterm{D-Dimer Test} A blood test measuring a substance released when blood clots are broken down. A normal result can help rule out a clot in carefully selected low-risk people; a high result has many possible causes and does not prove a clot.

\synonyms
D-dimer blood test

\medterm{De Lange Syndrome} A congenital developmental disorder, also called Cornelia de Lange syndrome, involving characteristic facial features, growth and developmental delay, limb abnormalities of variable severity, and intellectual disability.

\synonyms
Cornelia de Lange syndrome

\medterm{De Musset Sign} Repeated head nodding that matches the heartbeat. It can occur when the aortic valve leaks severely, but it is uncommon and cannot confirm the diagnosis without examination and heart tests.

\medterm{De Novo} Newly arising in an individual rather than inherited from either parent; the term is commonly used for genetic variants or disease presentations.

\medterm{De Novo Variant} A genetic change found in a person that was not found in either biological parent. It usually begins in an egg, sperm, or very early embryo and may cause disease; testing parents can help estimate the chance in future pregnancies.

\medterm{De Quervain Tendinitis} Alternate name; see \textit{De Quervain Tenosynovitis}.

\medterm{De Quervain Tenosynovitis} Painful swelling and narrowing around the tendons on the thumb side of the wrist. It causes pain with thumb movement, gripping, or lifting and may cause tenderness or swelling near the base of the thumb.

\synonyms
Tenosynovitis, De Quervain's

\medterm{De Quervain's Tendonitis} Alternate spelling; see \textit{De Quervain Tenosynovitis}.

\medterm{De Quervain's Tenosynovitis} Alternate spelling; see \textit{De Quervain Tenosynovitis}.

\medterm{De Winter ECG Pattern} An ECG pattern that can signal a sudden severe blockage in a major artery supplying the front of the heart, even without the usual ST elevation. It is a heart-attack emergency requiring immediate specialist assessment and treatment; exercise testing must not be used.

\synonyms
de Winter Pattern; de Winter T Waves; de Winter Sign

\medterm{Deacetylase Inhibitor} A medicine or research compound that blocks enzymes removing chemical groups called acetyl groups from proteins. Some alter gene activity and are used for selected cancers; effects and risks depend on the specific medicine.

\medterm{Dead Space} The part of a breath that does not exchange oxygen and carbon dioxide with the blood, including the nose, throat, and larger airways. It increases when air reaches lung areas with too little blood flow.

\medterm{Dead Space Ventilation} Breathing that moves air through areas of the lungs where little or no blood is available for gas exchange. It may increase with a pulmonary embolism, low blood flow, or emphysema and can make breathing less effective.

\medterm{Dead Time} The brief period after an event during which a detector or recording system cannot register another event, limiting the measured count rate.


\medterm{Deafferentation} Loss or interruption of nerve signals carrying sensation from a body area to the brain or spinal cord. It may cause numbness, unusual sensations, poor awareness of position, or persistent nerve pain.

\medterm{Deafness} Complete or severe loss of hearing in one or both ears. The problem may involve sound passing through the outer or middle ear, the inner ear or hearing nerve, or more than one of these areas.

\medterm{Deaner} A former brand name for deanol, also called dimethylaminoethanol, an older product once marketed for behavioral or learning problems. It is not a standard modern treatment, and evidence for medical benefit is limited.

\medterm{Death} The permanent end of the body's vital functions. It means that the heart, breathing, and brain functions have stopped and cannot be restored; medical and legal tests used to confirm death depend on the circumstances.

\medterm{Death Rate} The number of deaths in a defined population during a specified period, usually expressed per population size and used in public-health comparisons.

\medterm{Death Rattle} Noisy, rattling breathing caused by secretions accumulating in the throat or upper airways, often occurring near the end of life when swallowing and coughing weaken.

\medterm{Death Sudden Cardiac (Sudden Cardiac Death)} Alternate index label for Sudden Cardiac Death.

\medterm{Deauville Criteria} A five-point scale used with a PET scan to compare lymphoma activity with normal organs. The score helps clinicians judge whether treatment is working, but it must be interpreted with the type of lymphoma and other findings.

\medterm{Deaver Retractor} A hand-held surgical instrument with a broad, curved blade used to gently move organs or tissues aside. It is often used during chest or abdominal surgery to improve the surgeon's view.

\medterm{Debility} General physical weakness or reduced ability to function, often caused by illness, aging, poor nutrition, prolonged inactivity, or recovery from major stress. The underlying cause should be assessed when weakness persists.

\synonyms
Weakness

\medterm{Debride} To remove dead, damaged, infected, or contaminated tissue from a wound or body site to support healing and control infection.

\medterm{Debridement} Removal of dead, damaged, infected, or contaminated tissue from a wound. It may be done with instruments, irrigation, or other methods to reduce infection, reveal healthy tissue, and support healing.

\medterm{Debulk} To reduce the size or amount of a tumor or other abnormal tissue, usually by surgery or another treatment, without necessarily removing every abnormal cell.

\medterm{Debulking} Removal of as much of a tumor as possible when removing all of it is unsafe or impossible. The aim is to reduce symptoms or improve the effect of later treatment; benefit depends on the cancer and situation.

\medterm{Decalcification} Loss of calcium salts from bone, teeth, or other mineralized tissue. In teeth it may weaken enamel; in bone it can contribute to reduced strength.

\medterm{Decannulation} Removal of a tube, especially a tracheostomy tube, after a person can breathe safely, protect the airway, and clear secretions without it. The site is monitored for bleeding, narrowing, or breathing difficulty.

\medterm{Decanoic Acid} A ten-carbon saturated fatty acid, also called capric acid, found in some fats and studied for antimicrobial, metabolic, and antiseizure effects.

\medterm{Decapitation} Separation of the head from the body, seen in severe trauma (vehicle accidents, train
injuries, explosions), industrial accidents, hanging with judicial execution, or
homicide. Complete decapitation is almost always rapidly fatal due to exsanguination and
brain ischemia.

\medterm{Decay Correction} A calculation that adjusts a radioactive measurement for the gradual loss of activity over time. It allows imaging or laboratory results obtained at different times to be compared accurately.

\medterm{Deceased Donor} Individual who has died (usually from brain death or cardiac death) and donated organs, tissues, or both for transplantation.

\medterm{Decenoic Acid} An unsaturated ten-carbon fatty acid; its biologic effects depend on the position and geometry of its double bond and on tissue and dietary context.

\medterm{Decerebrate Posture} Abnormal straightening of the arms and legs caused by severe injury to the brainstem or nearby brain structures. It is a sign of serious neurologic injury and may occur after head injury, bleeding, or increased pressure in the skull.

\medterm{Decerebrate Posturing} Abnormal straightening of the arms and legs, sometimes with inward turning of the arms, caused by severe brain or brainstem injury. It indicates a medical emergency, but the posture alone cannot predict recovery.

\medterm{Decerebrate Rigidity} Severe, abnormal straightening of the arms and legs, sometimes with the arms turned inward. It usually indicates serious brain or brainstem injury or dangerously high pressure inside the skull and requires emergency assessment.

\medterm{Decerebrate State} A severe neurologic posture marked by extension and inward rotation of the arms and extension of the legs, usually indicating serious injury to the brainstem or structures above it.

\medterm{Decerebration} Severe brain dysfunction that causes abnormal straightening of the arms and legs. It usually results from major brain or brainstem injury or increased pressure inside the skull and is a medical emergency.

\medterm{Decibel} A unit used to describe sound intensity. Higher values mean louder sound; repeated or very loud exposure can damage hearing, depending on the level and duration.




\medterm{Deciding To Have Knee Or Hip Replacement} A decision based on pain, functional limitation, imaging, nonoperative treatment response, surgical risk, expectations, and shared discussion with an orthopedic clinician.

\medterm{Deciding To Quit Drinking Alcohol} The process of planning to stop or reduce alcohol use. Assessment should consider dependence, withdrawal risk, support, and treatment options; people with heavy use may need medical supervision.

\medterm{Decidua} The specialized lining of the uterus during pregnancy. It is shed after delivery and helps support the pregnancy and placenta.

\medterm{Deciduitis} Inflammation of the specialized uterine lining that forms during pregnancy. It may be associated with infection or other pregnancy complications and can require assessment when pain, fever, or bleeding occurs.

\medterm{Deciduoma} An old term for a mass-like thickening of the uterine lining caused by pregnancy hormones. It is not a true tumor and usually disappears when the hormonal stimulation ends.

\medterm{Deciduosis} Pregnancy-related tissue from the uterine lining growing in an unusually large area or outside the uterus, such as on the cervix or pelvic surfaces. It can look like a tumor during examination or surgery but is usually benign.

\medterm{Deciduous Teeth} The 20 primary teeth that erupt during childhood and are gradually replaced by permanent teeth. Caries, infection, trauma, or premature loss can affect nutrition, speech, space for permanent teeth, and oral development.

\medterm{Deciduous Tooth} A primary or baby tooth that is eventually shed and replaced by a permanent tooth. Early decay or loss can affect chewing, speech, and the space needed for permanent teeth.
\medterm{Decision Making} The process of comparing information, options, benefits, risks, and preferences to choose an action or plan, including shared decision-making in healthcare.

\medterm{Decision-Making Capacity} The ability to understand information about a medical decision, consider the possible results, compare choices, and communicate a choice. It applies to a specific decision, may change over time, and is different from legal competence.

\medterm{Decitabine} A cancer medicine that changes abnormal gene activity and interferes with growth of blood-forming cancer cells. It is used for selected blood cancers and may cause low blood counts, infection, bleeding, or fatigue.

\medterm{Declarative Memory} Memory for facts, information, and personal events that a person can consciously remember and describe.

\medterm{Decoction} A preparation made by simmering plant material in water to extract soluble substances. The strength and safety depend on the plant, preparation, dose, contamination, and interactions with medicines.

\medterm{Decompensated Cirrhosis} Advanced liver scarring that has begun to cause problems such as abdominal fluid buildup, internal bleeding, jaundice, or confusion from waste products. It requires specialist care and may require assessment for liver transplantation.

\medterm{Decomposition} The natural breakdown of a dead body after death. Cells and tissues first break down themselves, followed by the action of microorganisms; the speed and appearance depend on temperature, moisture, and other conditions.

\medterm{Decompression} Reduction of pressure on a tissue or body compartment, such as removing pressure from a spinal nerve, draining an obstructed organ, or treating pressure changes after diving.

\medterm{Decompression Sickness} Illness caused by nitrogen bubbles forming in the body when pressure falls too quickly, usually after diving. It may cause joint pain, skin changes, breathing problems, weakness, confusion, or paralysis and needs emergency care.

\medterm{Deconditioning} Loss of strength, fitness, balance, and function after inactivity, illness, or prolonged bed rest. It increases fall risk but can often improve gradually with safe movement, exercise, nutrition, and rehabilitation.

\medterm{Decongestant} A medicine that reduces swelling inside the nose and relieves blocked breathing. Some can raise blood pressure, cause fast heartbeat, or cause rebound congestion, so use should follow instructions and health precautions.

\medterm{Decontamination} Steps that remove or reduce a harmful substance from the skin, eyes, clothing, or digestive tract and prevent further absorption. The correct method depends on the substance; some exposures require washing, irrigation, or specialist treatment.

\medterm{Decorin} A small protein-sugar molecule found around cells. It helps organize collagen, influences tissue growth and repair, and may affect scar formation.

\medterm{Decorticate Posture} Abnormal bending of the arms with straightening of the legs, caused by severe injury to the brain. It is a sign of serious brain damage and reduced consciousness.

\medterm{Decorticate Posturing} Abnormal bending of the arms toward the chest with straightening of the legs, caused by severe brain injury. It signals impaired consciousness and requires emergency assessment; posture alone cannot predict outcome.

\medterm{Decreased Alertness} Decreased alertness is reduced awareness or responsiveness that may result from metabolic illness, intoxication, infection, seizure, stroke, trauma, or sleep-related causes.

\medterm{Decreased Breath Sounds} Softer-than-normal sounds heard when listening to the lungs, suggesting reduced air movement. Causes include fluid or air around the lung, partial lung collapse, blockage, or shallow breathing; sudden one-sided loss can be an emergency.

\medterm{Decreased Concentration} Reduced ability to sustain attention or mental focus; causes include sleep loss, stress, medications, mood disorders, substance use, and neurologic or medical illness.

\medterm{Decreased Dentin} Reduced thickness or formation of dentin, the mineralized tissue beneath tooth enamel; causes include developmental defects, wear, erosion, caries, or dental treatment.

\medterm{Decreased Libido} Reduced sexual desire that may be related to relationships, mood, stress, medications, hormones, chronic illness, pain, or other health factors.

\medterm{Decreased Tear Production} Alternate index label; see \textit{Dry eyes}.

\medterm{Decubitus} A lying or reclining position; the term also appears in decubitus ulcer, a pressure injury caused by prolonged pressure and shear.

\medterm{Decubitus Ulcer} Damage to skin and deeper tissue caused by prolonged pressure, sliding, or rubbing, usually over a bony area. It may begin with persistent redness and progress to an open wound, especially in people with limited movement.

\medterm{Deep} Located farther inside the body or a structure, away from its surface. For example, deep muscles lie beneath muscles closer to the skin.
\medterm{Deep Brain Stimulation} A treatment in which implanted electrodes deliver controlled electrical impulses to selected brain regions. It is used for conditions such as Parkinson disease, essential tremor, and some forms of dystonia.

\medterm{Deep Femoral Artery} The profunda femoris artery, a major branch of the femoral artery that supplies the deep muscles of the thigh and gives rise to perforating branches.

\medterm{Deep Infiltrating Endometriosis} Endometriosis that grows deeply into tissues around the uterus, bowel, bladder, or other pelvic structures. It may cause severe period pain, pain with sex or bowel movements, urinary symptoms, infertility, and sometimes requires complex specialist treatment.

\medterm{Deep Neck Space Infection} A bacterial infection in one or more of the tissue spaces deep in the neck. It may cause neck swelling, severe throat pain, fever, difficulty swallowing, or breathing problems and can spread to the chest or threaten the airway.

\synonyms
DIE; Deep Endometriosis

\medterm{Deep Sleep} The deepest stage of normal non-REM sleep, when a person is difficult to awaken and the body carries out important physical restoration. Too little may contribute to daytime sleepiness and poor function.
\medterm{Deep Tendon Reflex} An automatic muscle response caused by tapping a tendon with a reflex hammer. It helps assess nerves and the spinal cord. Common examples are the knee, ankle, biceps, and triceps reflexes.

\medterm{Deep Vein Thrombosis} A blood clot in a deep vein, usually in the leg, caused by reduced blood flow, increased clotting tendency, or injury to the vein wall. It may cause one-sided leg swelling, pain, warmth, or discoloration, although some clots cause no symptoms. Part of the clot can break off and travel to the lungs, causing a potentially life-threatening pulmonary embolism.

\synonyms
DVT (Deep Vein Thrombosis)

\medterm{Deep Vein Thrombosis Ultrasound} An ultrasound examination that checks whether deep veins compress normally and whether blood flows through them. A vein that does not compress may contain a clot; results are interpreted with symptoms and other tests.

\medterm{Deep Venous Reflux} Backward flow of blood through weakened valves in deep veins. It can cause long-term leg swelling, heaviness, skin changes, pain, and ulcers from increased pressure in the veins.

\medterm{Deep Venous Thrombosis} Formation of a blood clot in a deep vein, usually in the leg, which can cause swelling and travel to the lungs.

\medterm{De-Escalation} A set of verbal and nonverbal techniques used to reduce agitation and prevent violence
while preserving dignity and safety. It includes calm communication, appropriate
personal space, clear limits, environmental modification, and offering choices.

\medterm{Defaecation} The expulsion of stool from the rectum through the anus. Difficulty, pain, urgency, incontinence, bleeding, or a sustained change in bowel habit may indicate constipation, anorectal disease, neurologic dysfunction, or another condition.

\medterm{Defecation} The process of passing stool from the rectum through the anus. It involves signals from a full rectum, contraction of the bowel, and relaxation of the anal muscles, followed by voluntary control of timing.

\medterm{Defecation Syncope} Fainting during or immediately after defecation, usually a situational reflex involving straining, reduced venous return, and a temporary fall in blood pressure.

\medterm{Defecography} An imaging test that records the rectum and pelvic floor while a person passes a soft contrast material. It can show prolapse, a bulge, inward telescoping, or poor coordination causing difficult stool passage.

\medterm{Defective Virus} A virus missing genetic information needed to reproduce independently. It can multiply only with help from another virus or host system and may influence infection without always causing disease on its own.

\medterm{Defence Mechanisms} Mental processes that help a person manage conflict, anxiety, or distress, often without conscious awareness. They may be helpful or harmful depending on how rigidly they are used and how they affect daily life.

\medterm{Defense Mechanism} A psychological process, often unconscious, that helps a person manage conflict, anxiety, or distress; its usefulness depends on context and flexibility.

\medterm{Defense Mechanisms} Psychological processes, often outside conscious awareness, that help a person manage distress
or conflict. Examples include denial, projection, repression, humor, and sublimation.

\medterm{Defense Mechanisms (Forensic)} Injuries to the hands, arms, or other body areas that may occur when a person tries to protect against an assault. Their meaning must be assessed with the history and all other evidence.
\synonyms
Forensic defense injuries

\medterm{Defense Wounds} Injuries to the hands, wrists, arms, or other areas that may occur when a person tries to protect against an attack. Their appearance and location can support a forensic interpretation but cannot prove how an injury happened by themselves.

\medterm{Defense Wounds (Inscribed)} Cuts or other injuries on the hands or forearms that may occur when a person attempts to protect against an attack.
The distribution and morphology of defensive injuries may help
reconstruct an assault but are not independently proof of a particular event.

\synonyms
Defensive wounds

\medterm{Defenses} The body’s protective mechanisms against infection, injury, and other threats, including physical barriers and immune responses.

\synonyms
Body defenses

\medterm{Defensin} A small protective protein made by skin, lining tissues, and immune cells. It can damage microbes and is part of the body’s early defense against infection.

\medterm{Deferasirox} An oral iron-chelating medicine used to reduce chronic iron overload when clinically indicated; kidney, liver, hearing, and gastrointestinal toxicity require monitoring.

\medterm{Deferiprone} An oral iron chelator used for transfusional iron overload; neutropenia or agranulocytosis is an important adverse effect requiring blood-count monitoring.

\medterm{Deferoxamine} A parenteral iron chelator used for severe or chronic iron overload and sometimes aluminum toxicity; treatment requires monitoring for ocular, auditory, renal, and infectious complications.

\medterm{Defibrillation} Delivery of an electrical shock to stop ventricular fibrillation or pulseless ventricular tachycardia during cardiac arrest. The shock can restore an organized rhythm and must be combined with high-quality CPR and emergency care.

\medterm{Defibrillator} A device that delivers an electrical shock to treat certain life-threatening abnormal heart rhythms. External defibrillators are used during cardiac arrest; implanted devices can detect and treat dangerous rhythms automatically.

\medterm{Defibrillator Implant} An implanted device that continuously monitors heart rhythm and can use pacing or an electrical shock to stop dangerous fast rhythms. It may be recommended for selected people at high risk of sudden cardiac arrest.
\medterm{Deficiency Diseases} Illnesses caused by too little of a needed nutrient or other substance, such as iron, iodine, vitamin D, or vitamin B12. Symptoms depend on the missing substance.

\medterm{Deficiency, FAO} An inherited disorder in which cells cannot use certain fats normally for energy. Fasting or illness may cause low blood sugar, vomiting, muscle breakdown, weakness, heart problems, or sudden metabolic illness.

\medterm{Deficiency, Iron} Inadequate iron stores or availability, which can impair hemoglobin production and cause iron-deficiency anemia, fatigue, pallor, and reduced exercise tolerance.

\medterm{Deficiency Primary Immunodeficiency (Primary Immunodeficiency Disease PIDD)} Alternate index label for Primary Immunodeficiency Disease PIDD.

\medterm{Deficiency, UDP-Glucuronosyltransferase} Reduced activity of a liver enzyme that prepares bilirubin for removal from the body. Inherited deficiency can cause intermittent mild jaundice or, in severe forms, dangerous bilirubin buildup.

\medterm{Deficient Color Vision} Alternate index label; see \textit{Color blindness}.


\medterm{Deformation} An alteration in the shape or structure of a body part, tissue, or material caused by developmental forces, injury, disease, pressure, or mechanical stress.

\medterm{Deformity} An abnormal shape, alignment, or structure of a body part caused by development, injury, disease, pressure, or unusual growth. It may affect appearance, movement, comfort, or organ function.

\medterm{Degarelix} An injected hormone-blocking medicine that rapidly lowers testosterone. It is used for advanced prostate cancer; injection-site reactions, hot flushes, bone loss, and other hormone-related effects may occur.

\medterm{Degeneration} Progressive structural or functional deterioration of cells, tissues, or organs caused by aging, disease, injury, inherited factors, or other processes.

\medterm{Degenerative Arthritis} Alternate index label; see \textit{Osteoarthritis}.

\medterm{Degenerative Disc} Degenerative disc disease describes age- or wear-related changes in spinal discs that may cause neck or back pain, stiffness, or nerve compression. Imaging changes can occur without symptoms; treatment usually begins with activity, exercise-based rehabilitation, and pain management.

\synonyms
DDD (Degenerative Disc)

\medterm{Degenerative Disc Disease} Wear, drying, or cracking of the cushioning discs between spinal bones, often related to aging or injury. It may cause back or neck pain, stiffness, or nerve pressure; imaging changes can occur without symptoms.

\medterm{Degenerative Disease} A condition in which a body tissue or organ gradually becomes damaged or loses function. Causes include ageing, inherited problems, injury, and other disease; examples include osteoarthritis and some disorders of the brain or nerves.

\synonyms
Degenerative disorder

\medterm{Degenerative Disk Disease} Alternate spelling; see \textit{Degenerative Disc}.

\medterm{Degenerative Disorder} A disease in which progressive cellular or tissue deterioration causes loss of structure or function, as in neurodegenerative, joint, retinal, or muscle disorders.

\medterm{Degenerative Joint Disease} Gradual loss of joint cartilage with changes in nearby bone and tissues, commonly meaning osteoarthritis. It can cause pain, stiffness, swelling, reduced movement, and difficulty with daily activities.

\synonyms
Osteoarthritis

\medterm{Degenerative Spine Disease} Wear-related changes in the spinal discs, joints, or bones that can cause back or neck pain, stiffness, reduced movement, or pressure on nerves. Symptoms may include leg or arm pain, numbness, weakness, or walking difficulty.

\medterm{Degenerative Spondylolisthesis} Forward slipping of one spinal bone over another because joints and disks have worn down, sometimes causing back pain, leg pain, or nerve compression.

\medterm{Deglutition} Swallowing is the coordinated movement of food, liquid, or saliva from the mouth through the throat and esophagus into the stomach. It protects the airway and begins digestion.

\medterm{Deglutition (Swallowing)} The coordinated process of moving food, liquid, or saliva from the mouth through the throat and esophagus into the stomach. Problems can cause choking, coughing, or food entering the airway.

\synonyms
Swallowing

\medterm{Degranulation} Release of stored chemicals from cells such as mast cells, white blood cells, and platelets. The chemicals help respond to injury or infection but can also cause allergy, inflammation, blood-vessel changes, or clotting.

\medterm{Degrees Fahrenheit} A unit on the Fahrenheit temperature scale. In clinical care, body temperature may be reported in Fahrenheit; 98.6 degrees Fahrenheit is an approximate average normal oral temperature, not a universal threshold.

\medterm{Dehiscence} Partial or complete separation of a healing wound or surgical closure. It may involve only the skin or deeper tissues; increasing pain, drainage, or exposed tissue requires prompt medical assessment.

\medterm{DEHP} Di(2-ethylhexyl) phthalate, a chemical added to some plastics to make them flexible. Repeated or high exposure may affect reproduction, development, or hormone function, so medical exposure is limited when possible.

\medterm{Dehydration} Loss of more water and electrolytes than are replaced, often from vomiting, diarrhea, fever, sweating, or poor fluid intake. It can cause thirst, dry mouth, reduced urination, dizziness, fast heartbeat, abnormal electrolyte levels, or shock.

\medterm{Dehydration In Children} Loss of more body water than a child replaces, often from diarrhea, vomiting, fever, sweating, or poor intake. Warning signs include dry mouth, few tears, reduced urination, unusual sleepiness, dizziness, or rapid breathing.


\medterm{Dehydroascorbic Acid} The oxidized form of vitamin C. Cells can take it up and convert it back to active vitamin C, allowing it to participate in antioxidant and other chemical reactions.

\medterm{Dehydroepiandrosterone} A hormone made mainly by the adrenal glands that the body can convert into other sex hormones. Levels change with age and illness; supplements may cause acne, unwanted hair, mood changes, or medicine interactions.

\medterm{Deinstitutionalize} To help a person leave long-term institutional care and live in the community with suitable housing, healthcare, and support. The plan should protect safety, treatment needs, and independence.

\medterm{Deiodinase} An enzyme that removes iodine from thyroid hormones and changes how active they are. Different tissues use these enzymes to help control metabolism, growth, brain development, and responses to thyroid hormone.

\medterm{Deja Vu} A brief feeling that a new situation has happened before. Occasional episodes are common, but frequent episodes with loss of awareness, unusual movements, or confusion may occur with seizures and need assessment.

\medterm{Dejerine-Roussy Syndrome} Persistent burning, aching, or painful abnormal sensation on one side of the body after a stroke damages a deep brain sensory area. Normally harmless touch or temperature may become painful.

\medterm{Delavirdine} An older HIV medicine that blocks reverse transcriptase, an enzyme HIV needs to copy its genetic material. It was used with other HIV medicines, but resistance and drug interactions limit its use.

\medterm{Delayed Cord Clamping} Waiting, commonly 30–60 seconds or longer, before clamping the umbilical cord after birth. This allows extra blood to reach the newborn and may improve early circulation and iron stores; timing depends on both patients’ condition.

\medterm{Delayed Ejaculation} Persistent difficulty or inability to ejaculate despite adequate sexual stimulation, causing distress. Medicines, nerve or hormone problems, chronic illness, psychological factors, and relationship circumstances may contribute.

\medterm{Delayed Emergence} Failure to wake and respond as expected after anesthesia. Causes include lingering anesthetic medicines, low body temperature, abnormal blood sugar or salts, breathing problems, stroke, or other complications and require prompt assessment.

\medterm{Delayed Gastric Emptying} Alternate index label; see \textit{Gastroparesis}.

\medterm{Delayed Gastric Emptying (Gastroparesis)} Alternate index label for Gastroparesis.

\medterm{Delayed Growth} Slower-than-expected increase in a child’s height, weight, or head size. Nutrition, chronic disease, hormone problems, genetic conditions, and family growth pattern can contribute; growth should be assessed over time.

\medterm{Delayed Hypersensitivity} An immune reaction that develops hours or days after exposure rather than immediately. It can cause contact dermatitis, some medicine rashes, or a reaction to a tuberculosis skin test.

\medterm{Delayed Orgasm} Persistent difficulty or inability to reach orgasm despite adequate stimulation, causing distress. Medicines, illness, hormonal or nerve problems, anxiety, trauma, and relationship factors may contribute.

\medterm{Delayed Puberty} Puberty that has not begun or is not progressing at the expected age. Family pattern, poor nutrition, chronic illness, hormone disorders, or problems involving the ovaries, testes, brain, or pituitary may contribute.

\medterm{Delayed Puberty In Boys} Lack of testicular enlargement or other expected pubertal changes by the usual age. Causes include normal late family development, poor nutrition, chronic illness, low testosterone, or hypothalamus or pituitary disorders.

\medterm{Delayed Puberty In Girls} Lack of breast development or menstruation by the usual age, or puberty that stops progressing. Causes include normal late family development, low energy availability, chronic illness, ovarian failure, or hypothalamus or pituitary disorders.

\medterm{Delayed Sleep Phase} A body-clock pattern in which a person naturally falls asleep and wakes much later than desired. Sleep may be normal at the preferred times but interfere with school, work, or social activities.

\synonyms
delayed sleep-wake phase disorder.
Delayed Sleep-Wake Phase Sleep Disorder
DSPS
DSWPD

\medterm{Delayed Sleep Phase Syndrome} A persistent body-clock disorder in which sleep begins and ends much later than desired, despite the ability to sleep normally on that schedule. It causes difficulty meeting work, school, or social obligations.

\medterm{Delayed Sleep-Wake Phase Sleep Disorder} Alternate index label; see \textit{Delayed sleep phase}.

\medterm{Deletion} Loss of a section of DNA or part of a chromosome. The effects depend on the size, location, and genes involved; deletions may cause no problems or lead to inherited, developmental, or medical disorders.

\medterm{Deletion Mutation} A genetic change in which one or more DNA building blocks are missing. Its effect depends on the location and size of the deletion and whether it changes a gene or its control region.

\medterm{Delftia} A group of oxygen-using bacteria found mainly in soil and water. Human infection is uncommon but can occur as an opportunistic infection, especially in people with weakened immunity or medical devices.

\medterm{Delirium} A sudden change in attention, awareness, and thinking that fluctuates over hours or days. It usually results from illness, infection, medicines, substance use or withdrawal, dehydration, or another medical problem.

\medterm{Delirium Superimposed On Dementia} Delirium occurring in a person who already has dementia. The sudden, fluctuating worsening in attention and awareness differs from the person’s usual cognitive impairment and may be triggered by illness, medicines, dehydration, or hospitalization.

\medterm{Delirium In Anesthesia} Sudden confusion, poor attention, and changing awareness during recovery from anesthesia or surgery. Possible causes include medicines, infection, pain, low oxygen, dehydration, or abnormal body chemistry and require prompt assessment.

\medterm{Delirium Tremens} A life-threatening form of alcohol withdrawal causing severe confusion, agitation, hallucinations, shaking, sweating, fever, fast heartbeat, and high blood pressure. It commonly begins two to four days after heavy drinking stops and requires emergency treatment.

\medterm{Delivery} Childbirth, ending with birth of the baby and delivery of the placenta. It may occur through the vagina, sometimes with forceps or vacuum assistance, or through cesarean surgery.

\medterm{Delivery Of Posterior Arm} A maneuver for shoulder dystocia in which the clinician brings the baby’s back arm through the birth canal. This reduces the width of the shoulders and may help complete delivery.

\medterm{Delivery Presentations} The fetal body part positioned to enter the birth canal first. Common presentations are head-first, breech, and shoulder; the presentation affects delivery planning and risks.

\medterm{Delta Agent} An older name for hepatitis D virus, which can reproduce only when hepatitis B virus is present. Infection may cause severe acute hepatitis or worsen long-term liver disease.

\medterm{Delta Waves} Very slow electrical brain waves that are normal during deep non-REM sleep. Excessive delta activity while awake may indicate brain dysfunction and must be interpreted with symptoms and other tests.

\synonyms
Delta rhythm

\medterm{Delta-ALA Urine Test} A urine test measuring a substance used to make hemoglobin. Increased levels may occur with acute porphyria or lead poisoning and are interpreted with symptoms and other tests.

\medterm{Deltacoronavirus} A group of coronaviruses identified by their genetic features. Some infect birds or mammals, and their ability to infect people or cause disease depends on the specific virus and host.

\medterm{Deltaretrovirus} A group of viruses that includes human T-cell leukemia viruses. They copy their RNA into DNA inside infected cells; some can persist for years and cause blood-cell disorders or cancer.

\medterm{Delta-Tocopherol} One form of vitamin E found in foods and supplements. It has antioxidant activity, but alpha-tocopherol is the main form maintained in human blood and used to assess vitamin E status.

\medterm{Deltoid} The large, triangular muscle covering the shoulder. It helps lift the arm and stabilize the shoulder joint; injury or nerve damage can cause pain and weakness.

\medterm{Deltoid Muscle} A triangular shoulder muscle attached to the collarbone, shoulder blade, and upper arm. Its front, middle, and back fibers move the arm forward, outward, and backward; the axillary nerve supplies it.

\medterm{Delusion} A firmly held false belief that remains despite clear evidence to the contrary and is not explained by the person’s cultural background. It may involve persecution, grandiosity, illness, or other themes.

\medterm{Delusional Disorder} A mental disorder with one or more persistent delusions lasting at least one month, while overall functioning is relatively preserved. Hallucinations, if present, relate closely to the delusional belief.

\medterm{Demarcation} A clear boundary between two areas or conditions, such as the line separating healthy tissue from tissue that has died. It may guide diagnosis or treatment.

\medterm{Demeclocycline} A tetracycline antibiotic that blocks bacterial protein production. It was also used to reduce excess water retention in some patients with inappropriate antidiuretic hormone secretion; adverse effects include sun sensitivity and kidney or liver problems.

\medterm{Demecolcine} A colchicine-related substance that disrupts microtubules, structures needed for cell division. It has limited medical use and can be highly toxic, causing gastrointestinal, blood-cell, nerve, or organ injury.

\medterm{Dementia} A lasting decline in memory, thinking, language, judgment, or other mental abilities that interferes with independent daily life. Causes include brain diseases, blood-vessel disease, infections, medicines, and some treatable medical problems.

\synonyms
Major neurocognitive disorder

\medterm{Dementia And Driving} The effect of cognitive impairment on driving safety. Assessment considers memory, judgment, visuospatial skills, reaction time, medications, and local licensing requirements.

\medterm{Dementia - Behavior And Sleep Problems} Behavioral and sleep changes that can occur in dementia, including agitation, wandering, apathy, hallucinations, or altered sleep–wake timing; causes and reversible contributors should be assessed.


\medterm{Dementia Due To Metabolic Causes} Cognitive impairment caused or worsened by metabolic problems such as thyroid disease, vitamin deficiency, organ failure, electrolyte disturbance, or medication toxicity; some causes are treatable.

\medterm{Dementia, Frontotemporal} Alternate index label; see \textit{Frontotemporal dementia}.

\medterm{Dementia, Lewy Body} Alternate index label; see \textit{Lewy body dementia}.

\medterm{Dementia (Major Neurocognitive Disorder)} Alternate name; see \textit{Dementia}.

\medterm{Dementia Pugilistica} A progressive brain disorder associated with repeated head trauma, historically described in some professional boxers.
\medterm{Dementia, Vascular} Alternate index label; see \textit{Vascular dementia}.

\medterm{Demethylation} Removal of a small chemical group called a methyl group from a molecule such as DNA or a protein. In DNA, this change can alter which genes are active and affect cell development or disease.

\medterm{Demineralization} Loss of minerals from a tissue such as tooth enamel or bone. It can weaken the tissue, causing early tooth decay or contributing to fragile bones, depending on the cause.


\medterm{Demulcent} A substance that coats and soothes irritated moist body linings, temporarily reducing discomfort from dryness, inflammation, or minor injury.

\medterm{Demyelinating Disease} A disease in which the protective covering around nerve fibers is damaged. Nerve signals then slow or fail, causing problems such as weakness, numbness, vision changes, or poor coordination.

\medterm{Demyelinating Disease Of The Central Nervous System} A disorder in which the protective covering of nerves in the brain, spinal cord, or vision pathways is damaged. Causes include immune disease, infection, inherited conditions, toxins, and metabolic problems.

\medterm{Demyelination} Loss or damage of myelin, the protective coating that helps nerve fibers send signals quickly. It can disrupt movement, sensation, vision, balance, or thinking and may result from immune, infectious, inherited, or toxic disorders.

\medterm{Denaturant} A substance added to another product to change its taste, smell, or properties, often to make it unsuitable for drinking. Some denaturants are poisonous or irritating, so health effects depend on the product and amount swallowed or contacted.

\medterm{Dendrite} A branching extension of a nerve cell that receives signals from other nerve cells or sensory receptors and carries those signals toward the cell body.

\medterm{Dendrites} Branching extensions of nerve cells that receive signals from other neurons or sensory receptors and carry those messages toward the cell body.

\medterm{Dendritic Cell} An immune cell that captures foreign material and shows pieces of it to other immune cells. This helps start a targeted immune response against infections and abnormal cells.

\medterm{Dendritic Cells} Immune cells that collect material from germs or damaged cells and show identifying pieces to T cells. This helps the immune system start a targeted response and develop longer-lasting protection.

\medterm{Dendritic Spine} A tiny branch on a nerve cell that receives signals from other nerve cells. Its shape and number can change with learning, development, and some brain diseases.

\medterm{Dendritic Ulcer} A branching epithelial ulcer of the cornea, classically caused by herpes simplex virus. It may cause a red eye, pain, light sensitivity, tearing, or reduced vision and requires ophthalmic assessment.

\medterm{Denervate} To remove or interrupt the nerve supply to a tissue or organ, either deliberately during treatment or because of injury or disease. The change may reduce movement, sensation, pain, or automatic organ function.

\medterm{Denervation} Denervation is loss of nerve supply to a muscle or tissue, causing weakness, atrophy, altered sensation, or autonomic dysfunction.

\medterm{Dengue} A viral infection caused by a dengue virus and spread mainly by the bites of infected female \textit{Aedes} mosquitoes. Many infections cause no symptoms; symptomatic illness commonly causes sudden fever, severe headache, pain behind the eyes, muscle or joint pain, nausea, and rash. A small number of cases progress to severe dengue, which can cause leakage of fluid from blood vessels, fluid buildup with breathing difficulty, severe bleeding, serious organ impairment, shock, or death.

\medterm{Dengue Fever} Alternate name; see \textit{Dengue}.

\synonyms
Dengue Hemorrhagic Fever


\medterm{Dengue Hemorrhagic Fever} A severe form of dengue infection in which blood vessels leak and bleeding or shock may develop. Warning signs include severe abdominal pain, repeated vomiting, bleeding, extreme tiredness, cold skin, or reduced urine; urgent hospital care and careful fluids are needed.
\medterm{Dengue Shock Syndrome} Life-threatening circulatory shock caused by severe dengue with plasma leakage and reduced effective blood volume, often accompanied by bleeding or organ dysfunction.

\medterm{Dengue Virus} The virus that causes dengue fever and is spread mainly by infected Aedes mosquitoes. Four major types exist; protection after infection is strongest against the same type and does not reliably prevent later infection by other types.

\medterm{Denial} A psychological response in which a person cannot or will not accept a distressing fact, diagnosis, or situation. It may briefly reduce distress but can delay needed decisions or treatment when persistent.

\medterm{Denileukin Diftitox} A laboratory-made protein combining part of interleukin-2 with diphtheria toxin. It was used for selected cutaneous T-cell lymphomas by delivering the toxin to cells bearing the interleukin-2 receptor; toxicity limits use.

\medterm{Denominator} The number below the line in a fraction. In a medical rate or proportion, it represents the total population or group being compared with the number above it.

\medterm{Denosumab} A medicine that blocks a protein needed to activate bone-destroying cells. It is used for osteoporosis and some cancer-related bone problems; stopping it without follow-up can cause rapid bone loss and spine fractures.

\medterm{Denosumab For Osteoporosis} An injection given every six months that blocks a signal activating bone-destroying cells. It lowers fracture risk, including at the spine and hip, but calcium must be monitored and another treatment is needed when it is stopped.

\medterm{Dens} The tooth-shaped projection of the second neck vertebra. It fits into the first neck vertebra and acts as a pivot that allows the head to turn.

\medterm{Dens Evaginatus} An extra bump or cusp on a tooth, most often a premolar. It may contain an extension of the tooth’s nerve and can wear, break, or develop decay and pain.

\medterm{Dens Invaginatus} An inward folding of a tooth’s enamel and dentin that may reach the nerve. The pocket can trap bacteria and food, increasing the risk of early decay, nerve infection, and tooth loss.

\medterm{Densitometry} Measurement of how dense a tissue is, most often the mineral content of bone. A bone-density scan helps diagnose osteoporosis, estimate fracture risk, and monitor treatment.

\medterm{Density} The amount of mass in a given volume. In medicine, density measurements help describe bone mineral content, tissue composition, fluid collections, or findings on imaging.

\medterm{Dent Disease} An inherited disorder of the kidney’s filtering tubes, usually affecting boys. It causes loss of small proteins in urine, excess urinary calcium, kidney stones or calcium deposits, and sometimes progressive kidney failure.

\medterm{Dental Abscess} A pocket of pus caused by infection in a tooth or surrounding gum. It may cause severe toothache, swelling, fever, or difficulty opening the mouth; spreading facial or neck swelling needs urgent care.

\medterm{Dental Age} An estimate of a child’s developmental age based on tooth formation, hardening, and eruption. It can help assess growth but does not exactly determine the child’s actual age.

\medterm{Dental Alloy} A mixture of metals used to make dental fillings, crowns, bridges, dentures, orthodontic appliances, or other restorations. Its strength, appearance, corrosion resistance, and possible allergies depend on its composition.

\medterm{Dental Alveolus} The part of the upper or lower jawbone that forms a tooth socket and holds the tooth root. It can shrink after tooth loss or severe gum disease.

\medterm{Dental Amalgam} A durable filling material made by mixing liquid mercury with powdered metals, mainly silver, tin, and copper. It is strong but has a silver color and may be unsuitable in some situations.

\medterm{Dental Anesthesia} Use of local anesthetic, sedation, or occasionally general anesthesia to prevent pain or reduce anxiety during dental treatment. The method depends on the procedure, health, age, and patient preference.

\medterm{Dental Anxiety} Fear or distress about dental treatment that may cause avoidance, delayed care, or difficulty during appointments. Reassurance, communication, pain control, and behavioral or medicine-based support may help.

\medterm{Dental Avulsion} Complete displacement of a tooth from its socket, usually caused by an injury. A permanent tooth that is knocked out needs immediate dental care; hold it by the crown, keep it moist, and do not scrub the root.

\medterm{Dental Arch} The curved arrangement of teeth and the bone supporting them in the upper or lower jaw. Its shape affects tooth alignment, bite, speech, chewing, and space for dental restorations.

\medterm{Dental Articulator} A device that holds models of the upper and lower teeth and imitates jaw movements. It helps dental professionals design and adjust crowns, bridges, dentures, and other restorations so they fit the bite.

\medterm{Dental Braces} Fixed or removable appliances that apply gentle, controlled force to move teeth and improve their alignment or bite. Treatment requires cleaning and follow-up; possible problems include soreness, decay, gum disease, and root shortening.

\synonyms
Braces And Retainers (Dental Braces)

\medterm{Dental Braces (Orthodontics)} Fixed or removable appliances that gradually move teeth and may improve the bite or tooth alignment. Treatment length and the need for extraction or other appliances depend on the teeth, jaws, growth, and treatment goals.

\medterm{Dental Bridge} A fixed dental appliance that replaces one or more missing teeth by attaching artificial teeth to nearby teeth or implants. It can improve chewing and appearance but requires healthy supports and careful cleaning.

\medterm{Dental Calculus} Hardened plaque formed when minerals in saliva or gum fluid collect on teeth. It cannot be removed reliably by brushing alone and can contribute to gum inflammation and disease.

\medterm{Dental Caries} Tooth decay caused when bacteria use sugars and produce acids that remove minerals from enamel and dentin. Untreated decay can cause sensitivity, pain, infection, and tooth loss.
\medterm{Dental Caries (Cavities)} Alternate index label for Cavities.

\medterm{Dental Cavities} Permanent holes or weakened areas caused when plaque bacteria use sugars and make acids that dissolve a tooth. Decay can spread from enamel into the deeper tooth, causing sensitivity, pain, infection, or tooth loss.

\medterm{Dental Cement} A dental material used to attach crowns, bridges, or orthodontic bands, fill small spaces, or protect exposed tooth tissue. Its strength and use depend on the material and treatment.

\medterm{Dental Ceramic} A tooth-colored material used for crowns, veneers, bridges, fillings, and implant restorations. It provides good appearance and resistance to wear, but its strength and suitability vary by type and location.

\medterm{Dental Clasp} A part of a removable partial denture that rests against or around supporting teeth to help keep the denture in place. Poor fit or excessive pressure can irritate teeth or gums.

\medterm{Dental Crown} A custom-made covering placed over a damaged or weakened tooth to restore its shape, strength, chewing function, or appearance. It may be made from ceramic, metal, resin, or a combination of materials.

\medterm{Dental Crowns} Custom-made coverings placed over most or all of a damaged or weakened tooth. A crown restores shape, strength, chewing function, and appearance after decay, fracture, or root-canal treatment.

\medterm{Dental Deposit} Material that collects on a tooth, such as plaque, hardened tartar, food debris, or stain. Some deposits irritate the gums or support decay and may require professional cleaning.

\medterm{Dental Disease} A disorder of the teeth, gums, jaw, or supporting tissues, including decay, gum disease, infection, and tooth loss. Fluoride, daily cleaning, limited frequent sugar exposure, and dental visits reduce risk.

\medterm{Dental Dysplasia} Abnormal development of a tooth or its supporting tissue. It may change the tooth’s shape, color, strength, position, or eruption and can increase the risk of wear, decay, or infection.

\medterm{Dental Erosion} Gradual loss of tooth enamel caused by acid rather than bacteria. Acidic foods or drinks, reflux, vomiting, and some medicines may contribute, causing sensitivity or worn teeth; reducing acid exposure helps prevent progression.

\medterm{Dental Extraction} Removal of a tooth from its jaw socket because of severe decay, infection, gum disease, impaction, crowding, trauma, or treatment planning. Risks include bleeding, infection, nerve injury, fracture, and dry socket.
\medterm{Dental Drill} A powered dental handpiece with a rotating bur, driven by air or an electric motor, used to remove decayed tooth structure, shape teeth, prepare teeth for fillings or crowns, and perform other dental procedures. Different burs and speeds are selected for the material and procedure.

\synonyms
Dental Handpiece

\medterm{Dental Fillings} Materials placed in a tooth after decay or damage is removed to restore its shape and function. Resin, glass ionomer, metal amalgam, or other materials may be used depending on the tooth and clinical situation.

\medterm{Dental Fissure} A narrow natural groove or developmental defect in a tooth, especially on a chewing surface. Food and bacteria can collect there, increasing the risk of decay.

\medterm{Dental Fluorosis} Changes in tooth enamel caused by too much fluoride while the teeth are developing. It may appear as white lines or patches and, in more severe cases, brown discoloration or pits.

\medterm{Dental Floss} A thin thread used to remove plaque and food between teeth and just below the gumline. Daily cleaning between teeth helps prevent gum inflammation but does not replace brushing or dental care.

\medterm{Dental Hygiene} Daily care of the teeth, gums, tongue, and mouth, including brushing with fluoride toothpaste and cleaning between teeth. It removes plaque and helps prevent decay, gum disease, and bad breath.

\medterm{Dental Implant} A small biocompatible fixture placed in the jawbone to support a replacement tooth, bridge, or denture. Bone may grow around it for stability; risks include infection, poor healing, nerve injury, and implant failure.

\medterm{Dental Impression Material} A material placed in a tray to make an accurate mold of the teeth and mouth tissues. The mold helps make crowns, dentures, orthodontic appliances, restorations, or study models.

\medterm{Dental Injuries} Damage to teeth, gums, jawbone, or mouth tissues caused by a fall, blow, sports injury, accident, or biting force. Prompt assessment may save displaced or knocked-out teeth; severe bleeding or head injury is an emergency.

\synonyms
Dental Trauma (Dental Injuries)

\medterm{Dental Inlay} A custom-made filling fitted inside the grooves of a damaged tooth and bonded in place. It restores areas too large for a simple filling but does not cover the tooth’s cusps like an onlay or crown.

\medterm{Dental Material} A substance used in dental care, such as composite, amalgam, ceramic, metal, cement, bonding material, impression material, or implant material. Choice depends on the treatment, location, strength, appearance, and patient factors.

\medterm{Dental Occlusion} The way the upper and lower teeth meet when the mouth closes or the jaw moves. An abnormal bite may contribute to difficulty chewing, tooth wear, jaw pain, or muscle strain.

\medterm{Dental Overjet} The horizontal distance between the upper and lower front teeth when the teeth are closed. A large overjet may increase the risk of trauma to the upper front teeth and affect biting.

\medterm{Dental Pain (Toothache)} Pain from a tooth or nearby gum and jaw tissues. Common causes include decay, inflammation or infection inside the tooth, a crack, trauma, gum disease, or an impacted tooth.

\medterm{Dental Papilla} Soft embryonic tissue beneath a developing tooth. Its cells form the dentin and living pulp, which contains the tooth’s nerves and blood vessels.

\medterm{Dental Pellicle} A thin protective film of salivary proteins that forms on tooth enamel soon after cleaning. Bacteria can attach to it and build dental plaque.

\medterm{Dental Plaque} A sticky film of bacteria, saliva, and food material on teeth and along the gums. If not removed by brushing and cleaning between teeth, it can cause decay, gum inflammation, and gum disease.

\medterm{Dental Porcelain} A tooth-colored ceramic used for crowns, veneers, bridges, and other restorations. It provides a natural appearance but may chip or fracture under heavy force.

\medterm{Dental Procedure} An examination, preventive service, or treatment involving the teeth, gums, jaw, or mouth, such as cleaning, filling, extraction, root-canal treatment, or restoration. Benefits and risks depend on the procedure.

\medterm{Dental Prosthesis} An artificial replacement for one or more missing teeth or nearby mouth tissues, such as a denture, bridge, crown, or implant-supported restoration. It restores appearance, speech, and chewing.

\medterm{Dental Public Health} Efforts to prevent and control oral disease in communities. They include fluoride programs, oral-health education, screening, disease monitoring, access to dental care, and policies that reduce differences in oral health.

\medterm{Dental Pulp} The living soft tissue at the center of a tooth. It contains nerves, blood vessels, and connective tissue and can become inflamed or infected when decay or injury reaches it.

\medterm{Dental Pulp Necrosis} Death of the soft tissue inside a tooth, usually after untreated inflammation, deep decay, trauma, or loss of blood supply. The tooth may become discolored and develop infection around its root.

\medterm{Dental Pulpitis} Inflammation of the soft tissue inside a tooth, commonly caused by decay, trauma, a crack, or dental treatment. Brief pain may be reversible; spontaneous or lingering pain may require root-canal treatment or extraction.

\medterm{Dental Radiographs} X-ray images of the teeth and jaws used to find decay, bone loss, infection, impacted teeth, fractures, and developmental problems. Types include bitewing, periapical, occlusal, and panoramic views.

\medterm{Dental Radiography} Use of X-rays to examine teeth and jaws. Bitewing, periapical, panoramic, and three-dimensional scans answer different questions; exposure is kept as low as reasonably achievable while obtaining useful images.

\medterm{Dental Restorative Materials} Materials used to repair teeth, including composite resin, glass ionomer, amalgam, ceramic, and metal alloys. Selection depends on tooth location, size of damage, strength, appearance, moisture control, cost, and patient needs.

\medterm{Dental Sac} A layer of tissue surrounding a developing tooth. It forms the tissues that support the tooth, including the periodontal ligament, cementum, and part of the surrounding jawbone.

\medterm{Dental Screening} A check of the teeth, gums, and mouth to look for decay, gum disease, oral lesions, or risk factors before serious symptoms develop. Screening does not replace a complete dental examination.

\medterm{Dental Sealant} A thin protective coating placed in the grooves of teeth, especially molars. It blocks food and bacteria from collecting in deep pits and lowers the risk of decay.

\medterm{Dental Sealants} Thin resin or glass-ionomer coatings placed in the pits and grooves of teeth, especially molars. They reduce plaque and food retention and lower the risk of decay but do not replace brushing or fluoride.

\medterm{Dental Trauma (Dental Injuries)} Alternate index label for Dental Injuries.

\medterm{Dental Traumatology} The prevention, assessment, and treatment of injuries to teeth and surrounding mouth tissues. Injuries may involve cracks, fractures, displacement, root damage, or complete tooth loss and often need prompt care.

\medterm{Dental X-Rays} X-ray images of teeth and jaws used to detect decay, bone loss, infection, impacted teeth, fractures, developmental problems, and other disease not visible during an examination.

\medterm{Dentate Gyrus} A curved part of the hippocampus that helps process incoming information and form memories. It also produces some new nerve cells during adulthood.

\medterm{Dentatorubral-Pallidoluysian Atrophy} A rare inherited brain disorder caused by an extra-long repeated section of the ATN1 gene. Symptoms vary with age and may include seizures, poor coordination, involuntary movements, behavior changes, and worsening thinking ability.

\synonyms
DRPLA; Haw River Syndrome; Naito-Oyanagi Disease

\medterm{Dentifrice} A preparation, such as toothpaste or tooth powder, used to clean teeth. Fluoride dentifrice strengthens enamel and helps prevent tooth decay when used regularly.

\medterm{Dentigerous Cyst} A fluid-filled sac that forms around the crown of an unerupted tooth, usually in the jaw. It may cause swelling, move nearby teeth, or damage bone if it becomes large or infected.

\medterm{Dentin} The hard tissue beneath tooth enamel that forms most of the tooth. Tiny channels within it can transmit sensitivity, especially when dentin is exposed by decay, wear, gum recession, or injury.

\medterm{Dentin Sensitivity} Brief, sharp tooth pain caused by cold, heat, sweet foods, touch, or air reaching exposed dentin. Decay, worn enamel, gum recession, cracks, or aggressive brushing can expose it.

\medterm{Dentin Hypersensitivity} Alternate term; see \textit{Dentin Sensitivity}.

\medterm{Dentinogenesis} Formation of dentin by specialized cells during tooth development and repair. Abnormal dentin formation may produce weak, discolored, sensitive teeth that wear or fracture easily.

\medterm{Dentinogenesis Imperfecta} An inherited disorder in which dentin forms abnormally, causing teeth that may look blue-gray or amber and break, wear, or discolor easily. Some forms occur with fragile bones.

\medterm{Dentist} A licensed healthcare professional who examines, prevents, diagnoses, and treats diseases and injuries of the teeth, gums, jaw, and mouth.

\medterm{Dentistry} The healthcare field concerned with preventing, diagnosing, and treating disorders and injuries of the teeth, gums, jaw, and mouth, including restoration, replacement, and oral health education.

\medterm{Dentistry (Geriatric)} Dental care for older adults, who may have dry mouth, root decay, gum disease, tooth loss, denture problems, or oral cancer risk. Care includes medication review, prevention, examination, treatment, and adaptations for mobility or cognition.

\synonyms
Geriatric dentistry

\medterm{Dentition} The teeth present in the mouth, including their number, arrangement, development, and eruption. Children usually have primary teeth that are replaced by permanent teeth; disease, injury, or development can alter either set.

\medterm{Denture} A removable replacement for some or all missing teeth and nearby mouth tissues. Partial dentures replace selected teeth; complete dentures replace all teeth in one jaw. Correct fit, cleaning, and review help prevent sores and infection.

\medterm{Denture Base} The part of a denture that rests on the gums and mouth tissues and supports the artificial teeth. Its shape and fit distribute chewing pressure and help prevent soreness.

\medterm{Denture Problems} Difficulties caused by a removable denture, including poor fit, looseness, pain, sores, breakage, difficulty chewing or speaking, or infection. A dental review may correct the cause or determine whether replacement is needed.

\medterm{Dentures} Removable appliances that replace some or all missing teeth and nearby gum tissue. They can improve chewing, speech, and appearance but need proper fitting, daily cleaning, gradual adjustment, and regular dental review.

\medterm{Denturist} A dental professional authorized in some regions to make, fit, adjust, and repair removable dentures within the local scope of practice. Dentists may also assess oral disease and provide related treatment.

\medterm{Denver Developmental Screening Test} A screening questionnaire and task-based assessment for young children that checks personal-social, fine-motor, language, and gross-motor skills. It identifies possible developmental delay but does not diagnose its cause.

\medterm{Denys-Drash Syndrome} A rare genetic disorder caused by changes in the \textit{WT1} gene. It can cause severe kidney disease, differences in sex development, and a high risk of Wilms kidney tumor in childhood.

\medterm{Deodorant Poisoning} Harmful exposure to a deodorant product, usually by swallowing, inhaling, or getting it in the eyes or on the skin. Symptoms depend on the ingredients and may include irritation, vomiting, breathing difficulty, or chemical burns.

\medterm{Deoxycoformycin} An anticancer medicine that blocks an enzyme needed by some cells to make DNA. It was used for selected blood cancers and can suppress bone marrow, weaken immunity, and increase the risk of serious infections.

\medterm{Deoxycytidine} A DNA building block made of the base cytosine joined to the sugar deoxyribose. Cells use it to make and repair DNA.
\medterm{Deoxycytidine Kinase} An enzyme that chemically activates deoxycytidine and some similar medicines. The activated medicines can interfere with the DNA of cancer cells or viruses and stop them from multiplying.

\medterm{Deoxyglucose} A molecule resembling glucose that is used mainly in research and PET imaging. When labeled with a radioactive form of fluorine, it can show tissues using unusually large amounts of glucose, including many tumors and areas of inflammation.

\medterm{Deoxyguanosine} A DNA building block made of the base guanine joined to the sugar deoxyribose. Cells use it to make and repair DNA.
\medterm{Deoxyhemoglobin} Hemoglobin that is not carrying oxygen. A high amount in blood near the skin can make the lips, skin, or nail beds appear blue, a finding called cyanosis.

\medterm{Deoxyribonuclease} An enzyme that breaks DNA into smaller pieces. An inhaled laboratory-made form can thin thick mucus in some people with cystic fibrosis by breaking down DNA released from damaged cells.

\medterm{Deoxyribonucleic Acid} DNA, the molecule that stores hereditary information in cells. Its sequence helps instruct protein production and cell control; changes may be inherited or acquired during life and can affect health.

\synonyms
DNA

\medterm{Deoxyribonucleosides} Molecules made of one of the four DNA bases joined to the sugar deoxyribose. Cells use them as starting materials when making or repairing DNA.

\medterm{Deoxyribonucleotides} DNA building blocks made of a base, deoxyribose sugar, and phosphate groups. The four forms—dATP, dGTP, dCTP, and dTTP—are assembled into DNA strands.

\medterm{Dependence} Physical adaptation to repeated exposure to a medicine or substance, so that reducing or stopping it can cause withdrawal symptoms. Dependence can occur with or without addiction.

\medterm{Dependence, Nicotine} Alternate index label; see \textit{Nicotine dependence}.

\medterm{Dependent Personality Disorder} A long-term pattern of excessive need to be cared for, difficulty making decisions independently, fear of separation, and submissive behavior that causes distress or impairs relationships and daily life.

\medterm{Dependovirus} A group of small viruses that usually need another virus, such as an adenovirus or herpesvirus, to reproduce. They are generally not known to cause serious disease in healthy people.

\medterm{Depersonalization} A feeling of being detached from or observing one’s own body, thoughts, or emotions as if from outside. It can occur with anxiety, trauma, panic, substance use, or other mental health conditions.

\medterm{Depersonalization-Derealization Disorder} A disorder causing repeated or persistent feelings that one’s body, thoughts, emotions, or surroundings are unreal or distant, while the person knows these feelings are not literally true. Symptoms cause significant distress or impairment.

\medterm{Dephosphorylation} Removal of a phosphate group from a molecule. This can switch proteins and other cell components on or off and change how cells communicate, use energy, grow, or respond to signals.

\medterm{Depigmentation} Loss or reduction of normal pigment, especially melanin in the skin, hair, or mucous membranes. It may result from vitiligo, inflammation, injury, infection, chemical exposure, or treatment.

\medterm{Depilation} Removal of hair at or near the skin surface by shaving, cutting, waxing, or a chemical product. It differs from epilation, which removes hair together with its root.

\medterm{Depilatory Poisoning} Harmful exposure to a hair-removal product, usually by swallowing or contact with the eyes, skin, or airways. Ingredients may cause irritation, chemical burns, vomiting, or breathing problems.

\medterm{Depolarization} A change that makes the inside of a cell less electrically negative. In nerve and muscle cells, it can start or transmit an electrical signal that produces movement or sensation.

\medterm{Depot Implant} A small device placed in the body that releases medicine slowly over an extended period. It can provide steady treatment but may require a procedure for insertion and removal.

\medterm{Deprescribing} A planned, supervised reduction or stopping of a medicine when its harms, burden, or limited benefit outweigh expected benefit. The plan considers withdrawal, return of illness, safer alternatives, and patient goals.

\medterm{Depressed Fracture} A break in which bone is pushed inward, most often in the skull. It can injure the brain or other underlying tissue and requires urgent assessment, especially after head injury or with neurologic symptoms.

\medterm{Depression} A group of mental health disorders involving a persistent sad, empty, or irritable mood, or a loss of interest or pleasure in activities. It is more than the short-lived sadness or low mood that most people experience at times. Other symptoms may include changes in sleep, appetite, energy, concentration, movement, self-worth, or hope, as well as physical aches or digestive problems. Symptoms cause significant distress or interfere with daily life and may include thoughts of death or suicide. Major depression involves symptoms most of the day, nearly every day, for at least two weeks; persistent depressive disorder involves less severe symptoms lasting at least two years.

\synonyms
Clinical Depression (Depression)

\medterm{Depression (Geriatric)} Depression in an older adult may cause low mood, loss of interest, fatigue, sleep or appetite changes, slowed thinking, physical complaints, or memory problems. Medical illness, medicines, isolation, and bereavement may contribute.

\synonyms
Depression in older adults

\medterm{Depression In Children} Depression in a child may appear as persistent sadness or irritability, loss of interest, withdrawal, school decline, sleep or appetite changes, physical complaints, guilt, or thoughts of death. Evaluation includes development, family context, medical causes, and safety.

\medterm{Depression In Older Adults} Depression later in life may cause sadness, loss of interest, fatigue, sleep or appetite changes, slowed thinking, physical complaints, or memory problems. Assessment should also consider medical illness, medicines, grief, isolation, and suicide risk.

\medterm{Depression (Major Depressive Disorder)} Alternate index label; see \textit{Major depressive disorder}.



\medterm{Depression, Postpartum} Alternate index label; see \textit{Postpartum depression}.

\medterm{Depression Screening} Use of questions or a validated questionnaire to identify possible depression. A positive screen is not a diagnosis; a clinician must assess symptoms, duration, impairment, medical causes, treatment needs, and safety.

\medterm{Depression, Teen} Alternate index label; see \textit{Teen depression}.

\medterm{Depressive Disorder} A mood disorder involving persistent low mood, loss of interest, or both, with possible changes in sleep, appetite, energy, concentration, movement, or self-worth. Diagnosis depends on the pattern, duration, severity, and functional impact of symptoms.

\medterm{Depressive Illness} Broad, non-diagnostic term; see \textit{Depressive Disorder}.

\medterm{Depth Perception} The ability to judge how far away objects are and where they are in three-dimensional space. It depends mainly on using both eyes together and also on movement, size, shading, and perspective cues.

\medterm{Depurination} Loss of the DNA bases adenine or guanine, leaving a gap in the DNA strand. If the gap is not repaired correctly, it can cause a genetic mutation.

\medterm{DeQuervains Tenosynovitis (De Quervain's Tenosynovitis)} Alternate index label for De Quervain's Tenosynovitis.

\medterm{Derbyshire Neck} An older term for an enlarged thyroid gland, or goiter, historically associated with areas where iodine deficiency was common. The cause and significance require assessment of the thyroid and its function.

\medterm{Dercum's Disease} A rare disorder characterized by multiple painful fatty growths, often with obesity, fatigue, and widespread tenderness. The cause is uncertain, and diagnosis is based on symptoms and examination after other causes are excluded.

\medterm{Derealization} A feeling that the outside world is unreal, distant, dreamlike, or visually changed, while the person knows it is not literally unreal. It may occur with anxiety, panic, trauma, dissociation, or substance use.

\medterm{Dermabrasion} A procedure that uses a rapidly rotating instrument to remove outer skin layers. It may improve selected scars, wrinkles, sun damage, or lesions but can cause pain, pigment changes, infection, or scarring.

\medterm{Dermal} Relating to the dermis, the skin layer beneath the epidermis. The dermis contains connective tissue, blood vessels, nerves, hair follicles, and sweat glands that support and nourish the skin.

\medterm{Dermal Fillers} Injectable materials placed under or within the skin to restore volume or reduce wrinkles or scars. Common effects include swelling and bruising; infection, lumps, skin damage, or rare blockage of a blood vessel can be serious.

\medterm{Dermal Melanocytosis} A usually harmless blue-gray or slate-colored skin patch caused by pigment-producing cells located deep in the skin. It is often present at birth and commonly affects the lower back or buttocks.

\medterm{Dermal Papilla} A small structure at the base of a hair follicle that supplies signals and nutrients to developing hair cells. Its activity helps control hair growth and the hair-growth cycle.

\medterm{Dermal Structure} The dermis has a thin superficial layer and a deeper supportive layer. It contains collagen, elastic fibers, blood and lymph vessels, nerves, hair follicles, sweat glands, and oil glands.

\medterm{Dermatan Sulfate} A long chain of sugar molecules found in skin, connective tissue, blood vessels, and heart valves. It helps organize tissues and can influence wound healing and blood clotting.
\medterm{Dermatitis} Inflammation of the skin causing some combination of itching, redness, swelling, scaling, dryness, or blisters. It may result from irritation, allergy, infection, medicines, genetics, or another disease.

\medterm{Dermatitis, Atopic} Alternate index label; see \textit{Atopic dermatitis (eczema)}.

\medterm{Dermatitis, Cercarial} Alternate index label; see \textit{Swimmer's itch}.

\medterm{Dermatitis, Contact} Alternate index label; see \textit{Contact dermatitis}.

\medterm{Dermatitis Herpetiformis} A long-lasting, intensely itchy rash with groups of small bumps or blisters, commonly on the elbows, knees, buttocks, or scalp. It is strongly linked to celiac disease and usually improves with a strict gluten-free diet.

\medterm{Dermatitis, Scratch} Alternate index label; see \textit{Neurodermatitis}.

\medterm{Dermatitis, Seborrheic} Alternate index label; see \textit{Seborrheic dermatitis}.

\medterm{Dermatofibroma} A common harmless, firm skin nodule, often on the legs. It may be brown or pink and often dimples when pinched. It usually needs no treatment, but a changing or bleeding lesion should be examined.

\medterm{Dermatofibrosarcoma} A rare, slow-growing cancer of connective tissue in the skin. It usually begins as a firm plaque or nodule and can grow into nearby tissue, but it rarely spreads to distant organs.

\medterm{Dermatofibrosarcoma Protuberans} A rare cancer arising from connective tissue in the skin. It often appears as a slowly enlarging firm plaque or nodule, commonly on the trunk; it can recur locally but rarely spreads. Complete surgical removal is the usual treatment.

\synonyms
DFSP

\medterm{Dermatofibrosarcoma Protuberans Of The Breast} A rare, usually slow-growing skin cancer involving the breast skin or nearby superficial tissue. It can extend into surrounding tissue and recur if incompletely removed; treatment aims for complete surgical removal.

\medterm{Dermatoglyphics} The study of ridge patterns on the fingers, palms, toes, and soles. Patterns help with identification and may differ in some genetic or developmental conditions, but they cannot diagnose a condition by themselves.

\medterm{Dermatographia} A form of hives in which light scratching, rubbing, or pressure causes raised, red, itchy lines within minutes. It is usually harmless, though symptoms may be troublesome.

\medterm{Dermatographia (Dermatographism)} See \textit{Skin Writing}.

\synonyms
Skin Writing

\medterm{Dermatographic Urticaria} A form of hives in which scratching, rubbing, or pressure causes raised, itchy lines on the skin within minutes. The reaction results from release of histamine by skin immune cells.

\medterm{Dermatographism} Hives that appear as raised linear welts within minutes after the skin is stroked or scratched. They result from an exaggerated release of histamine and may occur with other allergic or hive disorders.

\medterm{Dermatologic} Relating to the skin, hair, nails, or medical care for these structures. Dermatologic problems include rashes, infections, wounds, hair loss, nail changes, acne, and skin cancers.

\medterm{Dermatologic Disorder} A disease affecting the skin, hair, nails, or skin glands. It may cause rash, scaling, pigment change, lumps, ulcers, pain, itching, hair loss, or abnormal nail growth.

\medterm{Dermatologist} A physician who specializes in diagnosing and treating disorders of the skin, hair, nails, and related mucous membranes, including infections, inflammatory diseases, tumors, and cosmetic concerns.

\medterm{Dermatology} The medical specialty concerned with preventing, diagnosing, and treating diseases and injuries of the skin, hair, nails, and related mucous membranes.

\medterm{Dermatome} Either a surgical tool that removes a thin layer of skin for a graft or a strip of skin supplied by one spinal nerve. The meaning depends on whether the term is being used in surgery or nerve anatomy.
\medterm{Dermatoscope} A handheld optical instrument that magnifies and illuminates skin lesions, often with polarized or nonpolarized light. It reveals structures below the skin surface and supports clinical assessment of pigmented and other lesions.

\medterm{Dermatomycosis} A fungal infection of the skin, hair, or nails. It may cause scaling, itching, discoloration, cracking, or hair loss; treatment depends on the site and may require a topical or oral antifungal medicine.

\medterm{Dermatomyositis} An autoimmune disease causing muscle inflammation and weakness with characteristic skin changes. It may affect swallowing or breathing muscles, and some adults have an increased risk of associated cancer.

\medterm{Dermatopathology} The study of skin disease by examining tissue samples with a microscope and related laboratory methods. It helps distinguish inflammatory, infectious, harmless, and cancerous skin disorders.

\medterm{Dermatophyte} A fungus that grows in keratin, the tough protein in skin, hair, and nails. It causes tinea infections such as athlete’s foot, ringworm, and fungal nail disease.

\medterm{Dermatophytosis} A superficial fungal infection of the skin, hair, or nails. It can cause ring-shaped rashes, scaling, itching, cracked skin, or thickened discolored nails.

\medterm{Dermatoplasty} Surgical repair or reconstruction of damaged skin, often using tissue rearrangement or a skin graft. It may restore coverage after injury, burns, removal of a lesion, or another operation.
\medterm{Dermatoses - Systemic} Skin changes caused by a disease affecting the whole body or several organs, such as autoimmune, hormone-related, infectious, nutritional, or cancerous disease. The skin finding may help identify the underlying illness.

\medterm{Dermatosis} Any disease or disorder of the skin, including inflammatory, infectious, allergic, inherited, or tumor-related conditions. It may cause rash, itching, scaling, pigment changes, lumps, ulcers, or pain.
\medterm{Dermis} The strong, flexible skin layer beneath the epidermis. It contains collagen and elastic fibers, blood vessels, nerves, hair follicles, and sweat and oil glands that support and nourish the skin.

\medterm{Dermoid} A usually benign cyst or tumor containing tissues such as skin, hair, oil glands, fat, or teeth. It may occur in the ovary or other sites and can cause pressure, rupture, infection, or twisting of an organ.

\medterm{Dermoid Cyst} A usually benign cyst lined by tissue resembling skin and sometimes containing hair, oil glands, fat, or other skin structures. It may be present from birth or develop during life.

\medterm{Dermoid Cyst Of The Ovary} A usually benign ovarian cyst that can contain skin, hair, fat, glands, or teeth. It may cause pain, pressure, rupture, or twisting of the ovary and is assessed with examination and imaging.

\medterm{Dermoscopy} Examination of a skin lesion with magnification and special lighting that reveals structures not visible to the unaided eye. It helps assess pigmented and other lesions but does not replace biopsy when cancer is suspected.

\medterm{Descending Colon} The left-sided part of the large intestine, extending downward from the transverse colon toward the sigmoid colon. It absorbs water and salts and helps compact stool before it moves toward the rectum.

\medterm{Descending Thoracic Aorta} The part of the body’s main artery that runs from the aortic arch through the chest to the diaphragm. It gives branches to the chest wall and organs and continues as the abdominal aorta.

\medterm{Desensitization} A treatment that gradually exposes a person to increasing or controlled amounts of a specific allergen or medicine. It can reduce immune reactions but must be supervised because serious reactions may occur.

\medterm{Desferrioxamine} A medicine that binds excess iron and helps remove it from the body. It is used for serious iron overload and selected aluminum toxicity; prolonged treatment can affect hearing, vision, kidneys, or growth.

\medterm{Desflurane} An inhaled medicine used to maintain general anesthesia. It acts quickly but may irritate the airway, lower blood pressure, cause nausea, or rarely trigger malignant hyperthermia in susceptible people.

\medterm{Desiccate} To dry out or remove moisture from a tissue, substance, surface, or specimen. Excessive drying can damage skin or wounds and can alter medicines or laboratory samples.

\medterm{Desiccated Thyroid} A thyroid-hormone medicine made from dried animal thyroid tissue. It contains mainly thyroxine and triiodothyronine; excess treatment can cause palpitations, bone loss, or other effects, so blood tests and monitoring are needed.

\synonyms
Desiccated thyroid extract; natural desiccated thyroid

\medterm{Desipramine} A tricyclic antidepressant that increases norepinephrine activity. It is used for selected mental health or pain conditions but may cause dry mouth, constipation, drowsiness, abnormal heart rhythms, or dangerous overdose.

\medterm{Desipramine Hydrochloride Overdose} Taking too much desipramine hydrochloride can cause severe drowsiness, confusion, seizures, low blood pressure, coma, and life-threatening heart-rhythm or conduction problems. It requires emergency assessment, even if symptoms are initially mild.

\medterm{Desmin} A structural protein that links parts of muscle cells and helps maintain their alignment. Changes in the desmin gene or abnormal desmin buildup can cause muscle weakness and heart-muscle disease.

\medterm{Desmoid Tumor} A locally aggressive growth of fibrous tissue that usually does not spread to distant organs but can invade nearby structures and recur. It may be associated with familial adenomatous polyposis; treatment depends on location and symptoms.

\medterm{Desmoid Tumors} Fibrous growths that usually do not spread to distant organs but can invade nearby tissues and recur. Some remain stable; observation, medicine, or surgery may be considered according to symptoms and location.

\medterm{Desmoid Tumour} A growth of fibrous connective tissue that usually does not spread elsewhere but can grow into nearby structures and return after treatment. Its management depends on size, site, growth, and symptoms.

\medterm{Desmoplasia} Formation of dense, scar-like tissue around a tumor, long-lasting inflammation, or an injury. Around some cancers it can make a lump feel firm or fixed, but its significance depends on the cells, location, and other findings.

\medterm{Desmoplastic} Describing the formation of dense, scar-like fibrous tissue, often around a tumor or in response to chronic inflammation or injury.

\medterm{Desmoplastic Fibroblastoma} A rare, usually harmless tumor made of collagen-producing cells. It is typically a slow-growing, painless lump in superficial or deep soft tissue and is treated according to its size, site, and symptoms.

\medterm{Desmoplastic Fibroma} A rare noncancerous but locally aggressive fibrous growth in bone. It can cause pain, swelling, or fracture and may return after treatment.

\medterm{Desmoplastic Infantile Ganglioglioma} A rare, usually slow-growing brain tumor of infancy containing both nerve-cell and support-cell components. It may cause seizures, increasing head size, vomiting, or other signs of pressure in the brain.

\medterm{Desmoplastic Melanoma} A rare form of skin melanoma with abundant scar-like tissue. It often develops on sun-damaged skin, may have little or no pigment, and can grow deeply or recur locally.

\medterm{Desmoplastic Reaction} Formation of dense, scar-like fibrous tissue in response to injury, chronic inflammation, or invasion by a cancer. It may make affected tissue firm and can complicate interpretation of imaging or biopsy.
\medterm{Desmoplastic Small Round Cell Tumors} Plural form; see \textit{Desmoplastic small round cell tumor}.

\synonyms
DSRCT

\medterm{Desmoplastic Small Round Cell Tumor} A rare and aggressive soft-tissue sarcoma that usually develops in the abdomen or pelvis of children, adolescents, or young adults, especially males. It often forms several masses on the lining of the abdomen and can spread within the abdomen or to other sites. Symptoms may include abdominal swelling or pain, nausea, vomiting, diarrhea, or constipation. Tumor cells typically have an \textit{EWSR1-WT1} gene fusion, which helps confirm the diagnosis.

\medterm{Desmoplastic Trichoepithelioma} A rare harmless tumor arising from a hair follicle, usually appearing as a small firm facial bump or ring-shaped lesion. It can resemble a skin cancer and may require biopsy for confirmation.

\medterm{Desmopressin} A medicine that mimics vasopressin, a hormone that reduces urine production and helps control bleeding. It is used for central diabetes insipidus, some bleeding disorders, and selected nighttime urination problems; excess fluid can dangerously lower blood sodium.

\medterm{Desmosome} A strong junction that fastens neighboring cells together, especially in the skin and heart muscle. It links internal support fibers and helps tissues resist stretching and friction.
\medterm{Desogestrel} A progestin used in some hormonal contraceptives. It prevents pregnancy mainly by suppressing ovulation and thickening cervical mucus; irregular bleeding, mood changes, breast tenderness, or other hormonal effects may occur.

\medterm{Desonide} A low-strength corticosteroid applied to the skin to reduce redness, itching, and inflammation from conditions such as eczema or dermatitis. Excessive or prolonged use can thin the skin or cause other local effects.

\medterm{Desoximetasone} A corticosteroid applied to the skin to reduce inflammation and itching in responsive rashes. Strength, body site, amount, and duration affect the risk of skin thinning and absorption into the body.

\medterm{Desoxycorticosterone} A steroid hormone with mineralocorticoid effects. It makes the kidneys retain sodium and water and lose potassium, thereby influencing blood volume and blood pressure.

\medterm{Desquamation} Shedding or peeling of the outer layer of skin or another body lining. It may be normal or result from dryness, irritation, infection, inflammation, burns, medicines, or another skin disorder.

\medterm{Desquamative Interstitial Pneumonia} A rare inflammatory lung disease strongly associated with smoking. Immune cells collect in the air sacs, causing cough and breathlessness; stopping smoking is essential, and additional treatment may be needed.

\medterm{Desulfovibrio} A group of bacteria that can live without oxygen and use sulfate during metabolism. They occur in the environment and intestine; invasive human infection is uncommon and must fit the patient’s symptoms and specimen findings.

\medterm{Desvenlafaxine Succinate} A form of desvenlafaxine, an antidepressant that increases serotonin and norepinephrine signaling. It is used for major depressive disorder and may cause nausea, insomnia, increased blood pressure, or withdrawal symptoms if stopped suddenly.

\medterm{Desynchronosis (Jet Lag)} Alternate index label for Jet Lag.

\medterm{Detached} Separated or no longer connected. In medicine it may describe a retina, tissue, device, or feeling of emotional distance; the clinical meaning depends on the surrounding term and findings.

\medterm{Detached Retina} Separation of the light-sensitive retina from the tissue beneath it. Flashes, new floaters, or a curtain-like shadow may occur; it is an emergency because delayed treatment can cause permanent vision loss.

\medterm{Detergent} A cleaning substance that helps remove oils and dirt from surfaces. Some detergents can irritate or burn the skin, eyes, airways, or digestive tract, especially when concentrated or swallowed.

\medterm{Detergent Poisoning} Harmful exposure to a detergent by swallowing, inhaling, or contact with the eyes or skin. Symptoms depend on the product and may include irritation, vomiting, chemical burns, coughing, or breathing difficulty.

\medterm{Deterministic Effects} Tissue injuries from radiation that usually occur only after exposure exceeds a threshold. Higher exposure generally causes more severe injury, such as skin damage, cataracts, or tissue death.

\medterm{Detoxification} Medical treatment that reduces the effects of a poison or helps the body remove it. Care may include supportive treatment, an antidote, activated charcoal, or dialysis; stomach washing is rarely appropriate.

\medterm{Detritic Synovitis} Inflammation of the lining of a joint caused by loose particles from damaged cartilage, bone, blood, crystals, or a worn implant. It can cause swelling, pain, stiffness, and reduced movement.

\medterm{Detrusor Instability} Involuntary contractions of the bladder muscle while the bladder is filling. They can cause sudden urinary urgency, frequent urination, nighttime urination, or urge leakage.

\medterm{Detrusor Muscle} The muscle forming the bladder wall. It relaxes while urine collects and contracts to empty the bladder; abnormal activity can cause urgency or leakage, while weak activity can make emptying difficult.

\medterm{Detumescence} Return of erectile tissue to its non-erect state after sexual stimulation or an erection ends. It occurs when blood leaves the erectile tissue and pressure falls.

\medterm{Deubiquitinating Enzymes} Enzymes that remove a small regulatory tag from other proteins. This can change whether a protein is broken down, where it works, and how cells control signals, repair DNA, and respond to infection.

\medterm{Deubiquitination} Removal of the ubiquitin tag from a protein. This process can prevent the protein from being broken down or can change its activity, location, or role in cell signaling.

\medterm{Deuteranopia} An inherited red-green color-vision deficiency caused by missing green-sensitive pigment in the retina. People may have difficulty distinguishing some shades of green, red, brown, and orange.

\medterm{Deuterium} A naturally occurring form of hydrogen whose nucleus contains one proton and one neutron. Small amounts occur in water; very high exposure to heavy water can disrupt cell function.

\medterm{Development} The progressive growth and change by which a person, organ, or ability acquires its structure and function. It includes physical, brain, language, emotional, social, and adaptive changes.

\medterm{Developmental Coordination Disorder} A childhood condition in which movement skills develop more slowly or less accurately than expected and interfere with dressing, writing, sports, school, or daily activities. It is not explained only by another medical or neurologic condition.

\medterm{Developmental Differences Of The Female Genital Tract} Differences in the formation or structure of the uterus, cervix, vagina, fallopian tubes, or external genitalia that arise before birth. Effects range from none to infertility, pain, menstrual problems, or pregnancy complications.
\medterm{Developmental Disabilities} Lifelong conditions that begin during childhood and affect movement, learning, thinking, language, behavior, or ability to manage daily activities. Support needs vary and may include medical, educational, communication, and social services.

\medterm{Developmental Disorder} A condition that begins during childhood and affects physical, intellectual, language, learning, behavioral, or social development. It may interfere with school, relationships, communication, or daily independence.

\medterm{Developmental Delay} A child’s acquisition of movement, language, thinking, social, or self-care skills is slower than expected for age. Delay may affect one or several areas and can result from genetic, neurologic, sensory, medical, or environmental causes.

\medterm{Developmental Disorders} Conditions beginning during childhood that affect thinking, language, movement, learning, behavior, or social functioning. Assessment considers developmental history, daily impact, coexisting conditions, strengths, and needed educational or clinical support.

\medterm{Developmental Dysplasia Of The Hip} Abnormal development of an infant’s or child’s hip, ranging from a shallow socket to partial or complete dislocation. Early detection and treatment help the hip develop normally and reduce later pain or arthritis.

\medterm{Developmental Expressive Language Disorder} A childhood disorder causing persistent difficulty expressing ideas with spoken words, sentences, or grammar. Hearing, development, autism, and other possible causes should be assessed before diagnosis.
\medterm{Developmental Milestone} A skill that most children acquire within a broad expected age range, such as smiling, sitting, walking, speaking, or interacting. A missed or lost milestone may require developmental assessment.

\medterm{Developmental Milestones} Skills children commonly acquire as they grow, including movement, language, thinking, social interaction, and self-care. Ages vary; delayed or lost skills should be discussed with a healthcare professional.

\medterm{Developmental Milestones Record} A written or electronic record of a child’s movement, language, thinking, social, and self-care skills. It helps follow progress, identify possible delay, and guide further assessment.

\medterm{Developmental Milestones Record - 12 Months} Skills commonly seen around one year include pulling to stand, moving while holding furniture, pointing or using gestures, responding to simple requests, and interactive communication. Children develop at different rates.

\medterm{Developmental Milestones Record - 18 Months} Skills commonly seen around 18 months include walking, using several words, following simple directions, pointing to show interest, and beginning pretend or imitative play. Individual variation is normal.

\medterm{Developmental Milestones Record - 2 Months} Skills commonly seen around two months include social smiling, watching a caregiver or moving object, reacting to loud sounds, making sounds, and holding the head up briefly. Individual development varies.

\medterm{Developmental Milestones Record - 2 Years} Skills commonly seen around two years include running, using two-word phrases, following simple two-step directions, pointing to body parts, and playing near other children. Individual development varies.

\medterm{Developmental Milestones Record - 3 Years} Skills commonly seen around three years include speaking in short sentences, climbing stairs with alternating feet, drawing simple shapes, using a fork, and joining brief interactive play. Individual development varies.

\medterm{Developmental Milestones Record - 4 Months} Skills commonly seen around four months include steady head control, reaching for objects, bringing hands to the mouth, laughing, smiling to gain attention, and responding to caregivers. Individual development varies.

\medterm{Developmental Milestones Record - 4 Years} Skills commonly seen around four years include speaking in understandable sentences, hopping on one foot, drawing a person with several parts, dressing with some help, and engaging in imaginative play.

\medterm{Developmental Milestones Record - 5 Years} Skills commonly seen around five years include hopping or skipping, telling a simple story, counting, drawing recognizable figures or letters, dressing with little help, and following game rules.

\medterm{Developmental Milestones Record - 6 Months} Skills commonly seen around six months include rolling, sitting with support, reaching for and transferring objects, babbling, responding to sounds, and recognizing familiar people. Individual development varies.

\medterm{Developmental Milestones Record - 9 Months} Skills commonly seen around nine months include sitting independently, moving purposefully, babbling, responding to a name, looking for a hidden object, and picking up small objects with finger and thumb.

\medterm{Developmental Process} The physical, brain, emotional, social, and functional changes through which a person or organ grows and acquires abilities. Development is influenced by genes, health, nutrition, relationships, learning, and environment.

\medterm{Developmental Psychology} The study of how thinking, emotions, relationships, behavior, and abilities change from infancy through old age. It examines the effects of genes, health, learning, experience, relationships, and environment.

\medterm{Developmental Reading Disorder} A learning disorder causing persistent difficulty learning accurate, fluent reading. Problems may include recognizing words, linking letters to sounds, spelling, or understanding written language despite appropriate teaching.

\medterm{Developmental Regression} Loss of skills a child previously acquired, such as speech, movement, social interaction, or self-care. Regression is a warning sign requiring prompt evaluation for neurologic, genetic, metabolic, infectious, or other causes.

\medterm{Developmental Stage} A period of growth characterized by typical physical, thinking, emotional, social, language, or self-care changes. Children vary, so stages describe broad patterns rather than exact deadlines.

\medterm{Developmental Surveillance} Ongoing monitoring of a child’s development through caregiver concerns, observation, history, and milestone review during routine care. It can identify concerns early but differs from formal screening with a standardized tool.

\medterm{Deviated Septum} A nasal partition displaced away from the center, often from normal development or injury. A marked deviation can block one side, causing difficulty breathing through the nose, congestion, nosebleeds, sinus problems, or sleep symptoms.

\medterm{Deviated Nasal Septum} Alternate name for deviated septum, in which the wall dividing the nasal passages is displaced from the midline.

\synonyms
Septum Deviated (Deviated Septum)

\medterm{Deviation} A difference from an expected position, direction, value, or pattern. In medicine, it may describe anatomy, movement, a laboratory result, or a measurement and must be interpreted in context.

\medterm{Device Embolisation} Movement of a medical device or broken piece of one through the bloodstream until it lodges in and blocks a blood vessel. It can reduce blood flow and may cause pain, organ injury, or an emergency requiring removal.

\medterm{Devices For Hearing Loss} Equipment that improves access to sound, including hearing aids, cochlear implants, bone-conduction devices, alerting systems, and communication accessories. Choice depends on the type and severity of hearing loss and personal needs.

\medterm{Dexamethasone} A strong, long-acting corticosteroid that reduces inflammation and immune activity. It is used for selected allergic, brain-swelling, cancer-related, endocrine, and other conditions; risks include infection, high blood sugar, mood changes, and adrenal suppression.

\medterm{Dexamethasone Suppression Test} A test of cortisol control. After taking a small dose of dexamethasone, blood or urine cortisol is measured; failure to suppress may suggest excess cortisol but can have other causes and needs clinical interpretation.

\medterm{Dexamphetamine} A stimulant that increases dopamine and norepinephrine activity in the brain. It is used for selected attention or sleep disorders and may cause reduced appetite, insomnia, fast heartbeat, high blood pressure, or dependence.

\medterm{Dexibuprofen} A nonsteroidal anti-inflammatory medicine used for short-term pain and inflammation. It can irritate the stomach and may cause ulcers, bleeding, kidney injury, fluid retention, or cardiovascular problems.

\medterm{Dexketoprofen} A nonsteroidal anti-inflammatory medicine used for short-term pain. It reduces substances that cause pain and inflammation but can cause stomach bleeding, kidney injury, fluid retention, heart problems, or allergic reactions.

\medterm{Dexlansoprazole} A proton-pump inhibitor that reduces stomach-acid production. It is used for acid reflux and related esophageal injury; prolonged use may contribute to infections, low magnesium or vitamin B12, fractures, or kidney inflammation.

\medterm{Dexmedetomidine} A sedative that activates alpha-2 receptors and is given by injection during monitored care or intensive care. It usually depresses breathing less than many sedatives but can slow the heart or lower blood pressure.

\medterm{Dexmethylphenidate} A stimulant used to treat attention-deficit/hyperactivity disorder. It increases dopamine and norepinephrine activity and may cause reduced appetite, insomnia, headache, increased heart rate or blood pressure, or misuse.

\medterm{Dexrazoxane} A medicine used in selected cancer treatments to reduce heart damage from anthracyclines or tissue injury after leakage of an anthracycline from a vein. It can suppress bone marrow and increase infection risk.

\medterm{Dexrazoxane Hydrochloride} A salt form of dexrazoxane used in selected settings to reduce anthracycline-related heart damage or tissue injury after medicine leakage from a vein. Its use requires specialist dosing and monitoring.
\medterm{Dextrin} A group of shorter carbohydrate chains formed when starch breaks down. Dextrins are used in foods, laboratory products, and medicines and are generally digested as carbohydrates.

\medterm{Dextroamphetamine} A stimulant that increases dopamine and norepinephrine activity in the brain. It is used for selected attention or sleep disorders and may cause insomnia, reduced appetite, fast heartbeat, high blood pressure, or dependence.

\medterm{Dextrocardia} A congenital condition in which the heart points or lies mainly on the right side of the chest. It may occur alone or with reversed organs and structural heart defects.

\medterm{Dextromethorphan} A cough suppressant that reduces activity in brain pathways controlling the cough reflex. Excessive amounts can cause drowsiness, agitation, hallucinations, seizures, breathing depression, or serotonin toxicity.

\medterm{Dextromethorphan Overdose} Taking too much dextromethorphan can cause drowsiness, confusion, agitation, hallucinations, fast heartbeat, seizures, breathing depression, or serotonin toxicity. It requires urgent poison or emergency assessment.

\medterm{Dextroposition Of The Heart} The heart is pushed toward the right side of the chest by nearby organs, the diaphragm, lung disease, or another space-occupying problem. The heart itself may be normally formed and is not necessarily reversed as in dextrocardia.

\medterm{Dextrose} D-glucose, a simple sugar used for energy. Oral or intravenous dextrose can treat or prevent low blood glucose, but excessive amounts can raise blood sugar and disturb body fluids.

\medterm{DFSP} Alternate index label; see \textit{Dermatofibrosarcoma protuberans}.

\medterm{DHA} Docosahexaenoic acid, an omega-3 fatty acid that is an important part of brain and retinal cell membranes. It is obtained mainly from oily fish, some algae, and supplements.

\medterm{DHEA} Dehydroepiandrosterone, a hormone precursor made mainly by the adrenal glands. The body can convert it into male and female sex hormones; levels vary with age, illness, and some hormone disorders.

\medterm{DHEA-Sulfate Test} A blood test measuring a hormone made mainly by the adrenal glands. It can help investigate early puberty, excess body hair, some hormone disorders, congenital adrenal hyperplasia, or an adrenal tumor; results vary with age and sex.

\medterm{Diabetes And Alcohol} Alcohol can cause delayed low blood glucose, interact with medicines, and make glucose control harder. Risk depends on amount, timing, food intake, diabetes type, liver health, and treatment.

\medterm{Diabetes And Exercise} Physical activity can improve glucose control, fitness, and heart health. People using insulin or medicines that lower glucose may need to check levels and adjust food or treatment with professional guidance to prevent low glucose.

\medterm{Diabetes And Eye Disease} Diabetes can damage retinal blood vessels, causing retinopathy and macular swelling that may reduce vision. Good glucose and blood-pressure control and regular dilated eye examinations help prevent or limit sight loss.

\medterm{Diabetes And Eye Problems} Diabetes may cause retinopathy, macular swelling, cataracts, glaucoma, and vision loss. Glucose and blood-pressure control, regular dilated eye examinations, and timely laser, injection, or surgical treatment reduce the risk of severe loss.

\medterm{Diabetes And Kidney Disease} Kidney damage caused by diabetes. It may begin with albumin leaking into urine and progress to reduced kidney function or kidney failure; glucose, blood-pressure, and kidney monitoring can slow progression.

\medterm{Diabetes And Metabolic Health} Diabetes is linked with insulin resistance, excess body fat, abnormal cholesterol, and high blood pressure. Together these problems increase the risk of heart disease and can worsen blood-glucose control.

\medterm{Diabetes And Nerve Damage} Long-term diabetes can injure nerves, causing burning pain, tingling, numbness, weakness, digestive or bladder problems, sexual difficulties, or loss of protective feeling in the feet.

\medterm{Diabetes Complications} Problems caused by diabetes, including low or high glucose emergencies and long-term damage to the eyes, kidneys, nerves, heart, blood vessels, teeth, and feet. Regular monitoring and risk-factor control reduce risk.

\medterm{Diabetes During Pregnancy (Gestational Diabetes)} Alternate index label for Gestational Diabetes.


\medterm{Diabetes Eye Exams} Eye examinations that look for diabetic retinopathy, macular swelling, cataracts, glaucoma, and other sight-threatening problems. Dilated examination or retinal photography may detect changes before vision is affected.

\medterm{Diabetes - Foot Ulcers} Open sores on the foot in a person with diabetes, often caused or worsened by reduced sensation, pressure, deformity, or poor circulation. They can become infected or reach bone, so prompt assessment and wound care are important.

\medterm{Diabetes, Gestational} Alternate index label; see \textit{Gestational diabetes}.

\medterm{Diabetes Insipidus} A disorder causing large amounts of very dilute urine and intense thirst because vasopressin is lacking or the kidneys do not respond to it. It is unrelated to diabetes mellitus and can cause dehydration.

\medterm{Diabetes - Insulin Therapy} Treatment with injected or infused insulin to replace missing insulin or improve glucose control. Doses are individualized and balanced with meals, activity, illness, and the risk of dangerously low glucose.


\medterm{Diabetes Management} Coordinated care using food choices, activity, medicines when needed, glucose monitoring, education, and screening for complications. The plan is adjusted for diabetes type, health conditions, life stage, and personal goals.

\medterm{Diabetes Medications} Medicines that lower blood glucose or replace insulin. They work in different ways, and the choice depends on the type of diabetes, kidney and heart health, other medicines, cost, and treatment goals.

\medterm{Diabetes Mellitus} A group of diseases in which blood glucose, also called blood sugar, stays too high. Insulin is a hormone made by the pancreas that helps glucose move from the blood into cells for energy. Diabetes develops when the body makes too little or no insulin, cannot use insulin properly, or both. The main types are type 1 diabetes, type 2 diabetes, and gestational diabetes, which develops during pregnancy; other forms can result from genetic conditions, pancreatic disease, or medicines. Persistently high blood glucose can damage the eyes, kidneys, nerves, heart, and blood vessels. Diabetes mellitus is different from diabetes insipidus, a disorder of water balance that is not caused by high blood glucose.

\synonyms
Diabetes

\medterm{Diabetes Mellitus Type 1} An autoimmune disorder in which the immune system destroys insulin-producing pancreatic cells, causing little or no insulin. It can begin at any age and requires insulin; untreated deficiency can cause ketoacidosis.

\medterm{Diabetes Mellitus Type 2} A disorder in which the body resists insulin and, over time, may not make enough of it. Risk is increased by family history, excess body weight, and inactivity; it can occur at any age and may have no early symptoms.


\medterm{Diabetes Pre (Prediabetes)} Alternate index label for Prediabetes.

\medterm{Diabetes - Preventing Heart Attack And Stroke} Reducing cardiovascular risk in diabetes through glucose, blood-pressure, and cholesterol control, not smoking, regular activity, healthy eating, weight management when appropriate, and indicated preventive medicines.




\medterm{Diabetes Prevention} The risk of type 2 diabetes can often be reduced through regular physical activity, nutritious meals, avoiding tobacco, adequate sleep, and weight loss when appropriate. People at increased risk may benefit from structured prevention programs and regular testing.

\medterm{Diabetes Symptoms In Men} Diabetes in men may cause thirst, frequent urination, fatigue, blurred vision, slow wound healing, recurrent infections, erectile dysfunction, or loss of sensation. Some people have no symptoms, so risk-based blood testing is important.

\medterm{Diabetes Symptoms In Women} Diabetes in women may cause thirst, frequent urination, fatigue, blurred vision, recurrent vaginal or urinary infections, slow wound healing, or pregnancy complications. Symptoms can be absent, and screening depends on age, risk factors, and pregnancy status.


\medterm{Diabetes Treatment} Diabetes treatment combines nutrition, physical activity, glucose monitoring, education, and medicines chosen for the type of diabetes and individual risks. Care also addresses blood pressure, lipids, kidneys, eyes, nerves, feet, vaccines, and cardiovascular health.

\medterm{Diabetes, Type 1} Alternate index label; see \textit{Type 1 diabetes}.


\medterm{Diabetes Type 1 (Type 1 Diabetes)} Alternate index label for Type 1 Diabetes.

\medterm{Diabetes, Type 2} Alternate index label; see \textit{Type 2 diabetes}.


\medterm{Diabetes Type 2 - Meal Planning} Meal planning for type 2 diabetes emphasizes suitable portions, fiber-rich foods, carbohydrate awareness, adequate protein, and unsaturated fats. The plan should fit medicines, kidney health, cultural foods, preferences, and glucose goals.

\medterm{Diabetes Type 2 (Type 2 Diabetes)} Alternate index label for Type 2 Diabetes.

\medterm{Diabetic Angiopathy} Damage to small or large blood vessels caused by diabetes. It can reduce circulation, injure the eyes, kidneys, heart, brain, or feet, and increase the risk of wounds, heart attack, or stroke.

\medterm{Diabetic Autonomic Neuropathy} Diabetes-related damage to nerves that control automatic body functions. It may cause dizziness on standing, fast heartbeat, sweating changes, delayed stomach emptying, bladder problems, sexual dysfunction, or reduced awareness of low glucose.

\medterm{Diabetic Cardiomyopathies} Heart-muscle disease associated with diabetes that is not fully explained by blocked heart arteries or high blood pressure. It can make the heart relax or pump poorly and may lead to heart failure.

\medterm{Diabetic Cardiomyopathy} Heart-muscle damage associated with diabetes, after other causes such as blocked arteries and high blood pressure are considered. It may cause poor relaxation, weak pumping, heart failure, or rhythm problems.

\medterm{Diabetic Coma} Profound unconsciousness caused by a severe diabetes-related problem, such as extremely low or high blood glucose, ketoacidosis, severe dehydration, or abnormal body chemistry. It is a medical emergency.

\medterm{Diabetic Complications} Acute or long-term problems caused by diabetes, including low or high glucose emergencies and damage to the eyes, kidneys, nerves, heart, blood vessels, teeth, and feet.

\medterm{Diabetic Diet} An individualized eating pattern for diabetes that balances carbohydrate portions, fiber, protein, and healthy fats with medicines, activity, culture, and health goals. It supports glucose control without requiring the same foods or schedule for everyone.

\synonyms
Diet Diabetes (Diabetic Diet)

\medterm{Diabetic Diet For Type 2 Diabetes} An eating pattern for type 2 diabetes emphasizing suitable portions, high-fiber foods, carbohydrate awareness, adequate protein, and unsaturated fats. Medicines, kidney function, culture, preferences, and glucose or weight goals guide the plan.

\medterm{Diabetic Foot} Foot problems associated with diabetes, including loss of sensation, deformity, ulcers, infection, poor circulation, and slow healing. Daily inspection, proper footwear, glucose control, and early care help prevent serious injury.

\medterm{Diabetic Foot Disease} Foot damage caused by a combination of diabetes-related nerve loss, poor circulation, deformity, and impaired healing or immunity. A small injury can progress to an ulcer, deep infection, or amputation without early care.

\medterm{Diabetic Foot Ulcer} An open wound on the foot caused or worsened by diabetes-related loss of sensation, pressure, deformity, poor circulation, or slow healing. It may be painless and can become deeply infected without prompt treatment.

\medterm{Diabetic Gastroparesis} A complication of diabetes in which damage to the stomach's nerves slows the movement of food from the stomach into the small intestine. It may cause early fullness, nausea, vomiting, bloating, abdominal pain, and difficult-to-control blood sugar levels.


\medterm{Diabetic Hyperglycemic Hyperosmolar Syndrome} A life-threatening diabetes emergency with extremely high blood glucose, severe dehydration, and concentrated blood, usually without major ketone buildup. It can cause confusion, seizures, coma, and shock and requires hospital treatment.

\medterm{Diabetic Hyperosmolar Hyperglycemic State} A severe diabetes emergency marked by extremely high glucose, severe dehydration, and increased blood concentration with little or no ketoacidosis. Infection, illness, or inadequate treatment may trigger it; urgent hospital care is required.

\medterm{Diabetic Hypoglycemia} Blood glucose that is too low in a person with diabetes, often because of insulin, certain medicines, missed food, alcohol, or unusual exercise. Shaking, sweating, confusion, seizures, or unconsciousness may occur.

\medterm{Diabetic Ketoacidosis} A life-threatening emergency caused by too little effective insulin, leading to high glucose, ketone buildup, and acidic blood. It can cause thirst, frequent urination, vomiting, abdominal pain, deep breathing, confusion, and coma.

\synonyms
DKA

\medterm{Diabetic Ketoacidosis Monitoring} Monitoring during treatment includes frequent glucose checks, blood salts, kidney function, ketones, acidity, fluid balance, urine output, and mental status. Close monitoring helps guide fluids, insulin, and electrolyte replacement safely.

\medterm{Diabetic Macular Edema} Swelling of the macula, the central part of the retina, caused by leaking blood vessels in diabetes. It can blur or distort central vision and may occur at any stage of diabetic retinopathy.

\medterm{Diabetic Mastopathy} A rare, noncancerous, firm breast thickening associated mainly with long-standing diabetes. It can feel or look like breast cancer, so examination and imaging—and sometimes a biopsy—may be needed to establish the diagnosis.

\medterm{Diabetic Microvascular Complications} Diabetes-related damage to small blood vessels, mainly affecting the eyes, kidneys, and nerves. Good glucose, blood-pressure, and cholesterol control, along with regular screening, lowers the risk of progression.
\medterm{Diabetic Nephropathy} Kidney damage caused by diabetes. It often begins with albumin leaking into urine and may progress to reduced filtering ability or kidney failure; regular urine and blood tests help detect and monitor it.

\synonyms
DKD

\medterm{Diabetic Nephropathy (Kidney Disease)} Kidney damage caused by diabetes, often beginning with albumin in the urine and potentially progressing to chronic kidney disease or kidney failure. Blood and urine monitoring can detect it early.

\medterm{Diabetic Neuropathy} Nerve damage caused by diabetes. It may cause burning pain, tingling, numbness, weakness, loss of protective feeling in the feet, or problems with digestion, blood pressure, sweating, bladder control, or sexual function.

\synonyms
Autonomic Neuropathy Diabetic (Diabetic Neuropathy)
Neuropathy, Diabetic
Neuropathy Autonomic Diabetic (Diabetic Neuropathy)
Neuropathy Diabetic (Diabetic Neuropathy)
Neuropathy Diabetic Peripheral (Diabetic Neuropathy)
Neuropathy Focal Diabetic (Diabetic Neuropathy)
Neuropathy Proximal Diabetic (Diabetic Neuropathy)
Peripheral Neuropathy Diabetic (Diabetic Neuropathy)
Proximal Neuropathy Diabetic (Diabetic Neuropathy)

\medterm{Diabetic Peripheral Neuropathy} Diabetes-related damage to nerves, usually beginning in the toes and feet and later affecting the hands. It can cause numbness, burning pain, weakness, and loss of protective sensation that increases injury risk.

\medterm{Diabetic Retinopathy} Damage to retinal blood vessels caused by diabetes. Early disease may have no symptoms; advanced disease can cause macular swelling, bleeding, retinal detachment, and vision loss. Regular dilated eye examinations are important.

\synonyms
Retinopathy, Diabetic

\medterm{Diabetogenic} Able to cause or promote diabetes or high blood glucose. The term may describe a disease, hormone, medicine, genetic change, or other factor that reduces insulin action or increases glucose production.

\medterm{Diabetologist} A physician specializing in diabetes prevention, diagnosis, and treatment, including glucose management and complications involving the heart, kidneys, nerves, eyes, and feet.

\medterm{Diabrosis} An old medical term for erosion or corrosion that eats through tissue and may cause an ulcer or perforation, especially in a blood-vessel or organ wall. Modern reports usually name the specific cause and injury.

\medterm{Diacetoxyscirpenol} A toxin produced by some Fusarium molds that can contaminate food or animal feed. Exposure may irritate the digestive tract, suppress blood-cell production, and impair immune function.


\medterm{Diacylglycerol} A fat-like molecule made from cell membranes that acts as an internal chemical signal. It helps control cell growth, movement, secretion, metabolism, and responses to hormones or other signals.

\medterm{Diagnosis} Identification of a disease or health condition using symptoms, medical history, examination, and appropriate tests. A clinician may compare several possible causes before reaching the most likely diagnosis.

\medterm{Diagnostic} Relating to identifying a disease, condition, or cause of symptoms. A diagnostic test provides information that must be interpreted with the person’s history, examination, and other findings rather than alone.

\medterm{Diagnostic Imaging} Use of images to detect, locate, or monitor disease or injury. Methods include X-ray, ultrasound, computed tomography, magnetic resonance imaging, and nuclear imaging; findings must be interpreted with the clinical situation.

\medterm{Diagnostic Laparoscopy} A minimally invasive procedure in which a camera is inserted through small abdominal incisions to inspect pelvic or abdominal organs. The clinician may take tissue samples or perform treatment during the same procedure.

\medterm{Diagnostic Microbiology} Laboratory testing used to identify microorganisms that may cause disease. Methods include microscopy, culture, antigen or antibody tests, and genetic tests; results must be interpreted with symptoms and the specimen site.

\medterm{Diagnostic Procedure} An examination or procedure used to identify, exclude, locate, or monitor a disease or injury. The purpose, preparation, risks, accuracy, and follow-up depend on the specific procedure and patient.

\medterm{Diagnostic Radiology} The medical specialty that uses imaging to diagnose and characterize disease or injury. It includes X-ray, ultrasound, computed tomography, magnetic resonance imaging, and selected nuclear-medicine studies.

\medterm{Diagnostic Service} A healthcare service that collects or interprets information to detect, identify, characterize, or monitor a disease or health condition, such as laboratory testing, imaging, or specialist assessment.

\medterm{Diagnostic Test} A test used to evaluate symptoms, signs, or an abnormal screening result and help confirm, exclude, or characterize a disease. Its value depends on accuracy, timing, the patient’s condition, and how the result will guide care.

\medterm{Diagnostic Trial} A clinical study that evaluates how accurately a test detects, rules out, or predicts a disease or important health outcome. Researchers may compare it with an established test or another accepted way of determining the diagnosis.

\medterm{Diagnostic Uncertainty} A situation in which available information does not establish one diagnosis with enough confidence. It is managed by explaining uncertainty, arranging follow-up, watching for warning signs, and ordering targeted further tests when appropriate.

\medterm{Diagnostics} The tests, examinations, and clinical reasoning used to detect, identify, classify, or monitor a health condition. Results must be interpreted in context because no test is perfect.

\medterm{Diakinesis} The final stage before the first division that produces eggs or sperm. Paired chromosomes become tightly packed and remain connected at points where they exchanged genetic material.

\medterm{Dialectical Behavior Therapy} A structured talking treatment that combines behavior-change skills with mindfulness and acceptance. It teaches ways to manage intense emotions, tolerate distress, improve relationships, and reduce harmful behaviors.

\medterm{Dialister Pneumosintes} A bacterium found in the mouth and respiratory tract that grows without oxygen. It may contribute to mixed infections, but a positive sample must be interpreted with symptoms and collection site.

\medterm{Dialysis} Treatment that replaces part of the kidneys’ work by removing waste, excess water, and some chemicals from the blood. The main types are hemodialysis and peritoneal dialysis.

\medterm{Dialysis Access} A route used to connect a person’s blood circulation to a hemodialysis machine. Options include an arteriovenous fistula, graft, or central venous catheter; each has different risks of infection, clotting, and inadequate flow.

\medterm{Dialysis Adequacy} How well dialysis removes waste and extra fluid and keeps body salts and acid levels within a safe range. Blood tests, treatment records, symptoms, nutrition, and weight help determine whether dialysis is sufficient.


\medterm{Dialysis Disequilibrium Syndrome} A rare brain reaction during or soon after starting or intensifying hemodialysis, caused by rapid changes in dissolved substances. Headache, nausea, confusion, seizures, or coma may occur; gradual treatment reduces risk.

\medterm{Dialysis - Hemodialysis} A form of dialysis in which blood travels through a filter outside the body to remove waste and excess fluid before returning to the patient. It requires reliable blood access and regular monitoring.

\medterm{Dialysis Machine} Equipment that moves blood or dialysis fluid across a filter to remove waste, excess water, and selected salts when the kidneys cannot do so adequately.

\medterm{Dialysis - Peritoneal} A form of dialysis in which cleansing fluid enters and leaves the abdomen through a catheter. The abdominal lining filters waste and water; cloudy fluid, pain, or fever may signal peritonitis and requires urgent care.

\medterm{Peritoneal Dialysis-Associated Peritonitis} Infection or inflammation of the abdominal lining during peritoneal dialysis, usually related to contamination of the dialysis catheter or fluid. It may cause cloudy drainage, abdominal pain, fever, or nausea.

\medterm{Dialysis Solutions} Specially prepared fluids used during hemodialysis or peritoneal dialysis. They allow waste and excess water to move out of the blood; their salt, glucose, and other contents are adjusted to the patient’s needs.

\medterm{Dialysis Therapy} Use of hemodialysis or peritoneal dialysis to replace part of kidney function by removing waste and excess fluid and correcting some salt and acid imbalances. Treatment is individualized according to kidney function, symptoms, health, and goals.

\medterm{Diamine} A molecule containing two amino groups. Diamines occur naturally in metabolism and are used in laboratory, pharmaceutical, and industrial chemistry; health effects depend on the specific compound and exposure.

\medterm{Diamond-Blackfan Anemia} A rare inherited disorder in which the bone marrow makes too few red blood cells, usually beginning in infancy. It may also cause short stature, thumb or facial differences, or heart defects.

\synonyms
DBA; Diamond-Blackfan syndrome; Congenital hypoplastic anemia


\medterm{Diamorphine} A strong opioid medicine used in selected settings for severe pain or during medical procedures. It can cause sleepiness, nausea, constipation, dependence, and dangerous breathing depression, especially in overdose or with other sedatives.

\medterm{Dianisidine} An industrial chemical used in some dyes and laboratory reagents. Significant or repeated exposure can irritate tissues and may damage blood or increase cancer risk; workplace protection is important.


\medterm{Diapedesis} Movement of blood cells, especially infection-fighting white cells, through the walls of tiny blood vessels into nearby tissue. This allows them to reach and respond to infection, injury, or inflammation.

\medterm{Diaper Rash} Inflammation of skin covered by a diaper, commonly caused by moisture, urine, stool, friction, or yeast infection. It causes redness and soreness; a persistent, severe, blistering, or spreading rash needs medical assessment.

\medterm{Diaphoresis} Excessive sweating. It may be a normal response to heat or exercise, or a sign of fever, pain, low blood glucose, anxiety, infection, heart problems, or another condition.

\medterm{Diaphoretic} Describing heavy sweating or something that causes sweating. Sudden cold sweating with chest pain, faintness, breathlessness, confusion, or weakness may signal a serious medical problem and needs urgent assessment.

\medterm{Diaphragm} The dome-shaped muscle separating the chest from the abdomen and the main muscle used for breathing. When it contracts, it flattens and draws air into the lungs; when it relaxes, air leaves.

\medterm{Diaphragm (Muscle)} The dome-shaped breathing muscle between the chest and abdomen. Contraction enlarges the chest cavity and draws air into the lungs; relaxation allows exhalation.

\medterm{Diaphragmatic Eventration} Abnormal thinning or weakness of part or all of the diaphragm, causing it to sit higher than usual. It may cause no symptoms or lead to breathlessness, especially if severe or present from birth.

\medterm{Diaphragmatic Hernia} Movement of abdominal organs through a hole or weak area in the diaphragm into the chest. It may be present from birth or follow injury and can impair breathing or cause bowel obstruction.

\medterm{Diaphragmatic Rupture} A tear in the diaphragm, usually caused by a severe blunt or penetrating injury to the chest or abdomen. Abdominal organs may move into the chest, causing breathing difficulty, obstruction, or strangulation.

\medterm{Diaphysis} The long central shaft of a bone such as the femur, humerus, or tibia. It contains a hollow marrow cavity and is surrounded by strong compact bone; fractures there may require prolonged healing.

\medterm{Diarrhea} Passing loose or watery stools more often than usual. Causes include infections, medicines, food intolerance, inflammation, malabsorption, or other illnesses. Diarrhea can lead to dehydration and, when prolonged, malabsorption.



\medterm{Diarrhea, Constipation, Gas, And Bloating} These digestive symptoms may result from diet, infection, medicines, food intolerance, constipation, or disorders such as irritable bowel syndrome. Persistent, severe, bloody, or unexplained symptoms require assessment.

\medterm{Diarrhea IBS (IBS-D Irritable Bowel Syndrome With Diarrhea)} Alternate index label for IBS-D Irritable Bowel Syndrome with Diarrhea.

\medterm{Diarrhea In Infants} Frequent loose or watery stools in a baby. Because infants can dehydrate quickly, poor feeding, fewer wet diapers, dry mouth, unusual sleepiness, fever, blood, or repeated vomiting requires prompt assessment.

\medterm{Diarrhea, Traveler's} Alternate index label; see \textit{Traveler's diarrhea}.

\medterm{Diarrhea Travelers (Traveler's Diarrhea)} Alternate index label for Traveler's Diarrhea.

\medterm{Diarrhoea} Alternate index label; see \textit{Diarrhea}.

\medterm{Diarthrosis} A freely movable joint in which the bones are separated by a space containing lubricating fluid. The knee, hip, shoulder, and elbow are examples.

\medterm{Diastase} An older name for enzymes that break starch into smaller sugars. Modern medical writing usually uses the term amylase for the related digestive enzyme activity.

\medterm{Diastasis} Separation or widening between structures that normally lie together. It commonly describes separation of the abdominal muscles after pregnancy and is not itself a true hernia, although both can occur together.

\medterm{Diastasis Recti} Widening and thinning of the tissue between the two front abdominal muscles, commonly during or after pregnancy or major weight change. It can cause a midline bulge when the abdomen is strained and is not a true hernia.

\medterm{Diastematomyelia} A birth defect in which the spinal cord is divided into two parts, often by bone or fibrous tissue. It may occur with abnormal vertebrae, a skin mark, weakness, altered sensation, or bladder problems.

\medterm{Diastole} The part of each heartbeat when the heart muscle relaxes and its chambers fill with blood. The lower blood-pressure number is the pressure in the arteries during this phase.

\medterm{Diastolic} Relating to diastole, when the heart relaxes and fills. Diastolic blood pressure is the lower number in a blood-pressure reading.

\medterm{Diastolic Blood Pressure} The pressure in the arteries while the heart relaxes between beats. It is the lower number in a blood-pressure reading and is interpreted with the upper number, symptoms, age, and other health conditions.

\medterm{Diastolic Dysfunction} Impaired relaxation or increased stiffness of a heart chamber, reducing its ability to fill normally. It can raise pressure in the lungs and cause breathlessness or heart failure, even when pumping strength appears preserved.

\medterm{Diastolic Heart Failure} An older term usually referring to heart failure caused by a stiff heart that does not fill easily, often with preserved pumping percentage. Diagnosis also requires symptoms, signs, and evidence of raised filling pressure.

\medterm{Diastolic Murmur} An extra or unusual heart sound heard while the heart relaxes between beats. It often reflects abnormal flow through a heart valve or nearby blood vessel and needs medical evaluation.

\medterm{Diastolic Murmurs} Extra heart sounds heard during the relaxation phase between beats. They commonly result from abnormal flow through a heart valve or major blood vessel and require clinical assessment.

\medterm{Diathermy} Use of high-frequency electrical or electromagnetic energy to produce heat in tissue. It may be used during surgery to cut tissue or stop bleeding, or in selected treatments to warm deeper tissues.

\medterm{Diathesis} A tendency to develop a particular disease, symptom, or body response. For example, a bleeding diathesis means a person is more likely to bleed than usual.

\medterm{Diatom Test} A forensic test that looks for microscopic algae in tissues after suspected drowning. It may support an investigation but cannot prove drowning by itself because algae can enter the body from other sources.

\medterm{Diatomaceous Earth} A powder made from the fossilized remains of microscopic algae. Inhalation of fine particles can irritate the lungs, and some forms containing crystalline silica can cause serious lung disease with repeated exposure.

\medterm{Diatrizoate Meglumine} An iodinated contrast substance used for selected X-ray or imaging studies. It improves visibility of body structures but may cause allergic-like reactions, kidney injury, aspiration problems, or thyroid effects.

\medterm{Diazepam} A benzodiazepine that calms brain activity and is used for selected anxiety, muscle-spasm, seizure, and alcohol-withdrawal conditions. It can cause drowsiness, falls, dependence, withdrawal, and dangerous breathing depression with opioids or alcohol.

\medterm{Diazepam Overdose} Taking too much diazepam can cause extreme sleepiness, confusion, poor coordination, low blood pressure, slowed breathing, coma, or death, especially when combined with alcohol, opioids, or other sedatives.

\medterm{Diazinon} An organophosphate insecticide that blocks acetylcholinesterase, causing excess nerve signals. Exposure can lead to sweating, salivation, vomiting, diarrhea, breathing difficulty, muscle weakness, seizures, or respiratory failure.

\medterm{Diazinon Poisoning} Poisoning by diazinon, causing excess cholinergic activity with sweating, salivation, vomiting, diarrhea, wheezing, slow heartbeat, muscle weakness, seizures, or breathing failure. It requires emergency treatment and decontamination.

\medterm{Diazonium Compounds} Chemicals containing a diazonium group, used in dyes, synthesis, and laboratory reactions. Some are unstable or toxic and may irritate tissues or release hazardous products.

\medterm{Diazoxide} A medicine that opens potassium channels and can reduce insulin release or relax blood vessels. It is used in selected cases of low blood sugar from excess insulin and severe high blood pressure; fluid retention and high glucose may occur.

\medterm{Dibekacin} An aminoglycoside antibiotic used in some regions for serious bacterial infections. It blocks bacterial protein production but can damage the kidneys, hearing, or balance, especially with prolonged or high-dose treatment.

\medterm{Dibenzepin} A tricyclic antidepressant used in some countries for depressive disorders. It can cause drowsiness, dry mouth, constipation, blurred vision, abnormal heart rhythms, and dangerous toxicity in overdose.

\medterm{Dibucaine} A long-acting local anesthetic applied to skin or mucous membranes to reduce pain or itching. Excessive use or absorption can cause numbness beyond the treatment area, seizures, heart problems, or other toxicity.

\medterm{Dicarboxylic Acid} An organic molecule containing two acid groups. These molecules take part in metabolism and may accumulate in some inherited or acquired metabolic disorders.

\medterm{Dichlorophen} A chemical formerly used in some disinfectants and medicines. Exposure can irritate the skin, eyes, or digestive tract; serious effects depend on the amount, route, and duration of exposure.

\medterm{Dichlorvos} An organophosphate insecticide that can block acetylcholinesterase. Poisoning may cause sweating, salivation, vomiting, diarrhea, muscle weakness, breathing difficulty, seizures, or respiratory failure.

\medterm{Dichorionic Diamniotic Twin Placenta} A twin pregnancy in which each fetus has its own outer membrane and amniotic sac, usually with a separate placenta. It generally has fewer shared-placenta complications, but routine maternal and fetal monitoring remains important.

\medterm{Dichorionic Twins} Twins that each have a separate outer membrane and amniotic sac. Most fraternal twins and some identical twins are dichorionic; risks still include premature birth, growth problems, and pregnancy complications.

\medterm{Diclofenac} A nonsteroidal anti-inflammatory medicine used for pain and inflammation. It can cause stomach irritation or bleeding, kidney or liver injury, fluid retention, and increased cardiovascular risk.

\medterm{Diclofenac Sodium} A salt form of diclofenac used for pain and inflammation. It can irritate or bleed the stomach and may affect the kidneys, blood pressure, liver, or heart.

\medterm{Diclofenac Sodium Overdose} Taking too much diclofenac sodium can cause nausea, vomiting, stomach bleeding, kidney injury, drowsiness, seizures, breathing problems, or abnormal body chemistry and requires urgent medical assessment.

\medterm{Dicloxacillin} A penicillin-type antibiotic that blocks bacterial cell-wall production. It is used for susceptible infections, especially some staphylococcal infections; allergy, diarrhea, and liver injury can occur.

\medterm{DICOM} An international format and communication standard for storing, sending, and viewing medical images and related information. It allows images from X-ray, CT, MRI, ultrasound, and other systems to be shared.

\medterm{Dicoumarol} An older oral anticoagulant that blocks vitamin K recycling and reduces production of several clotting factors. It has largely been replaced by medicines with more predictable effects and remains important mainly historically.

\medterm{Dicrotic} Describing a second small wave or notch in an artery's pressure pulse, usually produced when the aortic valve closes. It may be visible on a monitor or arterial pressure tracing.

\medterm{Dicrotophos} An organophosphate insecticide that can inhibit acetylcholinesterase. Poisoning may cause sweating, salivation, vomiting, diarrhea, muscle weakness, breathing difficulty, seizures, or respiratory failure.


\medterm{Dicyclomine} An anticholinergic medicine that relaxes intestinal muscle and may reduce cramping in selected bowel disorders. It can cause dry mouth, blurred vision, constipation, fast heartbeat, confusion, or difficulty urinating.

\medterm{Dieffenbachia Poisoning} Poisoning after chewing or swallowing parts of the Dieffenbachia plant. Needle-like calcium oxalate crystals can cause immediate burning, pain, drooling, and swelling of the mouth and throat; airway swelling requires emergency care.

\medterm{Diencephalon} A central part of the brain containing the thalamus and hypothalamus. It helps process sensory information and regulate hormones, temperature, appetite, sleep, movement, and other automatic functions.

\medterm{DIEP Flap} A reconstructive surgery using skin and fat from the lower abdomen, supplied by small blood vessels, to rebuild a breast or another body area while preserving most abdominal muscle.

\medterm{Diesel Oil} A petroleum fuel whose fumes or liquid can irritate the skin, eyes, and lungs. Swallowing or aspiration into the lungs can cause severe chemical pneumonia; suspected exposure requires prompt poison or medical advice.

\medterm{Diestrus} A phase of the reproductive cycle in many mammals between periods of sexual receptivity. It is characterized by activity and later regression of the corpus luteum.





\medterm{Diet} The foods and drinks a person regularly consumes. A health-supporting diet provides enough energy, protein, vitamins, minerals, fiber, and fluids and is adapted to health needs, culture, preferences, and available foods.

\medterm{Diet Diabetes (Diabetic Diet)} Alternate index label for Diabetic Diet.




\medterm{Diet Therapy} Use of an individualized eating plan to prevent, manage, or treat a health problem. It may address diabetes, kidney disease, food allergy, malnutrition, digestive disorders, or other needs and is tailored to the person's health and goals.

\medterm{Diet Ulcerative Colitis (Ulcerative Colitis Diet)} Alternate index label for Ulcerative Colitis Diet.



\medterm{Dietary Assessment} Systematic review of a person’s food, drink, nutrient, and supplement intake using methods such as a food diary, recall, or questionnaire. It helps identify deficiencies, excesses, patterns, and changes needed for health.

\medterm{Dietary Calcium} Calcium obtained from food or drinks, needed for bones, teeth, muscles, nerves, and blood clotting. Requirements vary with age and health; excess supplements may cause kidney stones or other problems.

\medterm{Dietary Carbohydrate} Carbohydrate in foods and drinks, including sugars, starches, and fiber. Digestible carbohydrate provides energy and affects blood glucose; fiber is only partly or not digested and supports bowel function.

\medterm{Dietary Fat} A nutrient made mainly of triglycerides and fatty acids that provides energy, supports cell membranes, and helps absorb vitamins A, D, E, and K. The type and amount of fat affect cardiovascular health.


\medterm{Dietary Fats} Fats obtained from food, including saturated, monounsaturated, polyunsaturated, and trans fats. They provide energy and support cell membranes and vitamin absorption; replacing saturated or trans fats with unsaturated fats is generally healthier.


\medterm{Dietary Fiber} Plant material that is not fully digested. It helps form stool, supports regular bowel movements, may improve cholesterol and glucose control, and can increase fullness.

\medterm{Dietary Fibre} Plant-derived carbohydrate that is not fully digested in the small intestine. Soluble and insoluble fiber support bowel function and may improve cholesterol, glucose control, and fullness; increase intake gradually with adequate fluids unless medically restricted.


\medterm{Dietary Iron} Iron from food or supplements, needed to make hemoglobin, muscle proteins, and enzymes. Too little can cause iron-deficiency anemia; too much can be toxic, especially in children or people with iron-storage disorders.

\medterm{Dietary Phospholipid} A food-derived fat containing phosphate that helps form cell membranes and transports fats in the blood. Its digestion and handling depend on food intake, absorption, and liver function.

\medterm{Dietary Phosphorus} Phosphorus from food, needed for bones, teeth, cell membranes, and energy use. People with reduced kidney function may need to limit highly absorbable sources to prevent high blood phosphorus.

\medterm{Dietary Protein} Protein from food that is digested into amino acids and small peptides. These support muscle and tissue repair, enzymes, hormones, immunity, and energy when needed.

\medterm{Dietary Reference Intakes} A set of nutrient-intake reference values used to plan and assess diets. They include the recommended allowance, adequate intake, estimated average requirement, and tolerable upper intake level, which vary with age, sex, and life stage.

\medterm{Dietary Sterol} A cholesterol-like substance found in foods, including cholesterol and plant sterols. Absorption and metabolism affect blood cholesterol; plant sterols can reduce intestinal cholesterol absorption in some people.

\medterm{Dietary Supplement} A product taken by mouth to provide vitamins, minerals, herbs, amino acids, or other dietary ingredients. Supplements may be useful for a documented need but can interact with medicines or cause toxicity.

\medterm{Dietary Supplementation} Deliberate addition of nutrients or other substances to the diet to correct a deficiency, meet increased needs, or support a health goal. The dose and need should be assessed because excess can be harmful.

\medterm{Dietary Supplements} Products containing vitamins, minerals, herbs, amino acids, or other substances intended to add to the diet. They do not replace balanced food and may interact with medicines or cause adverse effects.

\synonyms
Nutritional supplements

\medterm{Dietetics} The science and professional practice of using food and nutrition to promote health, prevent disease, and manage medical conditions through individualized dietary advice.

\medterm{Diethylcarbamazine} An antiparasitic medicine used mainly for lymphatic filariasis and some other filarial infections. Fever, rash, swelling, or inflammation may occur as the immune system reacts to dying parasites.

\medterm{Diethylhexylphthalate} Another name for DEHP, a chemical added to some plastics to make them flexible. Repeated or high exposure may affect reproduction, development, or hormone function.

\medterm{Diethylpropion} A stimulant appetite suppressant used short term in selected weight-management plans. It can cause insomnia, anxiety, fast heartbeat, high blood pressure, misuse, and dependence; medical supervision is required.
\medterm{Diethylstilbestrol} A synthetic estrogen formerly given to some pregnant people. Exposure before birth can increase the risk of reproductive-tract abnormalities, fertility problems, and certain rare cancers later in life.

\medterm{Dietitian} A nutrition professional who assesses food and nutrient needs and provides individualized advice for health, disease prevention, and medical conditions. Professional titles and licensing rules vary by region.

\medterm{Dieulafoy Lesion} An unusually large artery close to the surface of the stomach or another digestive organ that can suddenly cause severe bleeding, sometimes without a visible ulcer. It may produce vomiting of blood or black stools.

\medterm{Diff Dx} An abbreviation for differential diagnosis, the structured comparison of possible causes of a person’s symptoms, signs, or test results.

\medterm{Differences Of Sex Development} Variations present from birth in chromosomes, gonads, hormones, or internal or external genital anatomy. Medical evaluation explains the variation, addresses health needs, and supports informed decisions and well-being.

\medterm{Differential Blood Count} A blood test measuring the numbers or percentages of different white blood cells, including neutrophils, lymphocytes, monocytes, eosinophils, and basophils. Results help assess infection, inflammation, allergy, and blood disorders.

\medterm{Differential Diagnosis} The process of comparing possible diseases or conditions that could explain the same symptoms or findings. Clinicians consider serious causes first, use examination and targeted tests, and revise the possibilities as new information appears.

\medterm{Differentiated Cancer} A cancer whose cells still resemble the normal tissue from which they arose. The degree of resemblance, called differentiation, helps pathologists describe tumor grade and may provide information about behavior.

\medterm{Differentiation} The process by which cells develop specialized structures and functions. In cancer pathology, it describes how closely abnormal cells resemble the normal tissue from which they originated.

\medterm{Differentiation Antigens} Molecules on a cell’s surface that indicate its tissue type or stage of development. Laboratory tests for them help identify blood-cell disorders, classify cancers, and guide selected treatments.


\medterm{Difficile Clostridium (Clostridium Difficile Colitis)} Alternate index label for Clostridium Difficile Colitis.

\medterm{Difficult Airway} A situation in which a clinician may have trouble keeping a person breathing or placing a device into the airway. Causes include unusual anatomy, swelling, injury, illness, or limited neck movement and require careful planning.

\medterm{Difficult-to-Control Asthma} Asthma that remains poorly controlled despite appropriate treatment. Possible reasons include ongoing triggers, incorrect inhaler use, another illness, an incorrect diagnosis, severe airway inflammation, or a different response to medicines. Persistent symptoms need clinical review.

\medterm{Difficult Intubation} Difficulty placing a breathing tube in the windpipe despite appropriate technique and equipment. Limited mouth or neck movement, unusual anatomy, swelling, injury, obesity, or a previous difficult intubation may increase risk.

\medterm{Difficult Urination} Trouble starting or maintaining urine flow, a weak stream, pain, straining, or a feeling of incomplete emptying. Causes include infection, blockage, prostate or pelvic disorders, medicines, and nerve or bladder problems.

\medterm{Difficulty Speaking} Alternate index label; see \textit{Voice disorders}.

\medterm{Difficulty Swallowing} Trouble moving food or liquid safely from the mouth to the stomach. It may cause coughing, choking, pain, drooling, weight loss, or a feeling that food is stuck in the throat or chest.

\medterm{Difficulty Swallowing (Swallowing)} Alternate index label for Swallowing.

\medterm{Difficulty Walking} Trouble walking safely or normally because of pain, weakness, poor balance, numbness, joint disease, neurologic illness, or heart or lung limitations. Sudden onset requires urgent assessment.

\medterm{Diffuse} Spread broadly or throughout an area rather than limited to one small, clearly defined focus. The term may describe a disease, pain, swelling, imaging finding, or other process.

\medterm{Diffuse Astrocytoma} An older name for an infiltrating brain tumor made of astrocyte-like cells. Modern classification depends on molecular findings, including IDH status, which help determine diagnosis, outlook, and treatment.

\medterm{Diffuse Axonal Injury} Widespread damage to brain nerve fibers caused by rapid acceleration, deceleration, or rotation of the head. It can cause prolonged unconsciousness or severe brain dysfunction, and MRI may show the injury.

\medterm{Diffuse Esophageal Spasm} A disorder in which the esophagus contracts too early or unevenly, causing intermittent chest pain and difficulty swallowing liquids or solids. Testing helps distinguish it from blockage, reflux, and heart disease.

\medterm{Diffuse Idiopathic Skeletal Hyperostosis} A condition in which extra bone forms along spinal ligaments and where tendons attach to bone. It can cause stiffness, pain, swallowing difficulty, or fractures after minor injury.

\synonyms
DISH (Diffuse Idiopathic Skeletal Hyperostosis)

\medterm{Diffuse Large B-Cell Lymphoma} A fast-growing cancer of B lymphocytes, a type of white blood cell. It may cause painless swollen lymph nodes, fever, night sweats, or weight loss and usually requires prompt specialist treatment.

\medterm{Diffuse Midline Glioma} An aggressive brain or spinal-cord cancer arising in a midline structure, often in children. Symptoms depend on location and may include weakness, movement problems, vision changes, vomiting, or seizures.
\synonyms
DMG

\medterm{Diffuse Parenchymal Lung Diseases} Alternate index label; see \textit{Interstitial Lung Disease}.

\medterm{Diffuse Scleroderma} A form of systemic sclerosis in which skin thickening extends beyond the hands and face onto the trunk or limbs. It carries a higher risk of lung, kidney, heart, blood-vessel, and digestive complications.

\medterm{Diffusing Capacity} The ability of the lungs to move oxygen and other gases from the air sacs into the blood. It may be reduced by emphysema, interstitial lung disease, pulmonary vascular disease, or anemia.

\medterm{Diffusing Capacity Of The Lung} A measure of how efficiently gas crosses from the lung air sacs into nearby blood vessels. It is commonly measured as DLCO and helps evaluate lung tissue, blood vessels, and hemoglobin-related problems.

\medterm{Diffusion} Movement of molecules from an area where they are more concentrated to an area where they are less concentrated, caused by their random motion. It helps exchange gases and substances in the body.

\medterm{Diffusion Anisotropy} A difference in the direction water moves through tissue, especially along nerve fibers in the brain. MRI can use this pattern to study nerve-fiber structure and injury.

\medterm{Diffusion Capacity} A measure of how efficiently gases pass from lung air sacs into blood. It is commonly reported as DLCO and may be reduced in interstitial lung disease, emphysema, pulmonary vascular disease, or anemia.

\medterm{Diffusion Restriction} An area where water movement appears reduced on a special MRI sequence. It can occur with an acute stroke, abscess, densely packed tumor, seizure, or other problem and must be interpreted with the full examination.

\medterm{Diffusion-Weighted Imaging} An MRI technique that shows how water molecules move within tissue. Areas with restricted movement may occur in acute stroke, abscess, densely packed tumors, seizures, or other conditions.

\medterm{Diflunisal} A salicylate nonsteroidal anti-inflammatory medicine used for pain and inflammation. It can cause stomach irritation or bleeding, kidney injury, fluid retention, allergic reactions, or cardiovascular problems.

\medterm{Digastric Muscle} A two-part muscle beneath the jaw that helps open the mouth, lower the jaw, and raise the hyoid bone during swallowing and speech.

\medterm{DiGeorge Syndrome} A genetic disorder usually caused by loss of a small part of chromosome 22. It can affect the heart, immune system, calcium regulation, palate, hearing, learning, and development.

\synonyms
22q11.2 deletion syndrome

\medterm{DiGeorge Syndrome (22Q11.2 Deletion Syndrome)} Alternate name; see \textit{DiGeorge Syndrome}.

\medterm{Digestion} The process of breaking food into smaller substances that the body can absorb and use. It begins in the mouth and continues through the stomach and intestines with movement, acid, bile, and digestive enzymes.

\medterm{Digestive Diseases} Disorders affecting the digestive tract, liver, gallbladder, or pancreas. Examples include reflux, ulcers, infections, inflammatory bowel disease, blockage, malabsorption, and cancer.

\medterm{Digestive Disorder} A disease affecting the digestive tract or the process of digestion. It may cause abdominal pain, nausea, reflux, diarrhea, constipation, bleeding, vomiting, or poor nutrient absorption.

\medterm{Digestive Health} Healthy function of the digestive tract in breaking down food, absorbing nutrients and water, and removing waste. It is influenced by diet, fluids, medicines, movement, illness, and the gut lining and microbiome.

\synonyms
Gut Health

\medterm{Digestive System} The organs that process food, absorb nutrients and water, and remove solid waste. They include the mouth, esophagus, stomach, intestines, liver, gallbladder, pancreas, rectum, and anus.

\medterm{Digestive Tract} The continuous hollow passage through which food travels from the mouth to the anus. It includes the throat, esophagus, stomach, small intestine, large intestine, rectum, and anus.
\synonyms
Gastrointestinal tract; alimentary canal

\medterm{Digit} A finger or toe. Most fingers and toes contain three small bones, while the thumb and great toe usually contain two; muscles, tendons, nerves, and blood vessels allow movement and sensation.

\medterm{Digital Anorectal Exam - DARE} An examination in which a clinician uses a lubricated, gloved finger to feel the anal canal and rectum. It may assess tenderness, lumps, bleeding, muscle tone, and, when relevant, the prostate.

\medterm{Digital Clubbing} Enlargement and rounding of the fingertips with curved nails and loss of the normal angle at the nail base. It may be associated with lung, heart, digestive, liver, or other diseases and needs evaluation.

\medterm{Digital Fibrokeratoma} A rare, harmless, firm skin growth usually found on a finger or toe. It may resemble a wart; removal or microscopic examination may be needed when its appearance is unusual or it causes pain or interference.

\medterm{Digital Ischemia} Too little blood flow to a finger or toe, causing pain, pallor, coolness, numbness, or tissue death. It may result from arterial blockage, severe Raynaud phenomenon, inflammation, or an embolus and can be an emergency.

\medterm{Digital Mammography} Breast imaging that records X-ray pictures electronically. It helps detect masses, distortions, and calcium deposits and allows images to be enlarged, adjusted, stored, and compared.

\medterm{Digital Mucous Cyst} A usually harmless fluid-filled lump near the last joint of a finger, often linked to osteoarthritis. It may press on the nail, cause pain, or occasionally become infected.

\medterm{Digital Papillary Adenocarcinoma} A rare cancer of a sweat gland, usually on a finger, hand, or foot. It often appears as a slow-growing lump but can return locally or spread to other organs.

\medterm{Digital Rectal Exam} A physical examination using a lubricated, gloved finger to feel the anal canal and rectum and, when relevant, the prostate. It may detect tenderness, bleeding, lumps, hemorrhoids, or muscle problems.

\medterm{Digital Rectal Examination} An examination in which a clinician gently feels inside the rectum with a lubricated, gloved finger. It may check for pain, bleeding, lumps, hemorrhoids, prostate changes, or muscle problems after explanation and consent.

\medterm{Digital Subtraction Angiography} An X-ray study in which computer processing removes background structures from contrast images so blood vessels are clearer. It can diagnose narrowing or blockage and guide selected vascular procedures.

\medterm{Digitalis (Digoxin) Toxicity} Poisoning from digoxin or a related cardiac glycoside. Nausea, vomiting, confusion, visual changes, slow or fast abnormal rhythms, and dangerous potassium changes may occur; severe cases require emergency treatment.

\synonyms
Digoxin toxicity

\medterm{Digitalis Preparation} A preparation containing a cardiac glycoside such as digoxin or digitoxin. It can strengthen heart contraction and slow electrical conduction, but the dose range is narrow and toxicity can cause digestive, visual, or rhythm problems.

\medterm{Digitalis Toxicity} Harmful effects from digitalis or related cardiac glycosides, including nausea, vomiting, visual changes, confusion, slow or fast abnormal heart rhythms, and high potassium. Kidney disease and drug interactions increase risk.

\medterm{Digitate Warts (Verruca Digitata)} Finger-like or thread-like warts, usually on the face or neck, caused by human papillomavirus infection. They are generally harmless but can spread or recur.

\synonyms
Verruca Digitata

\medterm{Digitoxin} A cardiac glycoside that can strengthen heart contraction and slow electrical conduction. Its safe and toxic doses are close together, so abnormal rhythms, digestive symptoms, visual changes, and drug interactions are important concerns.

\medterm{Digoxin} A cardiac medicine used in selected people with heart failure or atrial fibrillation. It can strengthen contraction and slow the heart rate, but kidney function, electrolytes, other medicines, and blood levels affect toxicity risk.

\medterm{Digoxin Immune Fab} Antibody fragments given by injection that bind digoxin and related cardiac glycosides, reducing their effects. They are used for life-threatening abnormal rhythms, severe toxicity, or a potentially very large ingestion.

\medterm{Digoxin Test} A blood test measuring digoxin concentration to help guide dosing or assess possible toxicity. The result must be interpreted with timing, symptoms, kidney function, electrolytes, and interacting medicines; the number alone does not prove toxicity.

\medterm{Dihydrocodeine} An opioid medicine used for selected pain or cough conditions. It can cause constipation, nausea, drowsiness, dependence, dangerous breathing depression, and overdose, especially with alcohol or other sedatives.

\medterm{Dihydroergotamine} An ergot medicine used for selected acute migraine attacks. It narrows certain blood vessels and must be avoided in several vascular, heart, liver, kidney, and pregnancy-related conditions.

\medterm{Dihydrotestosterone} A potent androgen made from testosterone. It contributes to development of male external genitalia, facial and body hair, prostate growth, and other androgen effects.
\medterm{Dihydroxyacetone} A three-carbon sugar involved in metabolism. Its phosphate form is an intermediate used when cells break down glucose or make glucose.

\medterm{Dilantin Overdose} Taking too much phenytoin can cause involuntary eye movements, poor coordination, slurred speech, drowsiness, confusion, low blood pressure, abnormal heart rhythms, seizures, or coma.

\medterm{Dilatation} Widening or expansion of a hollow organ, blood vessel, duct, opening, or other body structure. It may be normal, caused by disease, or deliberately produced during a procedure.

\medterm{Dilatation And Curettage (D \& C)} A procedure that opens the cervix and removes tissue from inside the uterus using suction or an instrument. It may treat miscarriage or retained pregnancy tissue, investigate bleeding, or remove abnormal tissue; risks include bleeding, infection, and perforation.

\medterm{Dilate} To become wider or make a passage, blood vessel, opening, or body structure wider. Dilation may occur naturally, result from medicines or disease, or be created during a procedure.

\synonyms
Expand; enlarge

\medterm{Dilated Cardiomyopathy} A heart-muscle disease in which one or more chambers enlarge and pump weakly. Causes include inherited disease, infection, toxins, medicines, pregnancy, blocked arteries, and inflammation; it can lead to heart failure or abnormal rhythms.

\synonyms
Cardiomyopathy, Dilated

\medterm{Dilated Cardiomyopathy Prognosis} Outlook varies with heart-pumping strength, symptoms, cause, abnormal rhythms, other illnesses, and response to treatment. Medicines and devices can improve function for some people; severe disease may require advanced heart-failure care.

\medterm{Dilation} Widening or expansion of a hollow organ, blood vessel, duct, pupil, opening, or other body structure. It may be normal, caused by disease, or deliberately produced during a procedure.

\medterm{Dilation And Curettage} A procedure that opens the cervix and removes tissue from inside the uterus using suction or an instrument. It may treat miscarriage or retained pregnancy tissue, investigate abnormal bleeding, or remove abnormal uterine tissue.

\medterm{Dilation And Evacuation} A procedure that opens the cervix and removes pregnancy tissue from the uterus using suction and instruments. It is performed by trained clinicians for selected pregnancy-related indications, with risks including bleeding, infection, and uterine injury.

\medterm{Dilator} An instrument or device used to widen a body opening or passage gradually or temporarily. It may be used during examinations, surgery, childbirth care, dental treatment, or treatment of a narrowing.

\medterm{Diltiazem} A calcium-channel blocker used for high blood pressure, angina, and heart-rate control in some abnormal rhythms. It can slow the heart, lower blood pressure, cause constipation, or worsen heart failure in susceptible people.

\medterm{Dilute Russell Viper Venom Time} A blood-clotting test used mainly to detect lupus anticoagulant, an antibody associated with antiphospholipid syndrome. Anticoagulant medicines and other clotting problems can affect the result, so confirmatory testing is needed.

\medterm{Dilution} Lowering the concentration of a substance by adding solvent or another fluid. Accurate dilution is important when preparing medicines, laboratory reagents, standards, or specimens.

\medterm{Dimenhydrinate} An antihistamine used to prevent or treat motion sickness and nausea. It can cause drowsiness, dry mouth, blurred vision, constipation, confusion, or difficulty urinating.

\medterm{Dimenhydrinate Overdose} Taking too much dimenhydrinate can cause severe drowsiness or agitation, confusion, hallucinations, seizures, abnormal heart rhythms, breathing depression, or coma and requires urgent medical assessment.

\medterm{Dimercaprol} An injected medicine that binds certain toxic metals, including arsenic, mercury, and gold, so the body can remove them. It can cause high blood pressure, fast heartbeat, nausea, pain, and other adverse effects.

\medterm{Dimesna} A sulfur-containing compound studied as a protective or detoxifying medicine, including in some chemotherapy-related settings. Its clinical role depends on the specific treatment and available evidence.
\medterm{Dimethoate} An organophosphate insecticide that blocks acetylcholinesterase. Poisoning can cause sweating, excess saliva, vomiting, diarrhea, small pupils, muscle weakness, breathing problems, seizures, or coma.

\medterm{Dimethoxon} A poisonous breakdown product of some organophosphate pesticides. It can cause sweating, excess saliva, vomiting, diarrhea, small pupils, muscle twitching, weakness, breathing difficulty, or seizures and requires urgent poisoning advice.

\medterm{Dimethyl Fumarate} An oral medicine used for multiple sclerosis and selected skin conditions. It can cause flushing, digestive symptoms, low white-cell counts, liver injury, allergic reactions, or rare serious infections.

\medterm{Dimethyl Fumarate Allergy} An allergic or hypersensitivity reaction to dimethyl fumarate, a medicine used for some forms of multiple sclerosis. It may cause rash, itching, swelling, breathing difficulty, or other immune symptoms and requires medical assessment.

\medterm{Dimethyl Sulfoxide} An organic solvent used in laboratories and some specialized medical or investigational applications. Contact or ingestion can irritate tissues, and it may carry other substances through the skin.

\medterm{Dimethyl Sulphoxide} Another spelling of dimethyl sulfoxide, an organic solvent used in laboratories and selected medical or investigational applications. Exposure may cause irritation, a characteristic odor, or other adverse effects.


\medterm{Dimethylhydrazines} A group of highly toxic chemicals formerly used mainly in rocket fuel. Exposure can irritate or burn tissues and damage the nervous system, liver, blood, or other organs.

\medterm{Dimetilan} A poisonous organophosphate insecticide that disrupts nerve signals by blocking acetylcholinesterase. Exposure can cause sweating, vomiting, diarrhea, small pupils, weakness, breathing difficulty, seizures, or coma and needs urgent help.

\medterm{Diminished Ovarian Reserve} Fewer remaining eggs than expected for a person’s age, which may reduce response to fertility treatment and shorten the time available for conception. Tests estimate reserve but cannot precisely predict natural fertility.

\synonyms
Decreased Ovarian Reserve; DOR

\medterm{Dimocillin} A penicillin-type antibiotic used historically or in selected settings for susceptible bacterial infections. Choice depends on the organism, resistance testing, infection site, allergies, and local treatment guidance.

\medterm{Dimorphic} Existing in two distinct forms. Some fungi are dimorphic, growing as a mold in the environment and as a yeast in body tissues; this ability can contribute to infection.


\medterm{Dinophyceae} A group of aquatic microorganisms, including dinoflagellates. Some produce toxins during harmful algal blooms that can contaminate seafood or water and cause neurologic, digestive, or liver illness.

\medterm{Dinoprost} A prostaglandin medicine that stimulates uterine contraction and other smooth-muscle effects. It is used mainly in veterinary medicine and selected reproductive-health settings; pain, diarrhea, breathing difficulty, and blood-pressure changes may occur.

\medterm{Dinoprostone} A prostaglandin medicine used to soften the cervix and help start labor or manage selected pregnancy-related conditions. It can cause contractions that are too frequent, pain, fever, nausea, or fetal heart-rate changes.

\medterm{Dinoprostone (Cervidil, Prepidil)} A prostaglandin medicine used to soften the cervix and help start labor. It requires obstetric monitoring because it can cause overly frequent contractions, pain, fever, or changes in the fetal heart rate.

\synonyms
Cervidil, Prepidil

\medterm{Dinucleoside Phosphates} Molecules containing two nucleosides joined by phosphate groups. They occur in nucleotide metabolism and may take part in cellular signaling, energy transfer, or DNA and RNA-related reactions.

\medterm{Dinutuximab} An antibody-based cancer medicine that attaches to a marker found on many neuroblastoma cells and helps the immune system destroy them. It can cause severe nerve pain, allergic reactions, fever, and fluid leakage from blood vessels.

\medterm{Dioctyl Sodium Sulphosuccinate} A surfactant and stool-softening medicine also called docusate sodium. It helps water mix with stool and is used in selected constipation settings, although evidence of benefit is limited.

\medterm{Diopter} A unit of optical power equal to the reciprocal of focal length in meters. Eyeglass and contact-lens prescriptions use diopters to state the strength needed to focus light clearly.
\medterm{Dioxin} A group of long-lasting environmental chemicals produced by some industrial processes and burning. They can build up in body fat; long-term exposure may affect development, reproduction, immunity, hormones, and cancer risk.

\medterm{Dipeptide} A molecule made of two amino acids joined together. Dipeptides are produced during protein digestion and metabolism and can be absorbed or used in further chemical reactions.

\medterm{Dipeptidyl Peptidase-4 Inhibitors} Oral diabetes medicines that prolong the action of natural gut hormones. They help the body release insulin after meals and reduce glucagon, usually without causing weight gain or low blood glucose when used alone.


\medterm{Dipetalonema Infection} Infection with a Dipetalonema parasitic worm, usually acquired through an arthropod bite. It may affect skin or tissue beneath the skin; diagnosis and treatment depend on the species and clinical findings.

\medterm{Diphallia} A rare birth condition in which the penis is partly or completely duplicated. Abnormalities of the urinary tract, anus, spine, or other organs may also occur and require specialist assessment.

\medterm{Diphenhydramine Overdose} Taking too much diphenhydramine can cause severe drowsiness or agitation, confusion, hallucinations, seizures, abnormal heart rhythms, breathing depression, coma, or death and requires urgent medical assessment.

\medterm{Diphenoxylate} An opioid-related medicine that slows intestinal movement to reduce diarrhea, usually combined with atropine. Excessive use can cause drowsiness, constipation, dependence, breathing depression, or poisoning.

\medterm{Diphenylhydramine} An antihistamine with sedating and anticholinergic effects used in some allergy or sleep products. It can impair alertness and cause dry mouth, confusion, urinary retention, abnormal rhythms, or toxicity in overdose.

\medterm{Diphosphate} A chemical group containing two linked phosphate units. Diphosphate is involved in energy use, nucleotide formation, and other chemical reactions in cells.

\medterm{Diphosphonate} A chemical group found in compounds that can bind strongly to bone mineral. Some related medicines are used or studied for disorders of bone breakdown and calcium balance.

\medterm{Diphtheria} A serious infection caused by toxin-producing Corynebacterium bacteria. It can form a thick coating in the throat and cause airway blockage, heart inflammation, nerve damage, or death; vaccination prevents most cases.


\medterm{Diphtheria Toxin} A poisonous protein made by some Corynebacterium bacteria. It blocks protein production inside cells, causing local tissue injury and potentially severe heart, nerve, and airway complications.

\medterm{Diphtheroid Bacillus} A rod-shaped bacterium resembling Corynebacterium under the microscope. Many are normal skin bacteria or contaminants, but some can cause opportunistic infection in people with serious illness or medical devices.

\medterm{Diphyllobothriasis} An intestinal infection caused by a fish tapeworm acquired from raw or undercooked freshwater fish. It may cause abdominal symptoms and, rarely, vitamin B12 deficiency and anemia.

\medterm{Diphyllobothrium Latum} A fish tapeworm acquired by eating raw or undercooked freshwater fish. It lives in the intestine and may cause abdominal symptoms or compete for vitamin B12, occasionally leading to anemia.

\medterm{Diplacusis} Hearing one sound as two different pitches or qualities. It often reflects unequal hearing between the ears or inner-ear dysfunction and may follow infection, noise injury, or other ear disease.

\medterm{Diplegia} Severe weakness or paralysis affecting matching parts on both sides of the body. In spastic diplegia, the legs are affected more than the arms, often because of early brain injury.

\medterm{Diplococcus} A bacterium that appears as pairs of round cells under a microscope. The term describes shape, not one species; some diplococci cause infections in people.

\medterm{Diploid Cell} A cell containing two complete sets of chromosomes, one inherited from each parent. Most human body cells are diploid; eggs and sperm normally contain only one set.

\medterm{Diploidy} Having two complete sets of chromosomes, usually one from each parent. This is the normal chromosome state of most human body cells.

\medterm{Diplopia} Seeing two images of one object. It may result from eye misalignment, an eye-muscle or nerve problem, a brain disorder, or an eye-surface or lens problem; sudden onset needs urgent assessment.

\medterm{Diplopia (Double Vision)} Seeing two images of one object. It may occur with one eye or only when both eyes are open; causes include dry eye, cataract, eye misalignment, nerve disease, muscle disease, or brain disease.

\synonyms
Double Vision

\medterm{Dipole} A pair of equal and opposite electrical or magnetic charges separated by a distance. The concept is used in physiology, medical imaging, and physics to describe electrical or magnetic fields.

\medterm{Dipper} A person whose blood pressure falls during sleep, usually by about 10 percent or more compared with daytime readings. The pattern is measured with a portable 24-hour monitor and may help assess cardiovascular risk.
\medterm{Diprosopus} A rare birth defect in which part or all of the face is duplicated. It usually occurs with severe abnormalities of the brain, heart, or other organs and often causes major developmental problems.

\medterm{Dipsomania} Alternate historical term; see \textit{Alcohol Use Disorder}.

\medterm{Dipsosis} Excessive thirst or a persistent urge to drink. Causes include dehydration, high blood glucose, high blood concentration, medicines, dry mouth, or disorders affecting thirst regulation.

\medterm{Dipstick Device} A strip or small device containing chemical pads that change color when placed in a specimen, commonly urine. It provides a rapid screen for substances such as blood, protein, glucose, or infection markers.

\medterm{Diptera} The insect order containing flies and mosquitoes. Some species spread infections through bites or contamination, while others can cause infestations such as myiasis.

\medterm{Dipylidium Caninum} A tapeworm of dogs and cats that can occasionally infect people, especially children, after an infected flea is swallowed. Infection may cause mild digestive symptoms or visible worm segments in stool.

\medterm{Dipyridamole} A medicine that reduces platelet clumping and widens blood vessels. It is used in selected stroke-prevention plans and heart stress testing; headache, dizziness, low blood pressure, and bleeding may occur.

\medterm{Direct Action} An effect produced through an intervention’s immediate target, without requiring an intermediate pathway or mediator. The term may describe a medicine, hormone, nerve signal, or other biological process.

\medterm{Direct Antiglobulin Test} A blood test that detects antibodies or complement attached to a person’s red blood cells. A positive result may occur with autoimmune hemolysis, a transfusion reaction, or hemolysis affecting a newborn.

\medterm{Direct Coombs Test} Another name for the direct antiglobulin test, which detects antibodies or complement attached to red blood cells. It can help investigate autoimmune hemolysis, transfusion reactions, medicine-related hemolysis, or newborn hemolysis.

\medterm{Direct Hernia} An inguinal hernia in which abdominal tissue pushes through a weak area in the back wall of the groin canal. It may cause a groin lump or discomfort and can occasionally become trapped.


\medterm{Direct Oral Anticoagulants} Oral blood thinners that directly block a clotting protein, such as thrombin or factor Xa. They prevent and treat clots but can cause bleeding and require attention to kidney function, interactions, and procedures.

\medterm{Direct Transmission} Spread of an infectious organism through immediate contact or short-range respiratory droplets from an infected person or animal, without an intermediate object, food, water source, or vector.

\medterm{Direct Veneer} A tooth-colored composite layer shaped and hardened directly on a tooth to improve appearance or repair minor damage. It usually preserves more tooth than some alternatives but may stain, chip, or need replacement.


\synonyms
Direct composite veneer; composite dental veneer

\medterm{Directly Observed Therapy} A healthcare worker or trained supporter watches a person take each dose of medicine. It can support completion of complex treatment, especially for tuberculosis, and helps detect adverse effects or missed doses.

\medterm{Dirofilaria Repens} A mosquito-transmitted parasitic worm that rarely infects people. It usually causes a single lump beneath the skin or near the eye because humans are an accidental host; it does not usually mature or spread in the body.

\medterm{Dirt - Swallowing} Repeatedly eating soil or other nonfood substances, called geophagia and sometimes part of pica. It can expose a person to parasites, bacteria, toxic chemicals, or excess minerals and may cause constipation or blockage.

\medterm{Dirt-Eating} Repeatedly eating soil or other nonfood substances, a form of pica. It may be associated with iron deficiency, pregnancy, nutritional problems, or mental health conditions and can expose a person to infection or toxins.

\medterm{Dirty Necrosis} Dead tumor tissue mixed with broken cell material and many white blood cells. It is often seen in rapidly growing cancers and is interpreted by a pathologist with the rest of the tissue sample.

\medterm{Disability} A difficulty with body function, activity, or participation caused by a health condition and influenced by the person’s environment. Disability may be temporary or lifelong and varies in severity and support needs.

\medterm{Disability Evaluation} An assessment of how a health condition affects physical, mental, or daily functioning. History, examination, records, and appropriate tests may inform decisions about work, benefits, insurance, rehabilitation, or support.

\medterm{Disability-Adjusted Life Year} A population-health measure combining years lost from early death with years lived in less than full health. One DALY represents one lost year of healthy life under the method used.

\medterm{Disaccharidase} An intestinal enzyme that breaks a double sugar into simple sugars that can be absorbed. Lactase, for example, breaks lactose in milk into glucose and galactose.

\medterm{Disaccharide} A carbohydrate made of two linked simple sugars, such as lactose, sucrose, or maltose. It must be split by digestive enzymes before the intestine can absorb its component sugars.

\medterm{Disarticulation} Amputation performed through a joint rather than by cutting through the shaft of a bone. The level is chosen according to the injury, disease, circulation, function, and possibility of healing.

\medterm{Disc} A thin, round structure. In the spine, an intervertebral disc is a tough, flexible cushion between neighboring vertebrae that absorbs load and allows movement.

\medterm{Discectomy} Surgery to remove part of a damaged spinal disc that is pressing on a nerve. It may relieve leg or arm pain and weakness, but risks include infection, bleeding, nerve injury, and recurrent disc problems.


\medterm{Discitis} Inflammation or infection of a spinal disc, often involving the nearby vertebral bones. It may cause severe localized back or neck pain, stiffness, fever, or reduced movement and usually requires imaging and medical evaluation.


\medterm{Discography} An imaging procedure in which contrast is injected into a spinal disc to examine it and see whether the injection reproduces pain. Because it has limitations and risks, it is used selectively.

\medterm{Discoid} Coin-shaped or round. In medicine, it often describes a skin lesion or rash, such as the round plaques of discoid lupus; the associated condition must be identified separately.

\medterm{Discoid Lupus} A long-lasting form of skin lupus causing round, inflamed, scaly plaques. Lesions may scar, change skin color, or cause permanent hair loss on the scalp; some people later develop systemic lupus.

\medterm{Discoid Lupus Erythematosus} A chronic autoimmune skin disease causing persistent round, inflamed plaques that may scale, scar, or change pigment. It commonly affects the face, ears, and scalp and may cause permanent hair loss.
\medterm{Discoloration} An unusual change in the color of skin, mucous membranes, nails, or another tissue. Causes include bruising, low oxygen, jaundice, inflammation, infection, or altered pigment; sudden, painful, widespread, or unexplained change needs assessment.


\medterm{Discomfort} An unpleasant physical or emotional sensation that may be milder or less clearly located than pain. Its importance depends on the site, duration, associated symptoms, and underlying cause.

\medterm{Discordance} Lack of agreement between findings, measurements, genetic results, or the medical conditions of two related people. The significance depends on what is being compared and the clinical setting.

\medterm{Discordant} Not agreeing, matching, or corresponding. In medicine it may describe conflicting test results or different medical findings in related people, twins, or paired organs.

\medterm{Discordant Couple} A couple whose partners have different results or medical statuses relevant to their care, such as one partner having HIV and the other not. Counseling can reduce transmission risk and support reproductive decisions.


\medterm{Disease} An abnormal condition that affects the body or mind and changes normal structure or function. Diseases may be infectious, inflammatory, inherited, degenerative, metabolic, immune-related, or cancerous.


\medterm{Disease, Osgood-Schlatter} An overuse injury of the growing area where the knee tendon attaches to the shinbone. It causes pain and swelling below the knee, usually worse with running, jumping, or kneeling.

\medterm{Disease Progression} Change in a disease over time, often meaning worsening. It is assessed through symptoms, examination, laboratory results, imaging, function, complications, and response to treatment.

\medterm{Disease Reporting} Notification of specified diseases, infections, or health events to public-health authorities. Reports help detect outbreaks, protect contacts, arrange follow-up, and monitor health patterns; requirements vary by location.

\medterm{Disease Reservoir} A person, animal, environment, or material in which an infectious organism normally lives and from which it can spread to a susceptible host.

\medterm{Disease Response} The body’s biological or clinical reaction to a disease or its treatment. It is assessed using symptoms, signs, laboratory results, imaging, physical function, or survival.

\medterm{Disease, Subclinical} Disease that causes no obvious symptoms or signs but can be detected by laboratory tests, imaging, examination, or another evaluation.

\medterm{Disease Surveillance} Ongoing collection, analysis, interpretation, and reporting of health information to detect, monitor, and respond to disease patterns, outbreaks, and changes in population health.

\medterm{Disease Susceptibility} Increased likelihood of developing a disease because of genes, age, immune status, exposures, behavior, other illnesses, or environmental factors. Susceptibility raises risk but does not mean disease is certain.

\medterm{Disease Vector} A living carrier, commonly a mosquito, tick, flea, or other arthropod, that transfers an infectious organism between hosts. The vector may transmit infection without becoming ill itself.

\medterm{Disease-Free Survival} The length of time after treatment during which a person remains alive without detectable evidence of the disease. The exact definition and assessment method depend on the disease and study.

\medterm{Disfigurement} A visible change that substantially alters appearance, caused by injury, disease, surgery, a birth difference, or scarring. Its physical, emotional, and social effects vary between people.

\medterm{Disinfectant} A chemical or physical agent used on nonliving surfaces to destroy or inactivate many microorganisms. Disinfectants are not necessarily safe for skin, wounds, food, or internal use.

\medterm{Disinfection} Cleaning that destroys or inactivates most disease-causing microorganisms on nonliving surfaces. It may not destroy bacterial spores; the required level depends on the object, organism, and intended use.

\medterm{Disk} A flat, round structure. In the spine, an intervertebral disk lies between neighboring vertebrae, absorbs shock, and helps the back move.

\synonyms
Disc

\medterm{Disk Diffusion Susceptibility Test} A laboratory test of bacterial antibiotic sensitivity. Drug-containing disks are placed on growing bacteria, and the clear areas around them are measured and compared with standards to guide treatment.

\medterm{Disk Replacement - Lumbar Spine} Surgery replacing a damaged lower-back disk with an artificial disk in selected people. It aims to preserve movement and reduce pain, but risks include nerve injury, infection, persistent pain, and device problems.

\medterm{Diskectomy} Surgery to remove part or all of a damaged spinal disk pressing on a nerve or the spinal cord. It may relieve pain or weakness caused by disk herniation; risks include infection, bleeding, nerve injury, and recurrence.

\medterm{Diskitis} Inflammation or infection of a spinal disk and nearby vertebral bone. It commonly causes severe localized back pain, stiffness, fever, or reluctance to move and requires medical assessment and imaging.

\medterm{Dislocated Ankle (Ankle Dislocation)} Alternate index label for Ankle Dislocation.

\medterm{Dislocated Elbow} Loss of normal alignment of the elbow joint, usually after a fall or collision. It causes severe pain, deformity, and limited movement and needs urgent reduction with checks for fractures, nerve injury, and impaired circulation.

\medterm{Dislocated Hip} Displacement of the upper thighbone from its hip socket, usually after major trauma or around an artificial joint. It is an emergency because nearby nerves and the blood supply to the bone may be damaged.

\medterm{Dislocated Knee} Loss of alignment between the thighbone and shinbone, usually after severe injury. Major artery or nerve damage may be hidden, so urgent reduction and circulation checks are needed even if the knee returns to position.

\medterm{Dislocated Shoulder} Displacement of the upper arm bone from the shoulder socket. It causes pain, deformity, and inability to move normally and requires reduction with assessment of nerves, blood vessels, and fractures.


\medterm{Dislocation} Complete loss of normal contact between the surfaces of a joint, usually from injury or disease. It can cause severe pain and deformity and may occur with fractures or nerve and blood-vessel injury.

\medterm{Dislocation Of The Lens} Partial or complete movement of the eye’s natural lens away from its normal position because its supporting fibers are weak or torn. It may cause blurred or double vision, eye pain, or glaucoma.

\medterm{Disofenin Scan} A nuclear-imaging test that follows bile production and flow through the liver, gallbladder, and bile ducts. It can help detect blockage, leakage, inflammation, or selected causes of abdominal pain.

\medterm{Disopyramide} An antiarrhythmic medicine used selectively for certain abnormal heart rhythms. It can slow electrical conduction and weaken heart contraction and may cause dry mouth, urinary retention, low blood pressure, or new arrhythmias.

\medterm{Disorder} A disturbance of normal body or mind structure or function that causes symptoms, impairment, or important health risk. The term does not by itself indicate how severe the problem is.

\medterm{Disorder Antisocial Personality (Antisocial Personality Disorder)} Alternate index label for Antisocial Personality Disorder.

\medterm{Disorder Asperger (Asperger Syndrome)} Alternate index label for Asperger Syndrome.

\medterm{Disordered Thinking} A disturbance in the organization, clarity, logic, or flow of thoughts. It may occur with mental illness, brain disease, delirium, substance use, medicines, or other medical conditions.

\medterm{Disorders Sleep (Sleep)} Alternate index label for Sleep.



\medterm{Disorganized Schizophrenia} An older term for schizophrenia with markedly disorganized speech, behavior, or emotional expression. Modern classifications generally describe these features within schizophrenia-spectrum disorders rather than as a separate subtype.

\medterm{Disorientation} Difficulty knowing the time, place, situation, or one’s identity. It may occur with delirium, infection, abnormal body chemistry, intoxication or withdrawal, brain injury, or dementia; sudden onset needs urgent assessment.

\medterm{Dispensary} A place or service that stores, prepares, and supplies medicines and may provide basic medication advice. Safe dispensing includes checking the prescription, labeling the medicine, and explaining how to use it.

\medterm{Dispensed Amount} The quantity of medicine supplied to a patient. It is determined by the prescribed dose, treatment duration, formulation, refill rules, and applicable safety or legal requirements.

\medterm{Dispensing} The professional process of checking a prescription, selecting and preparing the medicine, labeling it, and supplying it with instructions for safe use.


\medterm{Displacement} Movement of an anatomical structure away from its normal position, as in a displaced fracture, joint injury, organ shift, or device migration. Effects depend on the structure involved and the amount of movement.

\medterm{Disruptive, Impulse-Control, And Conduct Disorders} A group of disorders involving persistent problems with emotional or behavioral self-control that violate the rights of others, major age-appropriate rules, or social norms. Examples include oppositional defiant disorder, conduct disorder, intermittent explosive disorder, kleptomania, and pyromania.

\medterm{Disruption} An interruption or breakdown of normal structure, function, continuity, or process. In medicine, it may describe damaged tissue, a broken barrier, disturbed development, or an interruption in the usual work of an organ or system.

\medterm{Dissect} To separate tissue layers or structures by cutting or careful blunt separation for surgery, examination, or anatomical study.

\medterm{Dissecting Aneurysm} Alternate index label; see \textit{Aortic dissection}.

\medterm{Dissection} Separation of tissue layers, either during surgery or because of disease. In an artery, a tear in the inner lining lets blood split the wall layers and can block blood flow or cause the artery to rupture.

\medterm{Dissection Aorta (Aortic Dissection)} Alternate index label for Aortic Dissection.

\medterm{Disseminate} To spread widely through the body, a tissue, or a population. Infection, cancer, inflammation, or another process may disseminate rather than remain in one limited area.

\medterm{Disseminated Coccidioidomycosis} A fungal infection in which Coccidioides spreads beyond the lungs to skin, bones, joints, brain coverings, or other organs. It can be severe and is more likely in people with impaired immunity.

\medterm{Disseminated Cutaneous Blastomycosis} A Blastomyces infection that spreads from the lungs or another site to multiple areas of skin. It may also involve bones or other organs and requires antifungal treatment.

\medterm{Disseminated Intravascular Coagulation} A life-threatening disorder in which clotting is activated throughout the bloodstream. Small clots can injure organs while using up platelets and clotting factors, causing serious bleeding at the same time.

\medterm{Disseminated Mycobacterium Avium–Intracellulare Infection} A widespread infection with Mycobacterium avium complex, usually in people with severe immune suppression. It may involve blood, bone marrow, lymph nodes, liver, spleen, or other organs.

\medterm{Disseminated Neuroblastoma} Neuroblastoma that has spread from its original site to distant organs, bone marrow, bones, liver, or lymph nodes. Treatment and outlook depend on age, spread, tumor biology, and response to specialist care.

\medterm{Disseminated Sclerosis} An older name for multiple sclerosis, an immune-mediated disease in which areas of damage develop in the brain, spinal cord, or optic nerves and cause symptoms that vary over time.

\medterm{Disseminated Tuberculosis} Tuberculosis that has spread through the bloodstream or lymphatic system to multiple organs. It may cause fever, weight loss, fatigue, organ damage, meningitis, or widespread small lesions and needs urgent treatment.

\medterm{Disseminated Varicella Infection} Chickenpox infection that spreads beyond the skin and may involve the lungs, brain, liver, or other organs. It is more likely and more severe in adults, pregnant people, and those with weakened immunity.

\medterm{Dissociation} A disruption in the usual connection between consciousness, memory, identity, emotions, perceptions, behavior, or sense of self. Mild forms may occur during daydreaming; severe or persistent forms can impair daily life.

\medterm{Dissociative Amnesia} Inability to recall important autobiographical information, often related to trauma or severe stress, beyond ordinary forgetfulness. Memory loss may involve a specific event, period, or broader part of a person’s life and is not better explained by ordinary forgetting, substances, a neurologic disorder, or another medical condition.

\medterm{Dissociative Disorder} Alternate index label for Dissociative Disorders.

\medterm{Dissociative Disorders} A group of mental disorders involving disruption in the usual integration of identity, memory, consciousness, perception, emotion, or behavior. Examples include dissociative amnesia, dissociative identity disorder, and depersonalization-derealization disorder. Temporary dissociative experiences can occur without constituting a disorder; a dissociative disorder causes significant distress or impairment.

\medterm{Dissociative Identity Disorder} A mental health disorder involving two or more distinct identity states and repeated gaps in memory for everyday events or personal information. It can cause changes in behavior, emotions, perception, and daily functioning.

\medterm{Distal} Farther from the point where a body part attaches or begins, or farther from the center of the body. For example, the hand is distal to the elbow.

\medterm{Distal Median Nerve Dysfunction} Abnormal function of the median nerve near the forearm, wrist, or hand, causing numbness, tingling, pain, or weakness in its area of supply. Compression at the wrist is a common cause.

\medterm{Distal Pancreatectomy} Surgery removing the body or tail of the pancreas, usually while leaving the head. It may be used for selected tumors, cysts, injury, or inflammation; diabetes and leakage from the pancreas are possible risks.

\medterm{Distal Radius Fractures} Breaks near the wrist end of the forearm bone, often after a fall. They may cause pain, swelling, deformity, or numbness; treatment depends on alignment, joint involvement, bone quality, and nerve or blood-vessel injury.

\medterm{Distal Renal Tubular Acidosis} A kidney disorder in which the kidneys cannot remove enough acid into the urine. Acid builds up in the blood and may cause low potassium, kidney stones, weak bones, muscle weakness, or poor growth in children.

\medterm{Distal Splenorenal Shunt} Surgery connecting the splenic vein to the left kidney vein to lower pressure in veins around the esophagus. It may reduce bleeding from varices while preserving some blood flow through the liver.

\medterm{Distal Tubule} A small kidney tube that adjusts the amount of salt, calcium, and water left in forming urine. Its activity helps regulate blood pressure and body chemistry and can be affected by some diuretic medicines.

\medterm{Distamycins} A group of antibiotic-like compounds that bind to DNA and have been studied for antimicrobial or anticancer effects. Their clinical use is limited and depends on the specific compound and evidence.

\medterm{Distant} Far from a reference point or original site. In cancer, distant disease means that cancer has spread to remote organs or lymph nodes rather than remaining near the primary tumor.


\medterm{Distended} Abnormally stretched, swollen, or enlarged by gas, fluid, blood, tissue, or pressure. The term may describe the abdomen, bladder, stomach, bowel, veins, or another body structure.

\medterm{Distension} Enlargement or stretching of a body part or hollow organ, often from gas, fluid, urine, stool, blood, pressure, or blockage. Abdominal distension with severe pain, vomiting, fever, or inability to pass stool or gas needs prompt assessment.

\medterm{Distention} Enlargement or stretching of an organ, cavity, or body part, such as abdominal distention from gas, fluid, swelling, or blockage.

\medterm{Distichia} Growth of an extra row of eyelashes, usually from openings of glands along the eyelid. Misguided lashes may rub the eye and cause irritation, tearing, corneal scratches, or ulcers.


\medterm{Distortion} An abnormal change in the shape, structure, appearance, or representation of tissue or an image. New breast, facial, or visual distortion may indicate disease and should be assessed with appropriate examination or testing.


\medterm{Distraction Osteogenesis} A surgical method for growing new bone by slowly separating two bone sections after a controlled cut. A device gradually widens the gap to lengthen or reshape the jaw, limbs, or other bones.

\medterm{Distress} Significant physical or emotional suffering, discomfort, or strain related to illness, injury, treatment, or difficult circumstances. Severe distress may affect sleep, functioning, decisions, or safety and may need support.

\medterm{Distributive Shock} Life-threatening shock caused by widespread blood-vessel dilation or abnormal blood-flow distribution. Effective circulation falls despite variable heart output; causes include sepsis, anaphylaxis, nerve injury, and some poisons.

\medterm{District Nurse} A community nurse who provides care in homes, clinics, or other local settings. Duties may include wound care, chronic-disease support, health education, medicine monitoring, prevention, and coordination with other services.

\medterm{Disulfiram-Like Reaction} An unpleasant reaction after alcohol is taken with disulfiram or another interacting substance. Flushing, headache, nausea, vomiting, palpitations, chest discomfort, or low blood pressure may occur.

\medterm{Disulfoton} A highly toxic organophosphate insecticide that inhibits acetylcholinesterase. Exposure can cause sweating, excess saliva, vomiting, diarrhea, muscle weakness, breathing failure, seizures, or coma.


\medterm{Dithiolethione} A sulfur-containing compound studied for activating cellular defense and detoxification pathways. It may influence inflammation and protection from chemical injury, but clinical uses depend on evidence for the specific compound.

\medterm{Ditiocarb} A term for dithiocarbamate compounds, some of which are used as pesticides or metal-binding chemicals. Toxic effects vary by compound and may involve the nervous system, liver, kidneys, or skin.

\medterm{Diuresis} Increased production of urine. It may result from excess fluid, high blood glucose, diabetes insipidus, medicines called diuretics, or other causes; excessive loss can lead to dehydration and electrolyte problems.

\medterm{Diuretic Drug} A medicine that makes the kidneys remove more sodium and water in urine. It is used for high blood pressure, heart failure, and swelling but can cause dehydration, low blood pressure, electrolyte changes, or kidney problems.

\medterm{Diurnal} Relating to daytime or occurring during a daily 24-hour cycle. The term may describe symptoms, hormone levels, body functions, or patterns of activity and sleep.

\medterm{Diverticula (Diverticulosis)} Alternate index label; see \textit{Diverticulosis}.

\medterm{Diverticular Disease} A condition involving small pouches in the wall of the colon. Diverticulosis means the pouches are present, while diverticular disease generally refers to symptoms or complications such as abdominal pain, altered bowel habits, bleeding, inflammation, infection, blockage, or perforation. Diverticulitis is inflammation or infection of one or more pouches and commonly causes lower-abdominal pain and fever.

\medterm{Diverticulitis} Inflammation of one or more small pouches in the wall of the colon. It commonly causes lower abdominal pain, fever, and changes in bowel habits and may lead to an abscess, perforation, or intestinal blockage.



\medterm{Diverticulosis} The presence of small pouches in the wall of the colon. It often causes no symptoms but can be associated with abdominal symptoms, bleeding, inflammation, or other complications.

\synonyms
Colonic diverticula

\medterm{Diverticulosis And Diverticulitis} Diverticulosis means small pouches have formed in the colon wall. Diverticulitis occurs when one or more pouches become inflamed or infected, causing abdominal pain, fever, and bowel changes.


\medterm{Diverticulum} A pouch-like projection from the wall of a hollow organ, especially the intestine. It may cause no symptoms or become inflamed, infected, blocked, bleed, or rarely perforate.

\medterm{Divisum Pancreas (Pancreas Divisum)} Alternate index label for Pancreas Divisum.

\medterm{Dix-Hallpike Maneuver} A bedside test for a common inner-ear cause of brief positional vertigo. The clinician moves the head and body into a specific position and watches for dizziness and characteristic eye movements.

\medterm{Dizygotic Twin} One of a pair of fraternal twins formed when two separate eggs are fertilized by two separate sperm. The twins share about half their genetic variation on average, like ordinary siblings.

\medterm{Dizygotic Twins} Twins formed when two separate eggs are fertilized by two separate sperm. They are genetically like ordinary siblings and may have separate or shared-appearing placentas; the pregnancy still requires appropriate monitoring.

\medterm{Dizziness} A sensation of spinning, lightheadedness, imbalance, faintness, or being unsteady. Causes include inner-ear, heart, blood-pressure, vision, medicine, blood, metabolic, or neurologic problems.

\medterm{Dizziness (Dizzy)} Alternate index label for Dizzy.


\medterm{Dizziness: When To See A Doctor?} Seek urgent care for sudden dizziness with weakness, numbness, trouble speaking, severe headache, chest pain, fainting, inability to walk, persistent vomiting, or new hearing loss. Recurrent or unexplained dizziness also needs medical evaluation.

\medterm{Déjà Vu} A feeling that a current situation has happened before, even though it is probably new. Occasional episodes are common; frequent episodes with impaired awareness or unusual movements may occur with seizures.

\medterm{DKD} Alternate index label; see \textit{Diabetic Nephropathy}.

\medterm{DLI} Donor lymphocyte infusion, in which immune cells from a stem-cell donor are given after transplantation to strengthen an immune attack against returning cancer. It can cause graft-versus-host disease and other complications.

\medterm{DM} An abbreviation most commonly meaning diabetes mellitus, a group of disorders with persistently high blood glucose. Because DM can have other meanings, the surrounding clinical context is important.

\medterm{DM2} An abbreviation for type 2 diabetes mellitus, in which the body resists insulin and may gradually produce too little of it. It can cause long-term damage to the eyes, kidneys, nerves, heart, and blood vessels.

\medterm{DMFT Index} A dental measure counting permanent teeth that are decayed, missing because of decay, or filled. It estimates a person's lifetime experience of tooth decay but does not show whether a current problem is active.

\medterm{DMG} Alternate index label; see \textit{Diffuse midline glioma}.

\medterm{DMSA Scan} A nuclear-imaging test using a small amount of labeled dimercaptosuccinic acid to show kidney tissue. It can identify differences in kidney function, scars, or injury, especially in children with repeated urinary infections.

\medterm{DNA} Deoxyribonucleic acid, the molecule that stores hereditary information in humans and most organisms. Its sequence helps control cell function and protein production; changes may be inherited or acquired during life.

\medterm{DNA Adduct} A chemical group attached to DNA after exposure to a reactive substance. If not repaired, it can alter DNA copying, cause mutations, and potentially contribute to cancer.

\medterm{DNA Alkylation} Attachment of an alkyl group to DNA. This can change base pairing, block copying or repair, cause mutations, or kill a cell; some anticancer medicines work partly this way.

\medterm{DNA Amplification} Production of many copies of a selected DNA sequence, naturally in cells or in the laboratory by methods such as polymerase chain reaction. It helps detect, identify, or study genes and microorganisms.

\medterm{DNA Binding} Attachment of a protein or other molecule to a DNA sequence. This can control gene activity, DNA copying, repair, chromosome organization, or the action of some medicines.

\medterm{DNA Damage} A chemical or structural change in DNA that can interfere with gene function or copying. Cells usually repair such damage, but persistent or incorrectly repaired damage can cause mutations, cell death, or cancer.
\medterm{DNA Demethylation} Removal or loss of small chemical groups called methyl groups from DNA. It can change how accessible a gene is and alter gene activity during development, cell reprogramming, or disease.

\medterm{DNA Double Helix} The twisted ladder-like structure of DNA, formed by two strands running in opposite directions. Sugar and phosphate form the outside rails, while paired bases join the strands inside.



\medterm{DNA Fragmentation} Breaking DNA strands into smaller pieces. It occurs normally during programmed cell death and can also result from injury, radiation, chemicals, disease, or laboratory processing.

\medterm{DNA Gyrase} A bacterial enzyme that twists and untwists DNA so it can be copied and separated. Fluoroquinolone antibiotics block this enzyme, stopping bacterial growth and survival.

\medterm{DNA Integration} Insertion of foreign or copied DNA into a cell’s genome. It occurs in some viral infections, evolution, and gene therapy and can alter nearby gene activity if the insertion affects a control region.

\medterm{DNA Ligases} Enzymes that join breaks in DNA strands during copying, repair, and genetic recombination. They help preserve the continuous structure of chromosomes.

\medterm{DNA Ligation} Joining DNA fragments by forming a chemical bond between them. Cells use it during DNA copying and repair, and laboratories use it in genetic testing and cloning.

\medterm{DNA Methylation} Addition of small chemical groups called methyl groups to DNA, usually at cytosine bases. It can change how strongly a gene is expressed without changing the DNA sequence.

\medterm{DNA Packaging} Arrangement of long DNA molecules around histone proteins and into compact structures called chromosomes. Packaging allows DNA to fit inside the nucleus and helps control access to genes.

\medterm{DNA Primase} An enzyme that makes short RNA starting pieces so other enzymes can begin copying DNA. It is required during cell division and DNA repair.

\medterm{DNA Primers} Short nucleic-acid sequences that provide a starting point for copying DNA. Cells make them during replication, and laboratories use designed primers to amplify selected genetic sequences.

\medterm{DNA Probe} A labeled single-stranded DNA sequence designed to bind a matching target. It can help detect a gene, mutation, chromosome region, or microorganism in a laboratory sample.

\medterm{DNA Profiling} Comparison of variable DNA markers from biological samples to assess identity or biological relationships. Results are statistical and require validated methods, suitable controls, and reliable sample handling.

\medterm{DNA Repair} Cellular processes that detect and correct damage or copying errors in DNA. Repair pathways address small base changes, bulky chemical damage, mismatched bases, and breaks in one or both DNA strands.

\medterm{DNA Repair Mechanisms} Cellular systems that correct DNA damage caused by copying errors, chemicals, radiation, or reactive oxygen. Failure of some pathways increases mutation risk and can contribute to inherited cancer syndromes.

\medterm{DNA Replication} The process by which DNA copies itself before cell division. Each new DNA molecule contains one original strand and one newly made strand, allowing genetic information to pass to daughter cells.

\medterm{DNA Sequencing} Determination of the order of DNA building blocks in a sample. It can examine a single gene, selected genes, all coding regions, or the whole genome to identify genetic changes.

\medterm{DNA Virus} A virus whose genetic material is DNA. DNA viruses vary widely in transmission, tissues affected, and disease; some cause persistent infection, while others cause short-term illness.

\medterm{D.O.} Doctor of Osteopathic Medicine, a physician degree whose graduates are licensed to practice medicine and surgery like MD graduates in the United States.

\medterm{Do Not Resuscitate} A medical order stating that cardiopulmonary resuscitation should not be attempted if a person’s breathing or heartbeat stops, according to the person’s informed wishes or legally authorized decision-maker.

\medterm{DOACs} Direct oral anticoagulants are oral blood thinners that directly block thrombin or factor Xa. They prevent and treat blood clots but can cause bleeding and require attention to kidney function, interactions, and planned procedures.

\medterm{DOB} An abbreviation for date of birth, the day, month, and year a person was born. It helps confirm identity and calculate age for records, screening, medicine doses, and treatment decisions.

\medterm{Dobutamine} An intravenous medicine that increases heart contraction and, sometimes, heart rate. It is used short term for selected cases of severe low-output heart failure or shock and requires close monitoring.

\medterm{Dobutamine Stress Echocardiography} An ultrasound heart test using dobutamine to increase heart workload when a person cannot exercise adequately. New areas of poor wall movement may suggest reduced blood flow, but results require clinical interpretation.

\medterm{Doc} Informal shorthand for doctor, a person qualified to practice medicine or another health profession. The exact role depends on the person’s training and professional license.

\medterm{Docetaxel} An anticancer medicine that interferes with microtubules, structures needed for cell division. It is used for several cancers and can cause low blood counts, infection, hair loss, nerve problems, or allergic reactions.

\medterm{Doconexent} A pharmaceutical name for docosahexaenoic acid, an omega-3 fatty acid found in oily fish, algae, and cell membranes. Its effects depend on the product, dose, and reason for use.

\medterm{Docosadienoic Acid} A 22-carbon fatty acid with two double bonds, present in small amounts in body and dietary lipids. It is studied mainly for possible effects on cell membranes and metabolism.

\medterm{Docosahexaenoic Acid} An omega-3 fatty acid that is an important structural part of brain and retinal cell membranes. It comes mainly from oily fish, algae, or supplements and is especially important during brain development.

\medterm{Docosapentaenoic Acid} A 22-carbon polyunsaturated fatty acid with five double bonds. It occurs as omega-3 and omega-6 forms in foods and tissues; its health effects depend on the form and amount.

\medterm{Docosatrienoic Acid} A 22-carbon fatty acid with three double bonds, found in small amounts in body and dietary lipids. It is studied mainly for its possible roles in cell membranes and metabolism.

\medterm{Doctor} A person who has completed appropriate professional training and is licensed or authorized to practice medicine. The title may also be used by other professionals with doctoral degrees, depending on local rules.



\medterm{Docusate} A stool-softening medicine that helps water mix with stool. It may be used when avoiding straining is important, but evidence for treating ordinary constipation is limited; diarrhea and abdominal cramps can occur.


\medterm{Dog Allergy} An immune reaction to proteins from a dog’s skin flakes, saliva, urine, or other secretions. It can cause sneezing, blocked or itchy nose, watery eyes, hives, cough, or asthma symptoms.

\medterm{Dog Bite} An injury caused by a dog’s teeth. Wash the wound promptly and seek assessment for tissue damage, infection, tetanus protection, and rabies risk, especially for deep bites or bites to the face or hand.

\synonyms
Dog bite injury

\medterm{Dog Bite Treatment} Care includes thorough wound cleaning, assessment of tendon, nerve, or bone injury, tetanus review, rabies-risk evaluation, and antibiotics when indicated. Deep bites, severe bleeding, infection, or bites to the hand or face need prompt care.

\synonyms
Bite Dog (Dog Bite Treatment)

\medterm{Dolasetron} A medicine that blocks serotonin 5-HT3 receptors and helps prevent nausea and vomiting in selected settings. It can cause headache, constipation, or dangerous heart-rhythm changes.

\medterm{Dolasetron Mesylate} A salt form of dolasetron, a medicine that blocks serotonin 5-HT3 receptors to prevent nausea and vomiting in selected settings. It can cause constipation, headache, or dangerous heart-rhythm changes.

\medterm{Dolichocephalic} Having a relatively long, narrow skull. It may be a normal family or positional variation, but marked or worsening skull narrowing in an infant may indicate early fusion of a skull suture and needs assessment.

\medterm{Dolichol} A long-chain fatty alcohol used inside cells to assemble sugar chains attached to proteins. This process is important for protein folding, transport, and function.

\medterm{Dolor} The Latin word for pain, one of the classic signs of inflammation. Pain may result from irritated nerves, swelling, injury, infection, or another cause and is interpreted with other findings.

\medterm{Domain} A defined region or area. In biology, a protein domain is a part with a particular structure or function; in genetics, a domain may be a genomic or regulatory region.

\medterm{Domestic Violence} A pattern of physical, sexual, emotional, psychological, or financial abuse or coercive control by an intimate partner or household member. It can cause injury, fear, isolation, and serious health effects; confidential support is available.

\medterm{Dominant} Describing a genetic variant whose effect can appear when only one copy is present. An affected parent with an autosomal dominant condition may pass the variant to each child with a one-in-two chance.

\medterm{Domperidone} A medicine that blocks dopamine receptors outside the brain and is used in some countries for nausea or slow stomach movement. It can prolong the heart rhythm and cause dangerous arrhythmias.

\medterm{Donath-Landsteiner Test} A blood test that detects an antibody associated with paroxysmal cold hemoglobinuria. The antibody binds red cells in the cold and destroys them when the blood warms.

\medterm{Donepezil} A medicine that increases acetylcholine activity in the brain and may temporarily improve symptoms of Alzheimer disease. It does not cure the disease and can cause nausea, diarrhea, slow heartbeat, fainting, or sleep problems.

\medterm{Donnan Equilibrium} The distribution of charged particles across a partly permeable membrane when large charged particles cannot cross. It helps explain cell volume, membrane electrical charge, and movement of salts and water.

\medterm{Donohue Syndrome} A rare inherited disorder in which the body cannot respond normally to insulin. It causes severe growth and feeding problems, very little body fat, unusual facial features, and serious metabolic complications beginning in infancy.

\medterm{Donor} A person who provides blood, cells, tissue, an organ, sperm, eggs, or another biological material for another person or for research. Screening, informed consent, compatibility, and traceability support safety.

\medterm{Donor Insemination} A fertility treatment in which sperm from a screened donor is placed in the reproductive tract to attempt pregnancy. The donor may be known or anonymous, and counseling and legal requirements vary by location.

\medterm{Donor Suitability Assessment} Evaluation of a potential donor’s medical history, examination, laboratory tests, infection risks, organ function, and compatibility to determine whether donation is safe and suitable for transplantation.

\medterm{Donor Testing} Laboratory and clinical evaluation of a potential donor for infections, cancer, organ function, blood type, and other factors that affect donation safety and recipient compatibility.
\medterm{Do-Not-Resuscitate Order} A medical order instructing clinicians not to perform cardiopulmonary resuscitation if a person’s breathing or heartbeat stops. It does not automatically refuse other treatments or comfort care.

\synonyms
DNR order

\medterm{Donovanosis} A sexually transmitted infection caused by Klebsiella granulomatis. It produces slowly enlarging, usually painless genital ulcers that may bleed easily and requires medical diagnosis and antibiotic treatment.

\synonyms
granuloma inguinale.

\medterm{Donovanosis (Granuloma Inguinale)} A sexually transmitted bacterial infection causing slowly progressive, usually painless genital ulcers that bleed easily. Diagnosis and antibiotic treatment require clinical evaluation, and sexual partners may need assessment.

\medterm{Dopa} A chemical precursor that the body can convert into dopamine. Levodopa, a form that reaches the brain, is used with another medicine to improve movement problems in Parkinson disease.

\medterm{Dopamine} A chemical messenger used by nerve cells and some hormone-producing cells. It helps control movement, motivation, reward, attention, learning, blood-vessel tone, and release of certain hormones.

\medterm{Dopamine Agonist} A medicine that stimulates dopamine receptors and mimics some effects of dopamine. It is used for selected movement, hormone, or other disorders and may cause nausea, sleepiness, low blood pressure, hallucinations, or impulse-control problems.
\medterm{Dopamine Agonists} Medicines that stimulate dopamine receptors and are used for selected Parkinson, restless-legs, or hormone disorders. Possible effects include nausea, sleepiness, low blood pressure, hallucinations, and impulse-control problems.

\medterm{Dopamine Antagonist} A medicine or compound that blocks dopamine receptors. It is used in selected psychiatric, nausea, or digestive treatments but can cause movement problems, drowsiness, hormone changes, or other adverse effects.

\medterm{Dopamine Receptor} A protein on or inside a cell that responds to dopamine. Different receptor types affect movement, mood, hormone release, blood vessels, and other functions and are targets of some medicines.

\medterm{Dopamine Transporter} A protein that clears dopamine from the space between nerve cells by moving it back into the releasing nerve cell. Stimulants and some imaging tracers affect or measure its activity.

\medterm{Dopaminergic Neuron} A nerve cell that makes and releases dopamine. These cells help control movement, motivation, reward, attention, and hormone regulation; loss of some contributes to Parkinson disease.

\medterm{Dopa-Responsive Dystonia} An inherited movement disorder causing sustained muscle contractions and abnormal postures, sometimes with Parkinson-like symptoms. It often improves markedly with a small dose of levodopa.

\medterm{Doppler Echocardiography} A heart ultrasound technique that measures the direction and speed of blood flow through the heart and major vessels. It helps assess valves, chambers, pressure, and abnormal flow.

\medterm{Doppler Ultrasound} An ultrasound technique that measures moving blood to assess its direction and speed. It can identify narrowed or blocked vessels, abnormal flow, blood clots, and circulation problems without ionizing radiation.

\medterm{Doppler Ultrasound Exam Of An Arm Or Leg} An ultrasound examination measuring blood flow through the arteries and veins of an arm or leg. It can help identify narrowing, blockage, a clot, aneurysm, or other circulation problem without ionizing radiation.

\medterm{Doppler Velocimetry} Ultrasound measurement of blood-flow speed and resistance in fetal or placental vessels. It helps assess circulation and fetal well-being when poor growth, placental problems, or other pregnancy risks are suspected.

\medterm{Dor Procedure} A heart operation that removes or excludes a scarred, bulging part of the left ventricle after a large heart attack. It reshapes the chamber in selected people with severe ischemic heart-muscle disease.

\synonyms
Endoventricular Circular Patch Plasty (EVCPP); Surgical Ventricular Restoration

\medterm{Doravirine} An HIV medicine that blocks reverse transcriptase, an enzyme HIV needs to copy its genetic material. It is used with other medicines for HIV-1 infection; interactions and resistance affect its choice.

\medterm{Dorea} A group of bacteria commonly found in the human intestine and able to grow without oxygen. Finding it usually reflects normal gut organisms unless it comes from a normally sterile site with compatible illness.

\medterm{Doripenem} An intravenous antibiotic that damages bacterial cell walls and is used for selected serious infections. It can cause allergy, diarrhea, seizures, or resistance and should be chosen according to the organism and patient.

\medterm{Dormant} Inactive or not multiplying. The term may describe a latent infection, quiet tumor cell, inactive disease process, or other biological state that can become active under certain conditions.

\medterm{Dorsal} Relating to the back or upper surface of a body part. In human anatomy it usually means toward the back, although in the hand or foot it refers to the opposite surface from the palm or sole.
\medterm{Dorsal Blocking Splint} A hand splint that limits wrist and finger extension to protect healing tendons, commonly after flexor-tendon repair. Position and duration are set by the treating surgical and therapy team.

\medterm{Dorsal Funiculus} A back column of the spinal cord containing nerve fibers that carry vibration, fine touch, and awareness of body position to the brain.

\medterm{Dorsal Kyphosis} Excessive forward curvature of the upper spine, causing a rounded upper back. It may result from posture, spinal wear, osteoporosis, or vertebral fractures and can cause pain, height loss, or breathing difficulty.
\medterm{Dorsal Root Ganglion} A swelling on the back root of a spinal nerve containing the cell bodies of sensory nerve cells. It relays information about touch, pain, temperature, and body position to the spinal cord.

\medterm{Dorsalis Pedis Artery} An artery running across the top of the foot as a continuation of the anterior tibial artery. Its pulse helps assess blood flow to the foot.

\medterm{Dorsiflexion} Upward movement of the foot at the ankle, bringing the toes toward the shin. It helps clear the foot during walking; weakness can cause foot drop and tripping.

\medterm{Dorsolateral} Describing a position toward the back and side of a body structure. For example, a dorsolateral area is behind and to one side rather than toward the front or middle.

\medterm{Dorsum} The back or upper surface of a body part, such as the back of the hand, top of the foot, or upper surface of the tongue.

\medterm{Dorzolamide} An eye medicine that lowers pressure inside the eye by reducing fluid production. It is used for glaucoma or raised eye pressure and may cause stinging, blurred vision, or a bitter taste.

\medterm{Dos And Donts During First Trimester Of Pregnancy} Early pregnancy care includes prenatal visits, folic acid, avoiding alcohol, tobacco, and non-prescribed drugs, reviewing medicines, following food-safety advice, and staying appropriately active. Severe pain, heavy bleeding, fainting, fever, or persistent vomiting needs urgent care.

\medterm{Dosage} The amount, frequency, route, and duration of a medicine given to achieve benefit while limiting harm. It may need adjustment for age, weight, kidney or liver function, interactions, and the condition treated.

\medterm{Dosage Forms} The physical form in which a medicine is supplied, such as a tablet, capsule, liquid, cream, inhaler, injection, patch, suppository, or implant. The form affects how and how quickly it is given.

\medterm{Dose} The amount of a medicine, radiation, or other substance given or received at one time. The appropriate dose depends on the substance, purpose, age, weight, organ function, interactions, and individual risk.

\medterm{Dose Rate} The amount of a medicine, radiation, or other exposure delivered during a specified time, such as milligrams per hour or grays per minute. It affects both benefit and toxicity.

\medterm{Dose-Limiting Toxicity} A treatment-related adverse effect severe enough to prevent increasing the dose or require reducing it. The term is used especially in early cancer trials to estimate a tolerable dose.

\medterm{Dose-Response Curve} A graph showing how the size of a biological effect changes as the dose of a medicine or substance increases. It helps describe potency, maximum effect, and the dose range that produces benefit or harm.

\medterm{Dose-Response Relationship} The relationship between the amount of a medicine or substance and the size of its effect. Increasing the dose may increase benefit, toxicity, or both, but the pattern depends on the substance and person.

\medterm{Dose Titration} Gradual adjustment of a medicine dose according to its effects, treatment goals, blood tests, or side effects. The dose may be increased, decreased, or kept the same as the person’s response is assessed.

\medterm{Drug-Food Interaction} A change in a medicine’s effect, absorption, or safety caused by food or drink. Food may increase or decrease absorption, while some foods can alter medicine metabolism or raise the risk of adverse effects.

\medterm{Pharmacodynamic Drug Interaction} An interaction in which two medicines have additive, opposing, or otherwise changed effects on the same body system, even when their blood concentrations are not altered.

\medterm{Pharmacokinetic Drug Interaction} An interaction in which one medicine changes another medicine’s absorption, distribution, metabolism, or elimination, altering its concentration or duration of action.

\medterm{Dosimeter} An instrument worn or placed near a person or object to measure or record exposure to ionizing radiation. It helps monitor occupational or treatment-related radiation.

\medterm{Dosimetry} Measurement and calculation of radiation dose received by a person, tissue, or object. In treatment, it helps plan and verify that the intended dose reaches the target while limiting exposure elsewhere.

\medterm{DOTATATE} A substance that binds somatostatin receptors found on many neuroendocrine tumors. When linked to a radioactive tracer, it is used in PET imaging to locate and stage tumors and assess possible targeted treatment.

\medterm{DOTS} A tuberculosis-control strategy involving reliable diagnosis, standardized combination treatment, medicine-supply support, monitoring, and sometimes directly observed doses. It aims to cure infection and limit drug resistance and transmission.

\medterm{Double Aortic Arch} A birth defect in which two aortic arches form a ring around the windpipe and esophagus. It may cause noisy breathing, cough, feeding difficulty, vomiting, or trouble swallowing.

\medterm{Double Blind} A study design in which participants and the researchers assessing outcomes do not know which treatment each participant receives during the blinded phase. It reduces expectation and assessment bias.

\synonyms
Double-masked

\medterm{Double Blind Trial} A clinical trial in which participants and the investigators measuring outcomes do not know the assigned treatment during the blinded phase. This helps reduce expectations from influencing results.

\medterm{Double Inlet Left Ventricle} A rare birth defect in which both upper heart chambers connect mainly to one lower chamber. It produces single-ventricle circulation and mixing of oxygen-rich and oxygen-poor blood and usually needs staged specialist treatment.

\medterm{Double Outlet Right Ventricle} A birth defect in which both major arteries arise mainly from the right lower heart chamber. An opening between the lower chambers allows blood from the left chamber to reach the arteries; anatomy and symptoms vary.

\medterm{Double Uterus} A birth difference in which the uterus develops as two separate cavities or two uteruses. Some people have no symptoms; others may have menstrual problems, infertility, miscarriage, or pregnancy complications.
\medterm{Double Vision} Seeing two images of one object. Causes include eye-surface or lens disease, eye misalignment, muscle or nerve problems, medicines, migraine, or brain disease; sudden onset with neurologic symptoms or eye pain is an emergency.

\medterm{Double-Balloon Enteroscopy} An endoscopic procedure using two balloons to move a flexible camera through much of the small intestine. It can locate bleeding or lesions, take biopsies, and perform selected treatments.

\medterm{Double-Blind Method} A research method in which treatment assignments are hidden from participants and relevant study personnel during the blinded phase to reduce expectation and assessment bias.

\medterm{Double-Blind Study} A clinical study in which participants and the investigators assessing outcomes do not know the assigned treatment group during the blinded phase.

\medterm{Double-Contrast Barium Enema} An X-ray examination in which barium coats the inside of the colon and air expands it. It can show polyps, masses, narrowing, or changes in the bowel lining.

\synonyms
Air-contrast barium enema

\medterm{Double-Outlet Right Ventricle} A birth defect in which both major arteries arise mainly from the right lower heart chamber. An opening between the lower chambers allows blood to reach the arteries; symptoms and treatment depend on the exact anatomy.
\medterm{Douche} A stream of water or solution directed into a body cavity, most often the vagina. Vaginal douching is not needed for hygiene and can irritate tissue or disturb normal bacteria and increase infection risk.

\medterm{Douglas, Pouch Of} A space in the pelvis between the uterus and rectum. In a person without a uterus, the comparable space lies between the bladder and rectum; fluid or blood can collect there.


\medterm{Dowager’S Hump} An exaggerated forward curve of the upper spine that creates a rounded back. It may result from osteoporosis and collapsed vertebrae; new pain, height loss, breathing difficulty, or worsening curvature needs assessment.

\synonyms
Kyphosis

\medterm{Down Syndrome} A genetic condition caused by an extra copy of chromosome 21. It may affect learning, muscle tone, growth, and appearance and is associated with heart, hearing, vision, thyroid, sleep, and digestive problems.

\synonyms
Trisomy 21 Syndrome
Down Syndrome Overview
Trisomy 21 (Down Syndrome Overview)
Down's Syndrome

\medterm{Downregulation} A decrease in the number, sensitivity, or activity of receptors, genes, or signaling pathways. It may occur after prolonged stimulation, medicine exposure, inflammation, or another regulatory change.

\medterm{Doxapram} A medicine that stimulates breathing and is used rarely in selected situations. It can cause agitation, high blood pressure, fast or abnormal heart rhythms, seizures, or other serious effects.

\medterm{Doxazosin} An alpha-1 blocker that relaxes blood vessels and prostate or bladder-neck muscle. It is used for high blood pressure or urinary symptoms from an enlarged prostate and can cause dizziness or fainting.

\medterm{Doxazosin Mesylate} A salt form of doxazosin, an alpha-1 blocker used for high blood pressure and urinary symptoms from an enlarged prostate. It can cause dizziness, low blood pressure, fainting, or falls, especially after the first dose.

\medterm{Doxepin} A tricyclic antidepressant also used in low doses for insomnia or skin itching. It can cause drowsiness, dry mouth, constipation, blurred vision, abnormal heart rhythms, and dangerous toxicity in overdose.

\medterm{Doxepin Overdose} Taking too much doxepin can cause severe drowsiness, confusion, seizures, low blood pressure, coma, and life-threatening heart-rhythm or conduction problems, especially with alcohol or other medicines.

\medterm{Doxercalciferol} A vitamin D-like medicine used in selected people with chronic kidney disease to reduce excess parathyroid hormone. Calcium and phosphate must be monitored because levels can become too high.

\medterm{Doxorubicin} An anthracycline cancer medicine that damages DNA and blocks cell division. It treats many cancers but can cause low blood counts, infection, nausea, hair loss, tissue injury after leakage, and dose-related heart damage.

\medterm{Doxorubicin Hydrochloride} A salt form of doxorubicin, an anthracycline cancer medicine used for many cancers. Important risks include low blood counts, infection, tissue injury after leakage, and dose-related heart damage.

\medterm{Doxycycline} A tetracycline antibiotic that blocks bacterial protein production. It treats many bacterial and tick-borne infections and can cause nausea, sun sensitivity, esophageal irritation, diarrhea, or tooth and bone effects in developing children.

\medterm{Doxycycline Hyclate} A salt form of doxycycline, a tetracycline antibiotic that blocks bacterial protein production. It is used for susceptible infections; nausea, sun sensitivity, esophageal irritation, diarrhea, and allergic reactions may occur.

\medterm{Doxylamine} A sedating antihistamine used in some sleep aids and pregnancy-related nausea treatments. It can cause drowsiness, dry mouth, blurred vision, constipation, urinary retention, confusion, or impaired driving.

\medterm{DPP-4 Inhibitor} An oral diabetes medicine that prolongs the action of hormones released after meals, increasing glucose-dependent insulin release and reducing glucagon. It is used mainly in type 2 diabetes and usually has a low risk of low glucose.

\medterm{DPP-4 Inhibitors} Oral medicines for type 2 diabetes that prolong hormones released after meals, helping the pancreas release insulin and reducing glucagon. They usually do not cause weight gain or low glucose when used alone but have modest glucose-lowering effects.

\medterm{DPT Immunization} Vaccination that protects against diphtheria, pertussis (whooping cough), and tetanus. Current childhood and booster vaccines use age-appropriate formulations and are given as a series according to the recommended schedule.
\medterm{DR Syndrome} An abbreviation that may refer to diabetic retinopathy syndrome or another context-specific disorder; the full clinical term should be specified.

\medterm{Dracunculiasis} Guinea-worm disease caused by a parasite acquired from unsafe drinking water containing infected tiny water fleas. About a year later, a painful blister usually forms on a leg as the adult worm emerges through the skin.

\medterm{Dracunculosis} Infection with the Guinea-worm parasite after drinking unsafe water containing infected water fleas. It causes a painful blister, usually on a leg, followed by emergence of a long adult worm through the skin.

\medterm{Dracunculus Medinensis} The parasitic worm that causes Guinea-worm disease. People become infected by drinking water containing tiny infected water fleas; roughly a year later, an adult female worm emerges through the skin.
\medterm{Drain} A tube or other device used to remove blood, pus, fluid, or air from a wound, body cavity, or organ. Drainage may prevent pressure, infection, or impaired healing after injury or surgery.

\medterm{Drain Cleaner Poisoning} Exposure to a caustic drain-cleaning product that can cause severe burns of the skin, eyes, airway, or gastrointestinal tract.

\medterm{Drain Opener Poisoning} Exposure to a caustic drain-cleaning product that can cause severe burns of the skin, eyes, airway, or gastrointestinal tract.

\medterm{Drain (Surgical)} A tube placed during or after surgery to remove blood, pus, or other fluid from an operative site. Possible complications include infection, blockage, leakage, bleeding, and movement of the tube.

\synonyms
Surgical drain

\medterm{Drainage} Removal of blood, pus, fluid, or air from a body site, wound, or abnormal collection. It may occur naturally or use a needle, catheter, tube, or surgical opening to relieve pressure and aid healing.

\medterm{Drainage Procedures} Treatments that remove pus, blood, or other collected fluid, sometimes with ultrasound or CT guidance. They may use a needle, catheter, tube, or surgery to treat abscesses and other fluid collections.

\medterm{Drainpipe Cleaners} Strong chemical products used to clear blocked pipes, commonly containing caustic alkalis or acids. Swallowing, inhaling, or splashing these products can cause severe burns to the mouth, throat, eyes, skin, or lungs and requires urgent poison-control or medical advice.


\medterm{Dream} A sequence of images, thoughts, sensations, or emotions experienced during sleep. Dreams are often most vivid during rapid-eye-movement sleep and may be influenced by stress, medicines, illness, or sleep disorders.

\medterm{Dreams} Experiences involving images, thoughts, sensations, or emotions during sleep, often vivid during rapid-eye-movement sleep. Dreams are normal, but distressing, violent, or disruptive dreams may occur with stress, medicines, or sleep disorders.

\medterm{Dreams (Geriatric)} Dreams in older adults may become more vivid or be affected by medicines, sleep disorders, depression, or changes in the sleep cycle. Repeated frightening dreams, acting them out, confusion, or sudden change requires clinical assessment.

\synonyms
Dreams in older adults

\medterm{Drepanocytosis} Sickle-cell disease or a related condition in which an abnormal form of hemoglobin makes red blood cells become rigid and sickle-shaped. This can cause anemia, painful blockages of blood flow, and damage to organs.

\medterm{DRESS Syndrome} A serious delayed allergic reaction to a medicine, usually beginning two to eight weeks after exposure. It may cause fever, widespread rash, facial swelling, abnormal blood counts, hepatitis, kidney injury, or other organ inflammation.

\medterm{Dressing} Material placed over a wound to protect it from contamination, absorb drainage, control minor bleeding, maintain moisture, or deliver local treatment. It should be changed as directed and monitored for infection.

\medterm{Dressler Syndrome} An immune-related inflammation of the sac around the heart that occurs weeks after a heart attack, heart surgery, or other heart injury. It can cause fever, sharp chest pain, and fluid around the heart or lungs.

\synonyms
Post-Cardiac Injury Syndrome
Postmyocardial Infarction Syndrome

\medterm{Dressler's Syndrome} Alternate spelling; see \textit{Dressler Syndrome}.

\medterm{Dressler’S Syndrome} Alternate spelling; see \textit{Dressler Syndrome}.

\medterm{DRI} Dietary reference values used to estimate the amounts of vitamins, minerals, protein, and other nutrients needed by healthy people. They help assess diets but are not individualized treatment targets for every illness.

\synonyms
Dietary Reference Intakes

\medterm{Dried Specimen} A biological sample collected and dried on a prepared card or other surface for transport and laboratory testing. Dried blood spots are commonly used for screening, but correct collection and handling are essential.


\medterm{Drift (Sit-To-Stand)} A tendency to lose balance or move away from a stable position while standing up from a chair. It may reflect weakness, poor coordination, dizziness, or impaired balance and should be interpreted with the full examination.

\synonyms
Sit-to-Stand

\medterm{Drinking Water} Water safe for human consumption. It should be free of disease-causing organisms and harmful chemicals; boiling, filtration, chlorination, or bottled water may be needed when the supply is contaminated or uncertain.

\medterm{Drip} A drop-by-drop flow of liquid, or fluid given slowly into a vein through tubing and a catheter. The infusion rate is controlled to deliver medicines, fluids, blood products, or nutrients safely.

\medterm{Drip (Intravenous Infusion)} Delivery of fluid, blood products, nutrients, or medicine directly into a vein through a catheter. The amount and speed are prescribed and monitored for reactions, fluid overload, infection, or vein irritation.

\synonyms
Intravenous Infusion

\medterm{Driving And Older Adults} Evaluation of whether an older person can drive safely, considering vision, attention, memory, hearing, movement, health conditions, and medicines. Age alone does not determine driving ability.

\medterm{Driving Evaluation} An assessment of driving ability after illness, injury, or a health change. It may examine vision, attention, memory, reaction time, strength, and vehicle control in a clinic and on the road; adaptive equipment may help.

\medterm{Driving Rehabilitation} Assessment and training to help a person resume driving safely or use adaptive vehicle controls after illness or injury. It may address vision, attention, reaction time, strength, vehicle operation, and road safety.

\medterm{Dronabinol} A medicine related to the active substance in cannabis, used for selected cases of chemotherapy-related nausea or severe appetite loss. It may cause dizziness, sleepiness, confusion, mood changes, or impaired coordination.

\medterm{Drool} Saliva that unintentionally escapes from the mouth. It may result from excess saliva, difficulty swallowing, poor lip closure, dental problems, or a neurologic condition affecting mouth control.

\medterm{Drooling} Unintentional loss of saliva from the mouth caused by excess saliva, difficulty swallowing, poor lip closure, or reduced control of the mouth. Persistent drooling may need assessment for dental, medication, or neurologic causes.

\medterm{Drop Foot} Alternate index label; see \textit{Foot drop}.

\medterm{Drop Hand Syndrome} Weakness or paralysis of the muscles that straighten the wrist and fingers, usually from injury or pressure on the radial nerve. The wrist hangs downward and gripping or releasing objects becomes difficult.

\medterm{Droperidol} A medicine used in selected clinical settings to treat severe nausea or agitation. It can prolong the heart’s electrical recovery time and rarely cause dangerous arrhythmias, so heart-risk factors and monitoring may be important.

\medterm{Droplet Precautions} Infection-control measures for infections spread by respiratory droplets released during coughing, sneezing, or talking. They may include a medical mask, hand hygiene, patient separation, and limiting transport according to the infection.

\medterm{Dropper Bottle} A bottle with a device that releases liquid one drop at a time. The medicine’s concentration, prescribed number of drops, route, and measuring device must be followed because drop size can vary.


\medterm{Drowning} Breathing impairment caused by submersion or immersion in liquid. It may be fatal or nonfatal. Lack of oxygen can injure the brain and other organs; complications include aspiration, lung injury, and respiratory failure.

\medterm{Drowning (Dry/Wet Signs)} Alternate index label; see \textit{Drowning}. The terms “dry drowning” and “wet drowning” are outdated and should not be used to classify drowning.

\medterm{Drowsiness} An increased tendency to fall asleep or reduced alertness. It may result from too little sleep, medicines, alcohol or other substances, metabolic illness, sleep disorders, infection, or neurologic disease; sudden severe drowsiness needs urgent assessment.

\medterm{Drug} A substance that changes body function or is used to prevent, diagnose, or treat disease. Drugs include prescription medicines, nonprescription medicines, biological products, and substances used recreationally or illegally.

\medterm{Drug Addiction} Legacy, nonpreferred label; see \textit{Substance Use Disorder}.

\medterm{Drug Caution Code} A label or code that warns about an important medicine risk, such as allergy, bleeding, interaction, overdose, or a need for special monitoring. Its meaning depends on the coding system and product information.
\medterm{Drug Delivery} The way a medicine is prepared and given so its active ingredient reaches the intended body site at an effective dose. Routes include swallowing, inhalation, skin application, injection, implantation, and targeted systems.

\medterm{Drug Desensitization} A supervised process that gives gradually increasing doses of a medicine to temporarily reduce an allergic reaction when that medicine is essential and no suitable alternative exists. The tolerance usually ends when doses are interrupted.

\medterm{Drug Discovery} The process of finding and developing substances that may prevent, diagnose, or treat disease. It includes laboratory research, safety testing, dose studies, clinical trials, manufacturing, and review by a regulatory authority.

\medterm{Drug Hypersensitivity} An immune reaction to a medicine that may cause rash, itching, swelling, breathing difficulty, fever, blood abnormalities, or organ inflammation. Severe forms include anaphylaxis and serious blistering or systemic skin reactions.

\medterm{Drug Implant} A solid medicine-containing device placed under the skin or in another tissue to release medicine slowly over time. It may provide local or whole-body treatment and can require removal when treatment ends.

\medterm{Drug Repurposing} Studying or using a medicine developed for one condition to treat a different condition. Earlier safety information may help, but the new use still needs evidence about effectiveness, dose, interactions, and risks.

\medterm{Drug Safety} The prevention, detection, and management of harmful effects, interactions, errors, and misuse involving medicines. It includes choosing the right medicine and dose, checking allergies and interactions, following instructions, monitoring effects, and reporting suspected problems.

\medterm{Drug Screening} Testing a sample such as urine, blood, saliva, or hair for selected drugs or their breakdown products. Initial results can be falsely positive or negative and may require confirmation with a more specific test.

\medterm{Drug Testing} Laboratory analysis of blood, urine, saliva, hair, or another sample to detect medicines, recreational drugs, their breakdown products, or poisons. The test’s purpose, detection limits, timing, and medicines taken affect interpretation.

\medterm{Drug-Induced Lupus} A lupus-like autoimmune reaction caused by certain medicines, often producing joint pain, muscle pain, fever, and inflammation around the lungs or heart. Symptoms and abnormal antibodies usually improve after the causative medicine is stopped under medical supervision.

\synonyms
Drug-Induced Lupus Erythematosus; DILE

\medterm{Drugs, Statin} A group of medicines that lowers LDL (“bad”) cholesterol by reducing cholesterol production in the liver. Statins lower the risk of heart attack and stroke; muscle symptoms and liver abnormalities can occur.

\medterm{Drugs That May Cause Erection Problems} Some medicines can contribute to difficulty getting or keeping an erection by affecting blood flow, hormones, nerves, or mood. A clinician should review the timing and all medicines before changing treatment.
\medterm{Drusen} Yellow-white deposits that form beneath the retina, often with aging. Small deposits may cause no symptoms; many or larger deposits can be associated with age-related macular degeneration and need eye examination.

\medterm{Dry Age-Related Macular Degeneration} Alternate index label; see \textit{Dry macular degeneration}.

\medterm{Dry AMD} Alternate index label; see \textit{Dry macular degeneration}.

\medterm{Dry Cell Battery Poisoning} Injury or poisoning caused by swallowing, leaking, or otherwise contacting a dry-cell battery. A swallowed button battery can rapidly burn the esophagus and is an emergency; do not induce vomiting and seek urgent medical help.
\medterm{Dry Cough} A cough that produces little or no mucus. Common causes include viral infection, asthma, acid reflux, allergies, smoking, certain medicines, and other airway or lung conditions; persistent cough needs assessment.

\medterm{Dry Cupping} A complementary treatment in which cups create suction on intact skin without making cuts. It may temporarily relieve some muscle pain, but evidence is limited; bruising, burns, skin injury, infection, and worsening of skin disease can occur.

\medterm{Dry Eye} A disorder in which tears are insufficient or evaporate too quickly, causing burning, grittiness, redness, watering, light sensitivity, or blurred vision. Long-lasting disease can damage the eye surface.

\synonyms
Dry-eye disease; keratoconjunctivitis sicca; Dry Eyes

\medterm{Dry Eye Disease} Alternate name; see \textit{Dry Eye}.

\medterm{Dry Eye Syndrome} Alternate name; see \textit{Dry Eye}.


\medterm{Dry Gangrene} Tissue death caused mainly by severely reduced arterial blood flow. The affected area becomes dry, shrunken, and dark, often with a clear boundary; infection can develop and urgent vascular assessment is needed.

\medterm{Dry Hair} Hair that feels rough, brittle, dull, or frizzy because its outer layer is damaged or lacks moisture. Heat, chemical treatments, weather, frequent washing, nutritional problems, and thyroid disease may contribute.

\medterm{Dry Ice} Solid carbon dioxide that changes directly into a cold gas. Contact can cause severe cold burns, and gas in an enclosed space can displace oxygen and cause unconsciousness or suffocation; never seal it in a container.

\medterm{Dry Macular Degeneration} The more common form of age-related macular degeneration, in which light-sensing cells in the macula gradually deteriorate without bleeding from new abnormal vessels. It can cause slow loss of detailed central vision.

\synonyms
Dry Age-Related Macular Degeneration
Dry AMD
Macular Degeneration, Dry

\medterm{Dry Mouth} A feeling or condition of having too little saliva. It can cause thirst, sticky mouth, difficulty speaking or swallowing, altered taste, tooth decay, mouth sores, and infections; medicines are a common cause.
\synonyms
xerostomia.

\medterm{Dry Mouth During Cancer Treatment} Reduced saliva during or after cancer treatment, especially head-and-neck radiation or certain medicines. It can cause mouth discomfort, difficulty chewing or swallowing, tooth decay, taste changes, and oral infection.

\medterm{Dry Powder Inhaler} A breath-powered device that delivers powdered medicine into the lungs. It does not require pressing and breathing at the same time, but the user must inhale strongly and correctly for the dose to reach the airways.


\medterm{Dry Skin} Skin that is rough, scaly, tight, or itchy because it has lost moisture or its protective barrier is damaged. Aging, cold or dry air, frequent washing, eczema, medicines, and some illnesses can contribute.


\medterm{Dry Socket} Painful inflammation after tooth extraction when the protective blood clot is lost or breaks down too early. Severe pain may begin a few days later and spread to the ear or jaw; dental treatment is needed.

\medterm{Alveolar Osteitis} Painful inflammation of a tooth socket after extraction, usually caused by loss or breakdown of the blood clot that normally protects the socket. It is also called dry socket.

\synonyms
alveolar osteitis.
Alveolar Osteitis

\medterm{Dry Socket (Alveolar Osteitis)} Alternate name; see \textit{Dry Socket}.

\medterm{Dry Socket Overview} Alternate name; see \textit{Dry Socket}.

\medterm{Dryness Vaginal (Vaginal Dryness And Vaginal Atrophy)} Alternate index label for Vaginal Dryness and Vaginal Atrophy.

\medterm{Dry-Powder Inhaler} A device that delivers powdered medicine to the lungs when a person inhales through it. Correct preparation, a strong breath, and proper technique are needed; the device should not usually be used with a spacer.

\medterm{DSM} The Diagnostic and Statistical Manual of Mental Disorders, a reference that describes recognized mental-health diagnoses and their criteria. Clinicians use it with history and examination; it does not replace individual clinical judgment.

\medterm{DSPS} Alternate index label; see \textit{Delayed sleep phase}.

\medterm{DSRCT} Alternate index label; see \textit{Desmoplastic small round cell tumor}.

\medterm{DSWPD} Alternate index label; see \textit{Delayed sleep phase}.

\medterm{DTaP} A vaccine for infants and young children that protects against diphtheria, tetanus, and pertussis (whooping cough). It is given as a series, with later booster vaccines used as children grow older.


\medterm{Dual Antiplatelet Therapy} Treatment with two medicines that reduce platelet clumping, usually after a heart attack or coronary stent. It lowers clot risk but increases bleeding risk; the duration depends on the procedure, heart condition, and bleeding risk.

\medterm{Dual Energy X-Ray Absorptiometry} See \textit{Bone Mineral Density Test}.

\synonyms
DXA; DEXA

\medterm{Dual-Energy CT} A CT scan that uses two X-ray energy levels to help distinguish materials in the body. It can improve evaluation of kidney stones, gout, blood vessels, bleeding, and some tissues while using the usual CT principles and risks.

\medterm{Dual-Energy X-Ray Absorptiometry} Alternate spelling; see \textit{Dual Energy X-Ray Absorptiometry}.

\medterm{Dual-Photon Absorptiometry} An older test that used two photon energy levels to estimate bone mineral content. It has largely been replaced by dual-energy X-ray absorptiometry for measuring bone density and fracture risk.
\medterm{Duane Retraction Syndrome} A birth-related eye-movement disorder caused by abnormal nerve control of an eye muscle. Turning the eye inward or outward may be limited, and the eyeball may pull backward and the eyelids narrow.

\medterm{Dubin-Johnson Syndrome} An inherited liver-transport disorder in which conjugated bilirubin is not moved normally into bile. It causes intermittent mild jaundice, often triggered by illness or stress, but usually does not cause serious liver damage.

\medterm{Dubowitz Syndrome} A rare developmental condition associated with poor growth, a small head, characteristic facial features, developmental or learning difficulties, and variable immune, skin, or behavioral problems. Features differ among affected people.

\medterm{Duchenne Muscular Dystrophy} An inherited muscle disease, usually affecting boys, caused by changes in the dystrophin gene. Progressive muscle weakness begins in childhood and may affect walking, breathing, the heart, and learning or behavior.

\medterm{Duck Embryo Vaccine} An older rabies vaccine made by growing inactivated rabies virus in duck embryos. It has largely been replaced by safer, more consistent cell-culture vaccines and is not commonly used today.

\medterm{Duct} A tube or channel that carries fluid within the body, such as bile, pancreatic juice, saliva, tears, or secretions from a gland. Blockage or injury can cause pain, swelling, infection, or impaired organ function.

\medterm{Ductal Carcinoma} A cancer that begins in cells lining a duct. The term most often describes breast cancer or pancreatic cancer; its behavior and treatment depend on the organ, spread, and tumor features.

\medterm{Ductal Carcinoma In Situ} Abnormal cancer-like cells confined to the milk ducts of the breast without invasion into nearby breast tissue. It has not spread elsewhere, but some cases can become invasive if untreated.

\synonyms
DCIS
Intraductal Carcinoma

\medterm{Ductal-Dependent Lesions} Serious congenital heart defects in which blood flow to the lungs or body depends on the fetal ductus arteriosus remaining open after birth. Affected newborns may rapidly become blue, breathless, pale, or shocked and need urgent specialist care.

\medterm{Ductless Glands} Glands that release hormones directly into the bloodstream rather than through a tube. Examples include the pituitary, thyroid, adrenal glands, and pancreatic islets; their hormones regulate growth, metabolism, stress, and reproduction.

\medterm{Ductus Arteriosus} A blood vessel in the fetus connecting the pulmonary artery to the aorta, allowing blood to bypass the lungs before birth. It normally closes soon after birth; failure to close can strain the heart and lungs.

\medterm{Ductus Venosus} A fetal blood vessel that directs oxygen-rich blood from the umbilical vein toward the heart while partly bypassing the liver. It normally closes after birth; abnormal fetal flow can indicate circulatory or placental problems.

\medterm{Duffy-Null Associated Neutrophil Count} A naturally lower measured neutrophil count associated with the Duffy-null blood-group pattern, common in some populations. It usually does not increase infection risk and should not automatically be treated as disease.

\medterm{Duke Criteria} A set of clinical, blood-culture, and heart-imaging findings used to classify suspected infective endocarditis. The criteria combine major and minor findings; the result supports diagnosis but must be interpreted with the patient’s history and examination.

\synonyms
Modified Duke Criteria; Duke Endocarditis Criteria

\medterm{Dukes Classification} An older system for describing how far colorectal cancer has spread through the bowel wall, into nearby lymph nodes, or to distant organs. Modern staging systems are more detailed and are generally preferred.

\medterm{Dull Pain} Aching, heavy, or pressure-like pain that is not sharp or stabbing. It can arise from muscles, organs, inflammation, or other tissues, and its significance depends on location, duration, associated symptoms, and examination findings.

\medterm{Duloxetine} A medicine that increases serotonin and norepinephrine signaling in the brain. It treats depression, anxiety, some nerve pain, fibromyalgia, and chronic muscle or bone pain; nausea, sleep changes, sweating, and withdrawal symptoms may occur.

\medterm{Duloxetine Hydrochloride} The salt form of duloxetine, a medicine used for depression, anxiety, some nerve pain, fibromyalgia, and chronic muscle or bone pain. It may cause nausea, sleep changes, sweating, blood-pressure changes, or withdrawal symptoms.

\medterm{Dumbbell Ossification} Bone formation or a bone lesion shaped like a dumbbell, often extending through an opening between the spinal canal and surrounding tissues. Imaging is needed to identify its cause and relationship to nerves or other structures.

\medterm{Dumdum Fever} An older name for visceral leishmaniasis, a parasitic infection transmitted by sandflies. It can cause prolonged fever, weight loss, enlarged spleen or liver, anemia, low blood counts, and life-threatening illness without treatment.

\medterm{Dumping Syndrome} Symptoms caused by food moving too quickly from the stomach into the small intestine, often after stomach surgery. Early symptoms include cramps, diarrhea, flushing, sweating, dizziness, and a fast heartbeat; later low blood sugar may occur.

\medterm{Dunbar Syndrome} Alternate index label; see \textit{Median arcuate ligament syndrome}.

\medterm{Duodenal} Relating to the duodenum, the first part of the small intestine immediately beyond the stomach. This segment receives bile and pancreatic juice and begins digestion and absorption of nutrients.
\medterm{Duodenal Adenocarcinoma} A cancer arising from gland-forming cells in the duodenum. It may cause abdominal pain, vomiting, bleeding, anemia, jaundice, or blockage; diagnosis usually requires imaging, endoscopy, and tissue sampling.

\medterm{Duodenal Atresia} A birth defect in which the duodenum is completely blocked, preventing food and digestive fluid from passing. Newborns usually develop green vomiting and feeding difficulty; it is often associated with other birth defects and requires surgery.

\medterm{Duodenal Bulb} The widened first part of the duodenum just beyond the stomach outlet. It receives partly digested food and is a common site for peptic ulcers; inflammation or scarring can narrow it and delay passage.

\medterm{Duodenal Disease} Any disorder affecting the duodenum, such as inflammation, ulcer, infection, bleeding, blockage, poor nutrient absorption, or a growth. Symptoms may include upper-abdominal pain, vomiting, anemia, or weight loss.

\medterm{Duodenal Diverticulum} A pouch that bulges outward from the duodenal wall, often near the openings of the bile and pancreatic ducts. It commonly causes no symptoms but can rarely lead to bleeding, inflammation, blockage, pancreatitis, or bile-duct infection.

\medterm{Duodenal Fistula} An abnormal connection between the duodenum and another organ, bowel segment, or skin. Digestive fluid may leak through it, causing infection, skin damage, dehydration, electrolyte loss, or poor nutrition.

\medterm{Duodenal Hemorrhage} Bleeding from the duodenum, often caused by a peptic ulcer, inflammation, abnormal blood vessel, injury, or cancer. It may cause black stools, vomiting blood, weakness, dizziness, anemia, or shock.

\medterm{Duodenal Infection} Infection affecting the duodenum, caused by bacteria, parasites, viruses, or spread from nearby digestive disease. Symptoms can include upper-abdominal pain, nausea, vomiting, diarrhea, fever, or bleeding.

\medterm{Duodenal Neoplasm} An abnormal growth in the duodenum that may be benign or cancerous. It can cause pain, bleeding, vomiting, blockage, anemia, or jaundice; evaluation depends on its size, location, appearance, and tissue type.

\medterm{Duodenal Obstruction} A partial or complete blockage that prevents food and digestive fluid from passing through the duodenum. It can cause repeated vomiting, abdominal swelling, pain, and dehydration and may need urgent imaging and treatment.

\medterm{Duodenal Papilla} A small raised opening in the duodenal wall where bile and pancreatic juice enter the intestine. Narrowing, blockage, or inflammation at this site can cause jaundice, pancreatitis, or digestive symptoms.

\medterm{Duodenal Perforation} A hole through the full thickness of the duodenal wall, usually caused by a severe ulcer, injury, or a medical procedure. Sudden severe abdominal pain, rigid muscles, and infection of the abdominal lining require emergency treatment.

\medterm{Duodenal Stenosis} Narrowing of the duodenum that partly blocks passage of food and digestive fluid. It may cause vomiting, abdominal swelling, pain, early fullness, poor feeding, or dehydration; treatment depends on the cause and severity.

\medterm{Duodenal Ulcer} An open sore in the duodenal lining, commonly related to Helicobacter pylori infection or anti-inflammatory medicines. It may cause upper-abdominal pain and can bleed, perforate the bowel, or obstruct food passage.

\medterm{Duodenitis} Inflammation of the duodenal lining, which may result from Helicobacter pylori, excess acid, anti-inflammatory medicines, infection, celiac disease, or inflammatory bowel disease. It can cause upper-abdominal pain, nausea, vomiting, or bleeding.

\medterm{Duodenoscope} A flexible viewing tube passed through the mouth and stomach to examine the duodenum and reach the openings of the bile and pancreatic ducts. It may be used during procedures to remove stones or place stents.

\medterm{Duodenostomy} A surgically created opening from the duodenum to the skin or another part of the intestine. It may be used in selected cases for feeding, drainage, or bypassing a blockage.

\medterm{Duodenum} The first part of the small intestine, between the stomach and jejunum. It receives bile and pancreatic juice, mixes them with stomach contents, and begins much of the digestion and absorption of nutrients.

\medterm{Duplex Doppler Ultrasound Scanning} An ultrasound test that shows the structure of blood vessels and measures the direction and speed of blood flow. It can detect narrowing, blockage, clots, aneurysms, or abnormal connections without radiation.
\synonyms
Duplex ultrasound

\medterm{Duplex Surveillance} Planned examinations and ultrasound checks of a blood vessel, graft, stent, aneurysm, or dialysis access. Monitoring can detect recurrent narrowing, clotting, leakage, or failure before serious symptoms develop.

\medterm{Duplex Ultrasound} An ultrasound examination that combines pictures of body structures with Doppler measurement of blood flow. It is commonly used to assess blood vessels and can identify narrowing, blockage, clots, or abnormal flow.

\medterm{Duplication} A genetic change in which a section of DNA or a chromosome is copied one or more extra times. The additional genetic material can alter how much of a gene is active and may cause a health condition, although some duplications have little or no effect.

\medterm{Dupuytren Contracture} Progressive thickening and shortening of tissue in the palm that pulls one or more fingers toward the palm and can limit hand function.

\synonyms
Dupuytrens Contracture (Dupuytren Contracture)

\medterm{Dupuytren, Guillaume} A French surgeon whose name is associated with Dupuytren contracture, progressive thickening and shortening of palmar fascia that bends the fingers toward the palm.

\medterm{Dupuytren's Contracture} A progressive fibrosing disorder of the palmar fascia that produces nodules and cords and may permanently flex one or more fingers toward the palm. Treatment depends on severity and functional impairment and may include observation, injection, needle fasciotomy, or surgery.

\medterm{Dupuytrens Contracture (Dupuytren Contracture)} Alternate index label for Dupuytren Contracture.

\medterm{Dupuytren's Disease} A usually progressive disorder in which the palmar fascia, a layer of fibrous tissue beneath the skin of the palm and fingers, becomes thick and forms nodules and cords. The cords can gradually prevent one or more fingers, especially the ring or little finger, from straightening and may limit hand function. It usually begins in middle age or later and is more common in men, although anyone can develop it. The condition begins in the palm and is not primarily a disease of the finger joints; the later bending deformity is called Dupuytren contracture.

\medterm{Dura} The tough outermost layer of protective tissue surrounding the brain and spinal cord. It helps contain cerebrospinal fluid, supports blood vessels, and can become inflamed, torn, thickened, or infected.

\medterm{Dura Mater} The tough outer layer of tissue surrounding and protecting the brain and spinal cord. It helps support blood vessels and contains the fluid-filled coverings around the central nervous system.

\medterm{Durable Power Of Attorney} A legal authorization allowing a designated person to act or make decisions for another person when specified conditions are
met. In healthcare, it may appoint a healthcare proxy.

\medterm{Dural Arteriovenous Fistulas} Abnormal links between arteries and veins in the protective covering around the brain or spinal cord. They can cause headaches, pulsating sounds, seizures, bleeding, or neurologic problems and may need vascular treatment.

\medterm{Dural Sac} The tough, tubular covering around the spinal cord, its surrounding fluid, and the spinal nerve roots inside the spinal canal. A tear can allow spinal-fluid leakage and cause positional headaches.

\medterm{Dural Venous Sinuses} Channels between layers of the brain’s tough covering that collect blood and cerebrospinal fluid and drain them toward the neck veins. A clot in these channels can cause headache, seizures, vision problems, or stroke-like symptoms.

\medterm{Duramorph} A brand name for preservative-free morphine used by injection, including in epidural or intrathecal analgesia; respiratory depression is a serious risk.

\medterm{Durapatite} Another name for hydroxyapatite, the calcium-phosphate mineral that gives bones and teeth much of their hardness. Its amount and arrangement contribute to bone strength and can change in metabolic bone disease.
\medterm{Durules} A formulation designed to release a medicine slowly over an extended period. Such tablets or pellets must be used exactly as directed because crushing or chewing may release too much medicine at once.

\medterm{Dust} Small solid particles suspended in air or settled on surfaces. Inhaled dust can irritate the nose and lungs; some types can cause allergy, asthma, infection, poisoning, or long-term lung disease.

\medterm{Dust Disease} Lung or other illness caused by repeated exposure to harmful dust, often at work. Effects depend on the dust and exposure and may include asthma, scarring of lung tissue, chronic cough, or cancer.

\medterm{Dust Mite Allergy} An immune reaction to proteins from tiny mites living in household dust. It can cause sneezing, blocked or itchy nose, watery eyes, cough, wheezing, eczema flares, or asthma symptoms.

\synonyms
Allergy, Dust Mite

\medterm{Dutasteride} A medicine that lowers the hormone dihydrotestosterone and gradually shrinks an enlarged prostate. It improves urinary symptoms and lowers retention risk but may reduce sexual function and lowers the blood-test level used to screen for prostate cancer.

\medterm{DVT} A blood clot in a deep vein, usually in the leg. It can cause swelling, pain, warmth, or redness and may travel to the lungs, causing a life-threatening pulmonary embolism.

\medterm{DVT Prevention} Steps to reduce the risk of a deep-vein blood clot, including early movement, leg exercises, compression devices, and preventive anticoagulant medicine when appropriate. The choice depends on clot risk and bleeding risk.

\medterm{Dwarfism} A medical condition causing unusually short adult height, often defined as about 4 feet 10 inches (147 cm) or less. Causes include skeletal disorders, growth-hormone problems, genetic conditions, and chronic disease; features vary by cause.

\medterm{Dwarfism (Achondroplasia)} Alternate index label for Achondroplasia.

\medterm{DWI FLAIR Mismatch} A pattern on brain MRI in which an area is abnormal on diffusion images but not yet clearly abnormal on FLAIR images. In selected people with an unknown stroke onset time, it may help identify potentially treatable recent stroke.
\medterm{Dx} A medical abbreviation for diagnosis or diagnostic assessment. Because abbreviations can be misunderstood, the full word should be used in patient instructions and other documents where clarity matters.
\medterm{DXA Scan (Dual-Energy X-Ray Absorptiometry)} Alternate name; see \textit{Dual Energy X-Ray Absorptiometry}.

\medterm{D-Xylose Absorption} An older test of how well the small intestine absorbs a simple sugar after it is swallowed. Low amounts in blood or urine may suggest intestinal lining disease, but kidney function and collection quality affect results.

\medterm{Dyad} A pair of related people, structures, cells, or measurements considered together. Medical examples include two linked findings or two parts of a cellular or body structure that function together.

\medterm{Dye} A colored substance used to stain tissue, identify cells, make structures visible, or provide contrast during a medical test. Some dyes can cause allergic reactions or affect the kidneys, especially when injected.

\medterm{Dye Lasers} Lasers that use a liquid organic dye to produce adjustable wavelengths of light. In medicine, selected wavelengths can treat blood-vessel or pigmented skin lesions; eye and skin protection is essential.

\medterm{Dye Remover Poisoning} Harm from swallowing, inhaling, or getting dye-removing chemicals on the skin or eyes. Depending on the ingredients, exposure can cause burns, breathing problems, vomiting, poisoning, or organ injury and needs poison-control advice.

\medterm{Dying} The physical process of approaching death, often involving increasing weakness, sleepiness, reduced appetite and drinking, changes in breathing, and less responsiveness. Emotional, social, cultural, and spiritual support may help the person and family.

\medterm{Dynamin Gtpase} A family of GTP-binding proteins that constrict membranes during vesicle budding, endocytosis, and other cell-membrane processes.

\medterm{Dynamometer} An instrument that measures force, commonly hand-grip or muscle strength. Clinicians use the result to compare sides, identify weakness, monitor recovery, and assess physical function during rehabilitation or neurologic examination.

\medterm{Dynein Atpase} The energy-using activity of dynein, a motor protein that moves materials inside cells along microtubules. Dynein also powers the movement of cilia and flagella, which help move fluid or cells.

\medterm{Dynorphin} A naturally occurring opioid peptide that mainly activates kappa-opioid receptors in the nervous system. It influences pain, stress, mood, movement, and hormone regulation.

\medterm{Dyphylline} A bronchodilator used in some countries to relax airway muscles in asthma or chronic obstructive lung disease. It may cause nausea, tremor, headache, nervousness, or a fast heartbeat.

\medterm{Dysaesthesia} An unpleasant abnormal sensation, such as burning, electric-shock pain, pins and needles, or excessive sensitivity to light touch. It commonly results from nerve injury, diabetes, medicines, or other nerve disease.

\medterm{Dysarthria} Slurred, weak, slow, or poorly coordinated speech caused by a problem in the brain, nerves, or muscles used for speaking. Sudden dysarthria can be a sign of stroke.

\medterm{Dysautonomia} A disorder of the automatic nervous system that controls blood pressure, heart rate, sweating, temperature, digestion, bladder function, and sexual function. Causes include nerve disease, diabetes, medicines, and inherited conditions.

\medterm{Dysbarism} Illness caused by a change in surrounding pressure, usually during diving, flying, or hyperbaric treatment. It includes decompression sickness and pressure injury to the ears, sinuses, lungs, or other tissues.

\medterm{Dysbiosis} An imbalance or change in the community of microorganisms living in the body, especially the intestine. Antibiotics, infection, diet, and illness may contribute; dysbiosis is associated with some digestive and inflammatory disorders, but it is not always a proven cause.

\synonyms
Gut Microbiome Imbalance

\medterm{Dyscalculia} A persistent difficulty understanding numbers, quantities, or arithmetic that is not explained by inadequate teaching or general intellectual disability. It may begin during development or follow brain injury and can affect school or daily tasks.

\medterm{Dyschezia} Difficult or painful passage of stool, often with excessive straining or a feeling of incomplete emptying. Causes include constipation, pelvic-floor coordination problems, anal disorders, and other bowel disease.

\medterm{Dyscrasia} An older broad term for an abnormal blood or bone-marrow condition. It may describe abnormal blood cells, clotting, or proteins in the blood, so the specific disorder should be named when possible.

\synonyms
Blood dyscrasia

\medterm{Dyscrasias} An older plural term for disorders affecting blood, bone marrow, or body-fluid composition. It may refer to abnormal blood cells, clotting, or abnormal proteins, so the specific disorder must be identified.

\medterm{Dysembryoplastic Neuroepithelial Tumor} A usually slow-growing, low-grade brain tumor that develops most often in the outer part of the brain in children or young adults. It commonly causes long-standing focal seizures and is often treatable with surgery.

\medterm{Dysentery} Severe intestinal infection causing frequent diarrhea with blood or mucus, abdominal cramps, fever, and an urgent feeling of needing to pass stool. Bacteria or parasites may cause it; dehydration and complications require medical care.

\medterm{Dysesthesia} An unpleasant abnormal sensation, such as burning, electric-shock pain, or painful sensitivity to light touch. It may result from injury or disease of nerves in the brain, spinal cord, or limbs.

\medterm{Dysferlinopathy} A group of inherited muscle diseases caused by changes in both copies of the DYSF gene, which makes a muscle-membrane repair protein. Weakness often begins in adolescence or early adulthood, affecting the hips, shoulders, or calves.

\synonyms
DYSF-Related Muscular Dystrophy; Dysferlin Deficiency

\medterm{Dysfibrinogenemia} An inherited or acquired condition in which fibrinogen, a blood-clotting protein, has abnormal structure or function. It may cause unusual bleeding, excessive clotting, or no symptoms and is diagnosed with clotting and specialized blood tests.

\medterm{Dysfunction} Abnormal, reduced, or poorly controlled function of an organ, tissue, body system, or process. It may be temporary or permanent and can result from disease, injury, medicines, inherited conditions, stress, or aging.

\medterm{Dysfunctional Voiding} Abnormal coordination between the bladder muscle and the muscles around the urethra during urination. It may cause a weak or interrupted stream, straining, incomplete emptying, urinary infections, wetting, or constipation, especially in children.

\medterm{Dysfunctional Uterine Bleeding} An older term for bleeding from the uterus that is unusually heavy, prolonged, frequent, or irregular without an identified structural cause. Current evaluation considers pregnancy, medicines, hormones, ovulation, and uterine disease.

\medterm{Dysgammaglobulinemia} An abnormal level or pattern of antibodies in the blood, with some antibody types low or absent and others normal or increased. It may cause repeated infections, autoimmune problems, or no symptoms and is identified by blood tests.

\medterm{Dysgerminoma} A cancer of the ovary arising from germ cells, the cells that can develop into eggs. It occurs most often in adolescents and young adults and may cause pelvic pain, swelling, or an abdominal mass.

\medterm{Dysgeusia} An altered sense of taste, including reduced, distorted, or unpleasant taste. Causes include viral infection, mouth or nose disease, nerve injury, nutritional deficiency, dry mouth, medicines, and other illnesses.

\medterm{Dysgraphia} A persistent difficulty with handwriting, spelling, or expressing ideas in writing despite appropriate learning opportunities. It may reflect problems with motor skills, language, attention, or brain function and can occur with other learning disorders.

\medterm{Dyshidrosis} A skin condition causing small, intensely itchy blisters on the palms, sides of the fingers, or soles. The skin may peel, crack, or become sore; triggers include sweating, irritation, allergy, and stress.

\synonyms
Eczema, Dyshidrotic

\medterm{Dyskeratosis Congenita} An inherited disorder affecting chromosome-end maintenance. It may cause abnormal skin coloring, nail changes, white patches in the mouth, bone-marrow failure, lung or liver disease, developmental problems, and increased cancer risk.

\medterm{Dyskinesia} Abnormal movement that is often involuntary, excessive, repetitive, twisting, or irregular. It may result from neurologic disease or medicines, including long-term treatment for movement disorders, and can interfere with walking, speaking, or daily activities.

\medterm{Dyslexia} A learning disorder involving persistent difficulty recognizing words accurately or fluently, decoding them, spelling, or linking letters with speech sounds. It is not caused by low intelligence or poor motivation and may continue into adulthood.

\synonyms
Reading Disability, Specific
Specific Reading Disability

\medterm{Dyslipidemia} An unhealthy level or pattern of fats in the blood, such as high LDL cholesterol, high triglycerides, low HDL cholesterol, or a combination. It can promote artery disease and, when triglycerides are very high, pancreatitis.

\medterm{Dyslipidemias} Conditions involving abnormal blood fats, including high LDL cholesterol, high triglycerides, low HDL cholesterol, or mixed changes. They can increase the risk of artery disease and pancreatitis when triglycerides are very high.

\medterm{Dysmenorrhea} Painful menstrual periods, usually causing cramping in the lower abdomen and sometimes back pain, nausea, diarrhea, headache, or dizziness. Severe, new, or worsening pain may indicate an underlying pelvic disorder.

\medterm{Dysmenorrhoea} Painful menstrual periods. Primary pain occurs without another pelvic disease, while secondary pain results from a condition such as endometriosis, fibroids, adenomyosis, pelvic infection, or an intrauterine device.

\medterm{Dysmetria} Inability to judge or control the distance, force, or speed of a voluntary movement. A person may overshoot or undershoot a target; it commonly reflects dysfunction of the cerebellum or its connecting pathways.

\medterm{Dysorexia} An abnormal appetite, including little desire to eat, excessive hunger, or disordered eating. It may result from illness, medicines, emotional distress, metabolic disease, or an eating disorder and can affect nutrition and weight.

\medterm{Dysostoses} Inherited or developmental conditions in which one or more bones form or grow abnormally. They may affect the skull, face, spine, ribs, arms, or legs and can alter body shape, movement, growth, hearing, breathing, or other functions.

\medterm{Dyspareunia} Recurrent or persistent genital or pelvic pain before, during, or after sexual activity, especially attempted vaginal penetration. Pain may be at the vaginal opening or deeper in the pelvis and can result from dryness, infection, skin disease, pelvic-floor muscle spasm, endometriosis, inflammation, injury, or emotional factors.

\medterm{Dyspepsia} Recurrent discomfort in the upper abdomen, such as burning, early fullness, bloating, nausea, or belching. Causes include ulcers, reflux, medicines, infection, and other illness; sometimes no structural cause is found.

\medterm{Dyspeptic} Describing indigestion symptoms such as upper-abdominal discomfort or burning, early fullness, bloating, belching, or nausea, especially after eating.
\medterm{Dysphagia} See \textit{Difficulty Swallowing}.

\synonyms
Difficulty Swallowing
Trouble Swallowing

\medterm{Dysphagia Diet} Food and drinks altered in texture or thickness to make swallowing safer after a swallowing assessment. The plan should reduce aspiration risk while providing enough calories, protein, fluids, and enjoyable foods.

\medterm{Dysphagia (Geriatric)} Difficulty swallowing in an older adult. It may cause coughing or choking during meals, food remaining in the mouth, weight loss, dehydration, or aspiration and can result from stroke, dementia, muscle disease, blockage, or medicines.

\synonyms
Swallowing difficulty in older adults

\medterm{Dysphagia Rehabilitation} Therapy to improve safe swallowing, nutrition, and hydration. It may include swallowing exercises, posture or pacing strategies, food-texture changes, and caregiver education based on the cause and assessment.

\medterm{Dysphasia} Difficulty using or understanding language because of a problem in the brain's language areas. It may affect speaking, understanding, reading, or writing; many clinicians use the clearer term aphasia.

\medterm{Dysphonia} An abnormal voice, such as hoarseness, breathiness, roughness, strain, weakness, or loss of voice. Causes include infection, vocal-cord swelling or growths, nerve problems, voice overuse, reflux, and cancer; persistent symptoms need evaluation.

\medterm{Dysphoria} A state of marked emotional distress, unease, or dissatisfaction. Gender dysphoria specifically means distress related to a difference between experienced gender and sex assigned at birth; it is not itself a mental illness.

\medterm{Dysplasia} Abnormal size, shape, or arrangement of cells in a tissue. It is not cancer, but some forms—especially when severe or persistent—can increase the chance of cancer and may need monitoring or treatment.

\medterm{Dysplasia Cervical (Cervical Dysplasia)} Alternate index label for Cervical Dysplasia.

\medterm{Dysplastic Nevi} Atypical moles with unusual size, shape, color, border, or microscopic features. Most do not become melanoma, but having many can increase melanoma risk; a changing, bleeding, or suspicious mole should be examined.

\medterm{Dysplastic Nevus} An atypical mole that may be larger, unevenly colored, or have an irregular border. Most do not become cancer, but having many atypical moles can increase melanoma risk; changing lesions should be examined.

\medterm{Dyspnea} The feeling of difficult, uncomfortable, or insufficient breathing. Causes include asthma, lung or heart disease, anemia, infection, muscle weakness, anxiety, and other conditions.

\medterm{Dyspnea On Exertion} Shortness of breath that occurs or worsens with physical activity. It may result from heart or lung disease, anemia, poor fitness, or other conditions; new, severe, or rapidly worsening symptoms need prompt assessment.

\medterm{Dyspnoea} The feeling of difficult, uncomfortable, or insufficient breathing. It may result from heart or lung disease, anemia, infection, poor fitness, anxiety, or other causes; sudden or severe symptoms require emergency care.

\medterm{Dyspraxia} Difficulty planning and coordinating purposeful movements. It may affect dressing, writing, balance, speaking, or sports; in children it may be part of developmental coordination disorder.

\medterm{Dyspraxia (Developmental Coordination Disorder)} A developmental difficulty planning and coordinating movements that interferes with daily activities, learning, or motor skills. It is not explained by muscle weakness, loss of sensation, or inadequate opportunity to learn.

\synonyms
Developmental Coordination Disorder

\medterm{Dysraphism} A birth defect caused by incomplete joining of developing body structures, especially the spinal cord or spinal bones. Effects range from a small skin mark to weakness, loss of sensation, or bladder and bowel problems.

\medterm{Dysreflexia (Autonomic Dysreflexia)} A dangerous sudden rise in automatic nervous-system activity, usually after spinal-cord injury above the upper chest. A full bladder, bowel problem, skin irritation, or other stimulus may cause severe high blood pressure, headache, sweating, and slow heartbeat.

\synonyms
Autonomic Dysreflexia

\medterm{Dysrhythmia} An abnormal heart rhythm that may be too fast, too slow, or irregular. Some cause no symptoms; others cause palpitations, dizziness, fainting, chest pain, stroke, or inadequate blood flow.


\medterm{Dyssomnia} A broad or older term for a sleep disorder involving difficulty sleeping, excessive sleepiness, or abnormal sleep timing. Modern diagnosis usually specifies the problem, such as insomnia, hypersomnia, or a circadian rhythm disorder.
\medterm{Dyssomnias} Sleep disorders that mainly cause difficulty falling asleep, staying asleep, waking at the desired time, or staying awake. Causes include sleep schedules, medical conditions, medicines, mood disorders, and abnormal sleep regulation.

\medterm{Dysthymia} Historical term; see \textit{Persistent Depressive Disorder}.

\medterm{Dysthymia (Persistent Depressive Disorder)} Alternate index label for Persistent Depressive Disorder.

\medterm{Dystocia} Difficult or unusually slow labor caused by ineffective contractions, the baby’s size or position, or the mother’s pelvis or soft tissues. It can increase risks for both parent and baby and may require assisted delivery or cesarean birth.

\medterm{Dystonia} A neurological movement disorder in which changes in the brain’s movement signals make muscles tighten or contract without the person intending them to. This can cause slow, repeated or twisting movements, tremor, or abnormal and sometimes painful postures. It may affect one body part, several connected areas, or most of the body, and symptoms may be triggered by movement, stress, or tiredness. Dystonia may be inherited, have no known cause, or result from brain disease or injury, infection, poisoning, or certain medicines. It can begin at any age and may remain stable or worsen over time.

\medterm{Dystonia, Cervical} Alternate index label; see \textit{Cervical dystonia}.

\medterm{Dystonia Disorder} A disorder causing repeated or sustained involuntary muscle contractions, which produce twisting movements, tremor, or abnormal postures. Symptoms may be limited to one area or involve the whole body.

\medterm{Dystonin} A protein made from the DST gene that helps attach the internal framework of cells to other cell structures. Harmful gene changes can cause inherited nerve disease, muscle weakness, and fragile or blistering skin.

\medterm{Dystrophia Unguium} Abnormal nail growth, shape, color, or texture. Causes include injury, fungal infection, inflammation, poor nutrition, medicines, reduced blood flow, and inherited disorders; treatment depends on the cause.

\medterm{Dystrophic Calcification} Buildup of calcium salts in damaged, dead, or chronically inflamed tissue even though the blood calcium level is normal. It can occur in scars, diseased valves, tumors, or injured muscles and may impair function.

\medterm{Dystrophin} A protein that strengthens the membrane of muscle cells during contraction. Too little or abnormal dystrophin damages muscle fibers and causes Duchenne or Becker muscular dystrophy.

\medterm{Dystrophin-Glycoprotein Complex} A group of proteins that connects the internal framework of muscle cells to their surrounding support tissue and protects the cell membrane during contraction. Changes in its components can cause several inherited muscular dystrophies.

\synonyms
DGC; Dystrophin-Associated Glycoprotein Complex (DAGC)

\medterm{Dystrophy} Abnormal development, maintenance, or degeneration of a tissue or organ, often caused by an inherited or metabolic condition. Effects depend on the tissue and may include muscle weakness, poor vision, fragile structures, or loss of function.

\medterm{Dysuria} Pain, burning, or discomfort during urination. Common causes include bladder or urethral infection, vaginal or penile inflammation, stones, irritation, and sexually transmitted infections; fever, flank pain, pregnancy, or blood in urine needs medical assessment.
\synonyms
Painful Urination
